\documentclass[a4paper,12pt]{report}
\usepackage[utf8]{inputenc}
\usepackage[ngerman]{babel}
\usepackage{amsmath}
\usepackage{amssymb}
\usepackage{amsthm}
\usepackage{geometry}
\usepackage{graphicx}
\usepackage[hidelinks]{hyperref}
\usepackage{listings}
\usepackage{xcolor}
\usepackage{enumitem}
\usepackage{mathtools}
\usepackage{booktabs}
\usepackage{float}
\usepackage{microtype}
\usepackage{ragged2e}

\tolerance=2000
\emergencystretch=3em
\hfuzz=2pt
\vfuzz=2pt
\usepackage{algorithm}
\usepackage{algpseudocode}
\usepackage{array}

\geometry{margin=2.5cm}

\theoremstyle{definition}
\newtheorem{definition}{Definition}[chapter]
\newtheorem{beispiel}{Beispiel}[chapter]
\newtheorem{satz}{Satz}[chapter]
\newtheorem{lemma}{Lemma}[chapter]
\newtheorem{bemerkung}{Bemerkung}[chapter]
\newtheorem{korollar}{Korollar}[chapter]
\newtheorem{proposition}{Proposition}[chapter]
\newtheorem{theorem}{Theorem}[chapter]

\setlist[itemize,1]{label=--}

% ---------- Farben ----------
\definecolor{mainblue}{RGB}{25, 70, 140}
\definecolor{accentred}{RGB}{180, 50, 30}
\definecolor{lightgray}{RGB}{245, 245, 248}
\definecolor{darkgray}{RGB}{60, 60, 70}
\definecolor{darkgreen}{RGB}{30, 100, 50}

\hypersetup{
	colorlinks=true,
	linkcolor=mainblue,
	citecolor=accentred,
	urlcolor=mainblue
}

% Code-Highlighting
\lstset{
	basicstyle=\ttfamily\small,
	keywordstyle=\color{blue},
	commentstyle=\color{gray},
	stringstyle=\color{red},
	numbers=left,
	numberstyle=\tiny,
	breaklines=true,
	frame=single,
	literate=%
	{"O}{{\"O}}1
	{Ä}{{\"A}}1
	{Ü}{{\"U}}1
	{ß}{{\ss}}1
	{ü}{{\"u}}1
	{ä}{{\"a}}1
	{ö}{{\"o}}1
}

\pagestyle{plain}

\title{%
	\textbf{Harte Echtzeit-Rechnersysteme}\\[0.5em]
	\large Vorhersagbare Scheduling-Algorithmen und Anwendungen
}
\author{Stephan Epp\\
	\small Universität Bielefeld\\
	\small Basierend auf Giorgio C. Buttazzo: \textit{Hard Real-Time Computing Systems}}
\date{2026}

\begin{document}

\maketitle

\tableofcontents
\newpage

%=============================================================
\chapter*{Vorwort}
\addcontentsline{toc}{chapter}{Vorwort}
%=============================================================

\section*{Motivation und Zielsetzung}

Dieses Buch behandelt harte Echtzeit-Rechnersysteme aus der Perspektive
vorhersagbarer Scheduling-Algorithmen. Es verfolgt ein doppeltes Ziel:
einerseits die präzise Darstellung der klassischen Theorie, die seit dem
Grundlagenaufsatz von Liu und Layland \cite{Liu1973} im Jahr 1973 entstanden
ist, andererseits die Integration neuerer Ergebnisse -- insbesondere des
hier erstmals ausführlich präsentierten EDF$^+$-Algorithmus \cite{Epp2025} --
in einen einheitlichen theoretischen Rahmen.

Die Hauptgrundlage ist das international maßgebliche Lehrbuch von Giorgio C.\
Buttazzo \cite{Buttazzo2011}, das seit seiner ersten Auflage 1997 als
Standardreferenz der Echtzeit-Scheduling-Theorie gilt. Es wird ergänzt durch
eine Vielzahl von Primärquellen, darunter die Originalarbeiten von Horn
\cite{Horn1974}, Dertouzos \cite{Dertouzos1974}, Jackson \cite{Jackson1955},
Lehoczky, Sha und Ding \cite{Lehoczky1989} sowie die einflussreichen Beiträge
zur Ressourcenverwaltung von Sha, Rajkumar und Lehoczky \cite{Sha1990} und
Baker \cite{Baker1991}.

\section*{Der fundamentale Kernsatz}

Der rote Faden dieses Buches ist eine Aussage, die sich aus dem Vergleich
zweier grundlegend verschiedener Scheduling-Philosophien ergibt -- der
statischen Prioritätszuweisung (RM) und der dynamischen Prioritätszuweisung
(EDF):

\begin{quote}
\textbf{Wenn es einen machbaren Schedule gibt, findet RM diesen sofort.
RM ist schwach optimal.}

\textbf{Wenn es einen machbaren Schedule gibt, findet EDF diesen sofort.
Die kürzeste Ausführungszeit aller Tasks findet EDF sofort.
EDF ist scharf optimal.}
\end{quote}

Diese Aussage ist nicht trivial. Sie unterscheidet zwei fundamentale Konzepte
der Optimalität: \emph{schwache} Optimalität, die RM innerhalb der Klasse
fester Prioritätsalgorithmen auszeichnet, und \emph{scharfe} Optimalität,
die EDF über alle präemptiven Einprozessor-Scheduling-Algorithmen hinweg
kennzeichnet. Die Schärfe von EDF -- im Sinne von: ``genau die schedulearbaren
Task-Sets werden gefunden, alle anderen nicht'' -- ist das stärkste
Optimalitätsergebnis, das für Einprozessor-Echtzeit-Scheduling bekannt ist.

Der Beweis der scharfen Optimalität von EDF, der auf Arbeiten von Horn
\cite{Horn1974} und Dertouzos \cite{Dertouzos1974} zurückgeht und später
durch Baruah et al.\ \cite{Baruah1990} für periodische Tasks präzisiert wurde,
ist mathematisch elegant und bildet das Herzstück von Kapitel~3 dieses Buches.

\section*{Historischer Kontext}

Die Echtzeit-Scheduling-Theorie entstand in den frühen 1970er Jahren als
Antwort auf die wachsende Verbreitung von Prozesssteuerungssystemen, die
zuverlässige zeitliche Garantien erforderten. Die Pionierarbeiten von Liu und
Layland \cite{Liu1973} legten 1973 mit dem Rate-Monotonic-Algorithmus das
Fundament. Fast gleichzeitig bewies Dertouzos \cite{Dertouzos1974} die
Optimalität von EDF für aperiodische Tasks.

In den 1980er und 1990er Jahren wurde dieses Fundament durch eine Reihe
wesentlicher Beiträge erweitert:
\begin{itemize}
	\item Lehoczky, Sha und Ding \cite{Lehoczky1989} lieferten eine exakte
	      Charakterisierung der RM-Schedulierbarkeit.
	\item Sha, Rajkumar und Lehoczky \cite{Sha1990} lösten das Prioritätsinversionsproblem
	      durch das Priority Inheritance Protocol.
	\item Sprunt, Sha und Lehoczky \cite{Sprunt1989} entwickelten den Sporadic Server,
	      der die Lücke zwischen periodischem und aperiodischem Scheduling überbrückt.
	\item Audsley et al.\ \cite{Audsley1993} etablierten die Response Time Analysis
	      als praktisches Werkzeug für feste Prioritätsschedules.
	\item Spuri und Buttazzo \cite{Spuri1996} übertrugen Server-Konzepte auf
	      EDF und entwickelten Total Bandwidth Server sowie Dynamic Sporadic Server.
	\item Abeni und Buttazzo \cite{Abeni1998} entwickelten den Constant Bandwidth
	      Server als Grundlage für ressourcenisoliertes Scheduling.
\end{itemize}

Im neuen Jahrtausend rückte das Mixed-Criticality-Problem (Vestal \cite{Vestal2007},
Baruah et al.\ \cite{BaruahMCS2012}) und das Mehrprozessor-Scheduling (Baruah
et al.\ \cite{BaruahPfair1996}) in den Vordergrund. Beide sind bis heute
aktive Forschungsfelder.

\section*{Über den EDF$^+$-Algorithmus}

Eine wichtige Erweiterung der klassischen EDF-Theorie ist der EDF$^+$-Algorithmus
\cite{Epp2025}, der in Kapitel~10 ausführlich behandelt wird. EDF$^+$ überträgt
die scharfe Optimalität von EDF auf Systeme mit stochastischen Blocking-Ereignissen:
Wenn Tasks mit Wahrscheinlichkeit $p$ blockieren (etwa durch I/O-Warten oder
Semaphore), erzeugt EDF$^+$ bei jedem Blocking-Ereignis sofort einen neuen
optimalen EDF-Schedule für die verbleibenden nicht-blockierten Tasks.

Der Beweis der Optimalität von EDF$^+$ beruht direkt auf dem EDF-Optimalitätstheorem
und ist konstruktiv: Durch vollständige Induktion über die Anzahl der
Blocking-Ereignisse wird gezeigt, dass EDF$^+$ stets einen Deadline-freien
Schedule liefert, sofern die Systemauslastung $U \leq 1$ beträgt. Dies macht
EDF$^+$ zu einem Algorithmus mit nachweisbarer formaler Korrektheit unter
realistischen Systemannahmen.

\section*{Struktur des Buches}

Das Buch ist in 13 Kapitel gegliedert, die aufeinander aufbauen:

\textbf{Kapitel~1 -- Allgemeine Betrachtung:} Grundlegende Eigenschaften von
Echtzeit-Systemen, Typen von Echtzeit-Tasks, Quellen des Nicht-Determinismus
und die zentrale Anforderung der Vorhersagbarkeit nach Stankovic \cite{Stankovic1988}.

\textbf{Kapitel~2 -- Grundlegende Konzepte:} Vollständiges Task-Modell mit
allen Zeitparametern (Lateness, Laxität, Slack), Klassifikation von Scheduling-Problemen
nach der Graham-Notation $\alpha|\beta|\gamma$, Gütekriterien, Scheduling-Anomalien
und der Dhall-Effekt auf Mehrprozessoren.

\textbf{Kapitel~3 -- Aperiodisches Task-Scheduling:} Jacksons EDD-Algorithmus
\cite{Jackson1955}, vollständiger Dertouzos-Beweis der EDF-Optimalität
\cite{Dertouzos1974}, Bratley's Branch-and-Bound für nicht-preemptives Scheduling
\cite{Bratley1971} und Lawlers LDF-Algorithmus \cite{Lawler1973} für Tasks
mit Vorrangrelationen.

\textbf{Kapitel~4 -- Periodisches Task-Scheduling:} Vollständige Analyse von
Rate Monotonic (Liu \& Layland \cite{Liu1973}) einschließlich aller Beweise,
Hyperbolic Bound \cite{Bini2001}, Response Time Analysis \cite{Audsley1993},
EDF-Schedulierbarkeitstheorem \cite{Baruah1990} sowie der präzise Vergleich
zwischen schwacher (RM) und scharfer (EDF) Optimalität.

\textbf{Kapitel~5 -- Fixed-Priority Server:} Vollständige Analyse von
Background Scheduling, Polling Server, Deferrable Server (Strosnider et al.\
\cite{Strosnider1995}), Priority Exchange, Sporadic Server \cite{Sprunt1989}
und Slack Stealing \cite{Lehoczky1992} mit allen Beweisen und Laufzeitanalysen.

\textbf{Kapitel~6 -- Dynamic Priority Server:} EDF-basierte Server -- DPE,
DSS, TBS, CBS \cite{Abeni1998} und EDL-Server -- mit Optimalitätsnachweisen
und dem fundamentalen Nicht-Existenz-Ergebnis für optimale Server \cite{Tia1995}.

\textbf{Kapitel~7 -- Ressourcenzugriffsprotokolle:} Prioritätsinversion, NPP,
PIP und PCP \cite{Sha1990}, Stack Resource Policy \cite{Baker1991}, vollständige
Beweise der Blockierzeitschranken.

\textbf{Kapitel~8 -- Limited Preemptive Scheduling:} Nicht-preemptives Scheduling,
Preemption Thresholds \cite{Wang1999}, Deferred Preemptions und Task Splitting
mit dem optimalen SelectEPP-Algorithmus \cite{Bertogna2011}.

\textbf{Kapitel~9 -- Überlastbehandlung:} Value-Based Scheduling, Cervin's
Theorem über EDF bei permanenter Überlast \cite{Cervin2003}, Elastic Scheduling
und Admission Control.

\textbf{Kapitel~10 -- EDF$^+$-Algorithmus:} Vollständiger Optimalitätsbeweis,
Laufzeitanalyse, Korollare (Dominanz über RM, Minimax-Eigenschaft) und
Integration in AUTOSAR/ISO~26262 \cite{Epp2025}.

\textbf{Kapitel~11 -- Kernel-Design:} Interrupt-Behandlung, Speicherverwaltung,
Schedulingwarteschlangen und Kernel-Overhead in der Schedulierbarkeitsanalyse.

\textbf{Kapitel~12 -- Anwendungsdesign:} Hierarchische Dekomposition,
Robotersteuerungsbeispiel, AUTOSAR-Konfiguration.

\textbf{Kapitel~13 -- Zusammenfassung und Ausblick:} Gesamtüberblick,
praktische Empfehlungen und ein sehr ausführlicher Ausblick auf offene
Forschungsfragen -- Mehrprozessor-Scheduling, Mixed-Criticality, maschinelles
Lernen, formale Verifikation, Hardware-Beschleunigung und autonome Systeme.

\section*{Voraussetzungen und Zielgruppe}

Das Buch richtet sich an Studierende der Informatik und Elektrotechnik ab dem
dritten Semester sowie an Praktiker in der Embedded-Systems-Entwicklung.
Vorausgesetzt werden grundlegende Kenntnisse der Algorithmentheorie (O-Notation,
Induktionsbeweise) und der Betriebssystemkonzepte (Prozesse, Semaphore,
Scheduling). Für ein vollständiges Verständnis der Kapitel über Ressourcenprotokolle
und CBS sind Grundkenntnisse über Nebenläufigkeit hilfreich.

Mathematisch werden Kenntnisse der Analysis (Grenzwerte, Ableitungen) benötigt,
die für das Verständnis des Liu-Layland-Beweises (Kapitel~4) und des
Cervin-Theorems (Kapitel~9) notwendig sind.

\section*{Danksagung}

Der Autor dankt den Begründern der Echtzeit-Scheduling-Theorie -- insbesondere
C.\,L.\ Liu, J.\,W.\ Layland, Giorgio C.\ Buttazzo, Lui Sha und Sanjoy Baruah --
für ihre grundlegenden Beiträge, ohne die dieses Buch nicht möglich wäre.
Besonderer Dank gilt der Universität Bielefeld für die wissenschaftliche
Infrastruktur und Unterstützung.

\bigskip
\noindent Bielefeld, Juni 2026 \hfill \textit{Stephan Epp}

%=============================================================
\chapter{Allgemeine Betrachtung}
\label{chap:allgemein}
%=============================================================

\section{Einführung}

Echtzeit-Systeme sind Rechnersysteme, die innerhalb präziser Zeitschranken auf Ereignisse
in ihrer Umgebung reagieren müssen. Die Korrektheit solcher Systeme hängt nicht nur
vom Wert des Ergebnisses, sondern auch vom Zeitpunkt ab, zu dem das Ergebnis
erzeugt wird \cite{Stankovic1988}. Eine Reaktion, die zu spät erfolgt, kann wertlos
oder sogar gefährlich sein. Echtzeit-Computing spielt heute eine entscheidende Rolle
in unserer Gesellschaft, da eine wachsende Anzahl komplexer Systeme vollständig oder
teilweise auf Computersteuerung angewiesen ist.

Anwendungsbeispiele umfassen:
\begin{itemize}
	\item Steuerung chemischer und nuklearer Anlagen
	\item Steuerung komplexer Produktionsprozesse
	\item Eisenbahn-Weichensysteme
	\item Automotive-Anwendungen (AUTOSAR, ISO~26262)
	\item Flugregelungssysteme
	\item Umwelterfassung und -überwachung
	\item Telekommunikationssysteme
	\item Medizinische Systeme
	\item Industrieautomation und Robotik
	\item Militärsysteme und Raumfahrtmissionen
\end{itemize}

\section{Was bedeutet Echtzeit?}

\subsection{Der Begriff Zeit}

Das Hauptmerkmal, das Echtzeit-Computing von anderen Arten der Berechnung unterscheidet,
ist die Zeit. Das Wort \emph{Zeit} bedeutet, dass die Korrektheit des Systems nicht nur
vom logischen Ergebnis der Berechnung, sondern auch vom Zeitpunkt abhängt, zu dem
die Ergebnisse produziert werden.

Das Wort \emph{real} zeigt an, dass die Reaktion der Systeme auf externe Ereignisse
während deren Ablauf erfolgen muss. Folglich muss die Systemzeit (interne Zeit) mit
derselben Zeitskala gemessen werden, die zur Messung der Zeit in der kontrollierten
Umgebung (externe Zeit) verwendet wird.

\begin{definition}[Echtzeit-System]
	Ein \textbf{Echtzeit-System} ist ein Rechnersystem, das innerhalb präziser,
	vorher festgelegter Zeitgrenzen auf Ereignisse in seiner Umgebung reagieren muss.
	Die Korrektheit hängt sowohl vom Wert als auch vom Zeitpunkt der erzeugten
	Ergebnisse ab.
\end{definition}

\subsection{Typen von Echtzeit-Tasks}

Abhängig von den Konsequenzen einer versäumten Deadline unterscheidet man:

\begin{definition}[Harte Echtzeit-Task]
	Eine Echtzeit-Task heißt \textbf{hart}, wenn die Erzeugung von Ergebnissen nach
	ihrer Deadline katastrophale Folgen für das System unter Kontrolle haben kann.
\end{definition}

\begin{definition}[Weiche Echtzeit-Task]
	Eine Echtzeit-Task heißt \textbf{weich}, wenn die Erzeugung von Ergebnissen nach
	ihrer Deadline zwar eine Leistungsverschlechterung verursacht, aber keinen Schaden.
\end{definition}

\section{Vorhersagbarkeit erreichen}

Eine der wichtigsten Eigenschaften, die ein hartes Echtzeit-System haben muss,
ist Vorhersagbarkeit \cite{Stankovic1988}. Das System soll basierend auf den
Kernel-Merkmalen und den Informationen jeder Task die Entwicklung der Tasks vorhersagen
und im Voraus garantieren können, dass alle kritischen Zeitbedingungen eingehalten werden.

\subsection{Quellen der Nicht-Determinismus}

Die Vorhersagbarkeit eines Systems wird durch mehrere Faktoren beeinflusst:

\begin{itemize}
	\item \textbf{DMA}: Der Prozessor muss warten, wenn ein DMA-Gerät einen Speicherzyklus
	      stiehlt. Die Anzahl solcher Wartezeiten ist nicht vorhersagbar.
	\item \textbf{Cache-Speicher}: Cache-Fehlzugriffe (Cache Misses) führen zu variablen
	      Zugriffszeiten. Präemptionen verschlechtern die Cache-Lokalität erheblich.
	\item \textbf{Interrupts}: Können unbeschränkte Verzögerungen in die Task-Ausführung
	      einführen, wenn sie nicht ordnungsgemäß behandelt werden.
	\item \textbf{Systemaufrufe}: Unbeschränkte Ausführungszeiten von Kernelprimitiven
	      verhindern präzise WCET-Abschätzungen.
	\item \textbf{Semaphore und Prioritätsinversion}: Klassische Semaphore können
	      hochpriore Tasks durch niederpriore Tasks für unbeschränkte Zeit blockieren.
	\item \textbf{Speicherverwaltung}: Demand-Paging verursacht große und unvorhersagbare
	      Verzögerungen durch Seitenfehler.
\end{itemize}

%=============================================================
\chapter{Grundlegende Konzepte des Echtzeit-Schedulings}
\label{chap:basic}
%=============================================================

\section{Einführung}

Dieses Kapitel legt das vollständige begriffliche Fundament für alle
Scheduling-Algorithmen der folgenden Kapitel. Es definiert präzise die
Parameter von Tasks, die verschiedenen Typen von Zeitschranken, die formale
Klassifikation von Scheduling-Problemen nach Graham et al.\ sowie die
grundlegenden Gütekriterien. Besonderes Gewicht liegt auf den Konzepten
der \emph{Lateness}, \emph{Laxity} und \emph{Slack}, die die Basis für
die Optimalitätsbeweise in Kapitel~\ref{chap:aperiodisch} und
Kapitel~\ref{chap:periodisch} bilden.

\section{Task-Modell und Zeitparameter}

\subsection{Aperiodische Tasks -- Das Job-Modell}

\begin{definition}[Job (Aperiodische Task)]
	Ein \textbf{Job} $J_i$ ist eine atomare Ausführungseinheit mit folgenden
	Zeitparametern:
	\begin{itemize}
		\item $a_i$ (Arrival Time, Ankunftszeit): Zeitpunkt, ab dem $J_i$ zur
		      Ausführung bereit ist. Auch: \emph{release time} $r_i$.
		\item $C_i$ (Worst-Case Execution Time, WCET): Maximale Ausführungszeit
		      unter allen möglichen Eingaben. Unveränderlich für jede Instanz.
		\item $d_i$ (Absolute Deadline): Spätester Fertigstellungszeitpunkt.
		\item $D_i$ (Relative Deadline): $D_i = d_i - a_i$. Zeitdauer ab
		      Ankunft bis zur Deadline.
		\item $f_i$ (Finishing Time, Fertigstellungszeit): Tatsächlicher
		      Abschluss der Ausführung im konkreten Schedule.
		\item $s_i$ (Start Time, Startzeitpunkt): Erster Zeitpunkt, zu dem
		      $J_i$ tatsächlich ausgeführt wird.
	\end{itemize}
\end{definition}

\begin{definition}[Abgeleitete Zeitgrößen]
	Aus den Grundparametern werden folgende Größen abgeleitet:
	\begin{align*}
		R_i &= f_i - a_i \quad\text{(Response Time, Antwortzeit)}\\
		L_i &= f_i - d_i \quad\text{(Lateness: positiv = zu spät, negativ = Puffer)}\\
		E_i &= \max(0, f_i - d_i) \quad\text{(Tardiness: Verspätung, nie negativ)}\\
		X_i &= d_i - a_i - C_i \quad\text{(Laxity, Spielraum: maximale Wartezeit)}\\
		\text{slack}_i(t) &= d_i - t - c_i(t) \quad\text{(Slack zur Zeit }t\text{,
		mit verbleib. Rechenzeit }c_i(t))
	\end{align*}
\end{definition}

\begin{bemerkung}[Zusammenhänge der Zeitgrößen]
	Die Laxität $X_i = D_i - C_i$ gibt an, wie lange $J_i$ insgesamt warten
	darf, ohne ihre Deadline zu verpassen. Der Slack $\text{slack}_i(t)$ ist
	die verbleibende Wartekapazität zum Zeitpunkt $t$:
	\begin{itemize}
		\item $\text{slack}_i(t) > 0$: $J_i$ kann noch warten.
		\item $\text{slack}_i(t) = 0$: $J_i$ muss \emph{sofort} starten, sonst
		      Deadline-Verletzung.
		\item $\text{slack}_i(t) < 0$: Deadline bereits verpasst (bei Betrieb oder
		      bei unrealisierbaren Parametern).
	\end{itemize}
\end{bemerkung}

\begin{beispiel}[Zeitparameter]
	$J_1$: $a_1=2$, $C_1=3$, $d_1=8$.
	\begin{itemize}
		\item $D_1 = 8-2 = 6$, $X_1 = 6-3 = 3$ (Laxität: darf 3 Einheiten warten)
		\item Wenn $s_1=4$, $f_1=7$: $R_1=5$, $L_1=-1$ (1 Einheit Puffer), $E_1=0$
		\item Wenn $s_1=6$, $f_1=9$: $R_1=7$, $L_1=1$ (1 Einheit zu spät), $E_1=1$
		\item $\text{slack}_1(4) = 8-4-3 = 1$ (kann noch 1 Einheit warten)
		\item $\text{slack}_1(5) = 8-5-3 = 0$ (muss sofort starten)
	\end{itemize}
\end{beispiel}

\subsection{Periodische Tasks -- Das $\tau$-Modell}

\begin{definition}[Periodische Task]
	Eine \textbf{periodische Task} $\tau_i$ erzeugt in regelmäßigen Abständen
	Jobs. Sie ist durch folgende Parameter vollständig beschrieben:
	\begin{itemize}
		\item $\Phi_i$: Phase (Ankunftszeit der ersten Instanz)
		\item $T_i$: Periode (Abstand zwischen aufeinanderfolgenden Instanzen)
		\item $C_i$: WCET jeder Instanz (konstant)
		\item $D_i$: Relative Deadline ($D_i \leq T_i$ im Allgemeinen;
		      $D_i = T_i$ unter Annahme A3)
	\end{itemize}
	Die $k$-te Instanz $\tau_{i,k}$ hat:
	\[
		r_{i,k} = \Phi_i + (k-1)T_i, \qquad
		d_{i,k} = r_{i,k} + D_i = \Phi_i + (k-1)T_i + D_i.
	\]
\end{definition}

\begin{definition}[Hyperperiode]
	Die \textbf{Hyperperiode} $H$ ist das kleinste gemeinsame Vielfache aller
	Perioden: $H = \mathrm{lcm}(T_1, \ldots, T_n)$.
	Nach der Hyperperiode wiederholt sich der Schedule exakt (bei synchronen
	Aktivierungen und $D_i = T_i$).
\end{definition}

\begin{definition}[Prozessorauslastung]
	\label{def:util}
	Die \textbf{Prozessorauslastung} $U$ einer Menge von $n$ periodischen Tasks ist:
	\[
		U = \sum_{i=1}^{n} \frac{C_i}{T_i}.
	\]
	$U$ gibt den Anteil der Prozessorzeit an, der für die periodischen Tasks
	\emph{im Durchschnitt} benötigt wird. Wenn $U > 1$, ist kein machbarer
	Schedule unter irgendeinem Algorithmus möglich.
\end{definition}

\begin{definition}[Synchroner vs. asynchroner Task-Set]
	\begin{itemize}
		\item \textbf{Synchron}: Alle Tasks starten gleichzeitig bei $t=0$:
		      $\Phi_i = 0$ für alle $i$.
		\item \textbf{Asynchron}: Tasks haben unterschiedliche Phasen $\Phi_i$.
	\end{itemize}
	Die kritische Analyse erfolgt stets für synchrone Aktivierungen (kritischer
	Augenblick -- schlimmster Fall für RM, vgl.\ Kapitel~\ref{chap:periodisch}).
\end{definition}

\subsection{Sporadische und aperiodische Tasks}

\begin{definition}[Sporadische Task]
	Eine \textbf{sporadische Task} ist eine Task mit unregelmäßigen Aktivierungen,
	jedoch mit einer Mindestankunftsabstand (Minimum Interarrival Time) $T_i^{\min}$:
	$a_{i,k+1} - a_{i,k} \geq T_i^{\min}$. Sporadische Tasks können offline
	garantiert werden, indem man $T_i^{\min}$ als Worst-Case-Periode verwendet.
\end{definition}

\begin{definition}[Firme aperiodische Task]
	Eine \textbf{firme} aperiodische Task hat eine Deadline. Sie wird online
	garantiert: Beim Eintreffen prüft ein Admission-Controller, ob sie ohne
	Deadline-Verletzung anderer Tasks integriert werden kann.
\end{definition}

\section{Klassifikation von Scheduling-Problemen}

\subsection{Graham-Notation}

\begin{definition}[Graham-Notation $\alpha|\beta|\gamma$]
	Jedes Scheduling-Problem wird durch drei Felder beschrieben:
	\begin{itemize}
		\item $\alpha$: \textbf{Maschinenmodell} -- z.\,B.\ $1$ (Einzelprozessor),
		      $P_m$ ($m$ identische Prozessoren), $Q_m$ (uniforme Geschwindigkeiten),
		      $R_m$ (unverwandte Prozessoren).
		\item $\beta$: \textbf{Task-Constraints} -- z.\,B.\ \texttt{preem} (preemptiv),
		      \texttt{no-preem} (nicht-preemptiv), \texttt{prec} (Vorrangrelationen),
		      \texttt{sync} (synchrone Ankunft), $D_i = T_i$ (Deadline gleich Periode).
		\item $\gamma$: \textbf{Gütekriterium} -- z.\,B.\ $L_{\max}$ (maximale Lateness),
		      $\sum L_i$ (Summe der Latenesses), $E_{\max}$ (maximale Tardiness),
		      \texttt{feasible} (machbarer Schedule).
	\end{itemize}
\end{definition}

\begin{beispiel}[Scheduling-Notationen der behandelten Probleme]
	\begin{itemize}
		\item $1 \mid \text{sync} \mid L_{\max}$: Einzelprozessor, synchrone
		      Ankunft, minimiere maximale Lateness $\to$ \textbf{EDD/Jackson} (Kap.~\ref{chap:aperiodisch})
		\item $1 \mid \text{preem} \mid L_{\max}$: Einzelprozessor, preemptiv,
		      beliebige Ankunftszeiten, minimiere $L_{\max}$ $\to$ \textbf{EDF/Horn} (Kap.~\ref{chap:aperiodisch})
		\item $1 \mid \text{preem}, D_i=T_i \mid \text{feasible}$: Periodisch,
		      preemptiv $\to$ \textbf{RM, EDF} (Kap.~\ref{chap:periodisch})
		\item $1 \mid \text{no-preem} \mid \text{feasible}$: Nicht-preemptiv,
		      NP-schwer $\to$ \textbf{Bratley} (Kap.~\ref{chap:aperiodisch})
	\end{itemize}
\end{beispiel}

\subsection{Gütekriterien}

\begin{definition}[Gütekriterien im Überblick]
	\begin{itemize}
		\item $L_{\max} = \max_i L_i$: \textbf{Maximale Lateness}. Negativ $\Rightarrow$
		      alle Deadlines eingehalten. Optimalitätsmaß für EDF und EDD.
		\item $E_{\max} = \max_i E_i$: \textbf{Maximale Tardiness}.
		\item $\sum L_i$: \textbf{Summe der Latenesses}. Equivalent zur Minimierung
		      der Gesamtfertigstellungszeit ($\sum f_i$).
		\item $n_{\text{late}} = |\{J_i \mid f_i > d_i\}|$: \textbf{Anzahl verspäteter
		      Tasks}. NP-schwer für beliebige Ankunftszeiten.
		\item \textbf{Feasibility}: Existenz eines Schedules mit $L_{\max} \leq 0$.
	\end{itemize}
\end{definition}

\begin{bemerkung}[Zusammenhang Gütekriterien]
	$L_{\max} \leq 0 \iff$ Schedule ist machbar $\iff E_{\max} = 0$.
	Ein Algorithmus, der $L_{\max}$ minimiert, findet automatisch einen
	machbaren Schedule, wenn einer existiert. Deshalb ist die Minimierung
	von $L_{\max}$ das stärkere und nützlichere Kriterium gegenüber
	bloßer Machbarkeit.
\end{bemerkung}

\section{Klassifikation von Scheduling-Algorithmen}

\begin{definition}[Taxonomie der Scheduling-Algorithmen]
	\begin{itemize}
		\item \textbf{Preemptiv vs. Nicht-Preemptiv}: Darf eine laufende Task
		      unterbrochen werden?
		\item \textbf{Statisch vs. Dynamisch}: Werden Prioritäten vor oder
		      während der Ausführung berechnet?
		\item \textbf{Online vs. Offline}: Sind alle Task-Parameter von Anfang
		      an bekannt (offline) oder werden Tasks dynamisch aktiviert (online)?
		\item \textbf{Optimal vs. Heuristisch}: Findet der Algorithmus stets
		      die beste Lösung, oder nur eine gute Näherung?
		\item \textbf{Clairvoyant vs. Non-Clairvoyant}: Hat der Algorithmus
		      vollständige Kenntnis aller zukünftigen Ankunftszeiten?
	\end{itemize}
\end{definition}

\begin{definition}[Optimalitätsbegriffe -- Zusammenfassung]
	\label{def:optimalitaet}
	\begin{itemize}
		\item \textbf{Schwach optimal} (innerhalb einer Klasse): Findet den
		      besten Schedule unter allen Algorithmen derselben Klasse.
		      Beispiel: RM ist schwach optimal unter festen Prioritäten.
		\item \textbf{Scharf optimal} (global): Findet jeden machbaren Schedule,
		      der überhaupt existiert. Beispiel: EDF ist scharf optimal für
		      preemptive Einzelprozessoren.
		\item \textbf{Optimal bezüglich $\gamma$}: Minimiert das Gütekriterium
		      $\gamma$ unter allen möglichen Schedules.
	\end{itemize}
\end{definition}

\section{Vorrangrelationen und Ressourcenabhängigkeiten}

\begin{definition}[Vorranggraph (Precedence Graph)]
	Ein \textbf{Vorranggraph} $G = (J, \prec)$ ist ein gerichteter azyklischer
	Graph (DAG), in dem $J_i \prec J_j$ bedeutet: $J_j$ darf erst starten,
	wenn $J_i$ vollständig abgeschlossen ist. Notationen:
	\begin{itemize}
		\item $J_i \to J_j$: $J_i$ ist unmittelbarer Vorgänger von $J_j$
		\item $J_i \prec J_j$: $J_i$ ist (direkter oder indirekter) Vorgänger
	\end{itemize}
\end{definition}

\begin{definition}[Ressourcenkonflikt]
	Zwei Tasks $J_i$ und $J_j$ stehen in einem \textbf{Ressourcenkonflikt},
	wenn sie dieselbe exklusive Ressource $R_k$ benötigen. Ressourcenkonflikte
	führen zu Blockierung und sind der Hauptgrund für Prioritätsinversion
	(Kapitel~\ref{chap:ressourcen}).
\end{definition}

\section{Scheduling-Anomalien -- Formale Analyse}

\begin{definition}[Scheduling-Anomalie nach Graham]
	Eine \textbf{Scheduling-Anomalie} tritt auf, wenn eine Verbesserung der
	Eingabeparameter -- weniger Tasks, kürzere WCETs, schwächere Vorrangrelationen,
	mehr Prozessoren -- zu einem \emph{schlechteren} Schedule-Ergebnis (z.\,B.\
	längerer Makespan) führt.
\end{definition}

\begin{beispiel}[Graham-Anomalie auf Mehrprozessorsystemen]
	Betrachte 6 Tasks $J_1,\ldots,J_6$ mit WCETs $C_i = 3$ und Vorrangrelation
	$J_1 \prec J_3$, $J_1 \prec J_4$, $J_2 \prec J_5$, $J_2 \prec J_6$ auf
	3 Prozessoren. List Scheduling (LPT) erzeugt Makespan 3. Wenn die Vorrangrelation
	entfernt wird, kann List Scheduling Makespan 4 produzieren -- trotz weniger
	Einschränkungen. Graham (1969) bewies: Solche Anomalien sind bei allgemeinen
	Mehrprozessorsystemen unter List Scheduling unvermeidbar.
\end{beispiel}

\begin{beispiel}[Anomalie bei Selbstsuspension]
	Task $\tau_1$ mit $C_1=1$, $T_1=4$, Deadline $d_1=4$ hat Slack~3. Führt
	$\tau_1$ eine Selbstsuspension von 2 Einheiten aus, verpasst sie ihre Deadline.
	Obwohl $\tau_1$ selbst weniger Rechenzeit benötigt (wegen Suspension), wird
	das System schlechter -- ein direktes Beispiel für eine Anomalie bei
	Selbstsuspensionen unter RM.
\end{beispiel}

\begin{bemerkung}[Anomalien und Konsequenzen für den Systementwurf]
	Scheduling-Anomalien haben direkte praktische Konsequenzen:
	\begin{enumerate}
		\item \textbf{WCET-Erhöhung} einer Task kann durch zusätzliche Preemptionen
		      zu Deadline-Verletzungen führen, obwohl das System ursprünglich
		      schedulierbar war.
		\item \textbf{Selbstsuspensionen} (z.\,B.\ für I/O) müssen explizit in
		      der Schedulierbarkeitsanalyse berücksichtigt werden.
		\item \textbf{Mehrprozessorsysteme}: Globale Scheduling-Anomalien
		      (Dhall-Effekt) zeigen, dass EDF auf Mehrprozessoren nicht mehr
		      optimal ist.
	\end{enumerate}
\end{bemerkung}

\section{Der Dhall-Effekt -- EDF auf Mehrprozessoren}

\begin{theorem}[Dhall-Effekt, 1978]
	EDF ist auf Mehrprozessorsystemen \emph{nicht optimal}: Es gibt Task-Sets
	mit Auslastung beliebig nah an $1$, die EDF auf $m$ Prozessoren nicht
	schedulieren kann.
\end{theorem}

\begin{proof}[Beweisidee]
	Betrachte $m+1$ Tasks mit $C_i = \varepsilon$, $T_i = \varepsilon + \delta$
	für $i=1,\ldots,m$ (kurze Tasks, frühe Deadlines) und $C_{m+1} = 1$,
	$T_{m+1} = 1+\varepsilon$ (lange Task, spätere Deadline). Auslastung:
	$U = m \cdot \varepsilon/(\varepsilon+\delta) + 1/(1+\varepsilon) \to 1$
	für kleine $\varepsilon, \delta$. EDF belegt alle $m$ Prozessoren mit
	den $m$ kurzperiodischen Tasks, da diese immer die früheste Deadline haben,
	und lässt die lange Task verhungern. \hfill$\blacksquare$
\end{proof}

\begin{bemerkung}
	Der Dhall-Effekt zeigt, dass der Wechsel von Uniprocessor- zu
	Multiprocessor-Scheduling grundlegend neue Algorithmen erfordert:
	Partitioned Scheduling, Global EDF mit Deadline-Adjustment (z.\,B.\
	GLLF, Pfair), oder Clustered Scheduling. Diese werden in diesem Buch
	nicht behandelt, sind aber aktives Forschungsgebiet.
\end{bemerkung}

\section{Vollständige Notation und Annahmen}

Zur Vereinheitlichung gelten in diesem Buch folgende Standardannahmen,
sofern nicht ausdrücklich anders angegeben:

\begin{itemize}
	\item[\textbf{A1}] Alle periodischen Task-Instanzen haben dieselbe WCET $C_i$.
	\item[\textbf{A2}] Die Perioden $T_i$ sind konstant.
	\item[\textbf{A3}] Relative Deadlines sind gleich den Perioden: $D_i = T_i$.
	\item[\textbf{A4}] Tasks sind unabhängig (keine Vorrangrelationen, keine gemeinsamen Ressourcen).
	\item[\textbf{A5}] Tasks suspendieren sich nicht (kein $I/O$-Warten).
	\item[\textbf{A6}] Tasks werden sofort bei ihrer Ankunft freigegeben.
	\item[\textbf{A7}] Kernel-Overhead ist null.
\end{itemize}

Annahmen A3 und A4 werden in späteren Kapiteln gezielt aufgehoben:
A3 in Kapitel~\ref{chap:periodisch} (Deadline Monotonic, EDF mit $D_i < T_i$),
A4 in Kapitel~\ref{chap:ressourcen} (Ressourcenzugriffsprotokolle).

\begin{table}[H]
	\centering
	\caption{Übersicht: Zeitparameter periodischer Tasks}
	\begin{tabular}{cll}
		\toprule
		\textbf{Symbol} & \textbf{Bezeichnung} & \textbf{Definition} \\
		\midrule
		$\Phi_i$ & Phase & $r_{i,1}$ (Ankunft der 1. Instanz) \\
		$T_i$ & Periode & Abstand zwischen Instanzen \\
		$C_i$ & WCET & Max.\ Ausführungszeit pro Instanz \\
		$D_i$ & Relative Deadline & $\leq T_i$ \\
		$r_{i,k}$ & Release Time (Instanz $k$) & $\Phi_i + (k-1)T_i$ \\
		$d_{i,k}$ & Absolute Deadline (Instanz $k$) & $r_{i,k} + D_i$ \\
		$U_i$ & Auslastung & $C_i/T_i$ \\
		$U$ & Gesamtauslastung & $\sum_i C_i/T_i$ \\
		$H$ & Hyperperiode & $\mathrm{lcm}(T_1,\ldots,T_n)$ \\
		$R_i$ & Antwortzeit & $f_i - r_i$ \\
		$L_i$ & Lateness & $f_i - d_i$ \\
		$X_i$ & Laxität & $D_i - C_i$ \\
		\bottomrule
	\end{tabular}
\end{table}

%=============================================================
\chapter{Aperiodisches Task-Scheduling}
\label{chap:aperiodisch}
%=============================================================

\section{Einführung}

Dieses Kapitel behandelt Algorithmen zur Planung aperiodischer Echtzeit-Tasks
auf einem Einzelprozessorsystem. Jeder Algorithmus stellt eine Lösung für ein
bestimmtes Scheduling-Problem dar.

Der fundamentale Kernsatz gilt bereits für aperiodische Tasks:

\begin{quote}
\textbf{Wenn es einen machbaren Schedule gibt, findet EDF diesen sofort. EDF ist scharf optimal.}
\end{quote}

Dies wird im folgenden streng bewiesen.

\section{Jacksons Algorithmus (EDD)}

Das durch Jackson gelöste Problem ist $1 \mid \text{sync} \mid L_{\max}$: Eine Menge
$J$ von $n$ aperiodischen Tasks soll auf einem einzelnen Prozessor eingeplant werden,
wobei die maximale Lateness minimiert wird. Alle Tasks haben synchrone Ankunftszeiten.

\begin{theorem}[Jackson 1955]
	\label{thm:jackson}
	Gegeben sei eine Menge von $n$ unabhängigen Tasks. Jeder Algorithmus, der die
	Tasks in Reihenfolge nichtabsteigender Deadlines ausführt, ist optimal bezüglich
	der Minimierung der maximalen Lateness.
\end{theorem}

\begin{proof}
	Der Beweis erfolgt durch ein Austauschargument. Sei $\sigma$ ein Schedule, der
	von einem beliebigen Algorithmus $A$ produziert wird. Wenn $A \neq \text{EDD}$,
	existieren zwei Tasks $J_a$ und $J_b$ mit $d_a \leq d_b$, sodass $J_b$ in
	$\sigma$ unmittelbar vor $J_a$ läuft. Sei $\sigma'$ ein Schedule, der durch
	Vertauschen von $J_a$ und $J_b$ in $\sigma$ entsteht. Das Vertauschen kann die
	maximale Lateness nicht erhöhen: In beiden Fällen ($L'_a \geq L'_b$ oder
	$L'_a \leq L'_b$) gilt $L'_{\max}(a,b) < L_{\max}(a,b)$. Durch endlich viele
	solcher Transpositionen wird $\sigma$ in $\sigma_{\text{EDD}}$ umgewandelt, der
	optimal ist. \hfill$\blacksquare$
\end{proof}

\section{Horns Algorithmus (EDF) -- Scharf optimal}

\subsection{Das EDF-Optimalitätstheorem}

Wenn Tasks dynamische Ankunftszeiten haben, liefert das Earliest Deadline First
Prinzip die Grundlage. Der fundamentale Satz lautet:

\begin{quote}
\textbf{Wenn es einen machbaren Schedule gibt, findet EDF diesen sofort. Die kürzeste
Ausführungszeit aller Tasks findet EDF sofort. EDF ist scharf optimal.}
\end{quote}

\begin{theorem}[Horn 1974 / Dertouzos 1974 -- EDF Optimalität]
	\label{thm:edf_optimal}
	Gegeben sei eine Menge von $n$ unabhängigen Tasks mit beliebigen Ankunftszeiten.
	Jeder Algorithmus, der zu jedem Zeitpunkt die Task mit der frühesten absoluten
	Deadline unter allen bereiten Tasks ausführt, ist optimal bezüglich der
	Minimierung der maximalen Lateness. Insbesondere gilt:
	\[
		\text{Wenn ein machbarer Schedule existiert, findet EDF diesen sofort.}
	\]
\end{theorem}

\begin{proof}[Beweis (Dertouzos 1974)]
	Sei $\sigma$ der Schedule eines beliebigen Algorithmus $A$ und $\sigma_{\text{EDF}}$
	der EDF-Schedule. Da Preemption erlaubt ist, kann jede Task in getrennten
	Zeitscheiben ausgeführt werden.

	\textbf{Notationen:}
	\begin{itemize}
		\item $\sigma(t)$: Task, die in der Scheibe $[t, t+1)$ ausgeführt wird.
		\item $E(t)$: Bereite Task, die zum Zeitpunkt $t$ die früheste Deadline hat.
		\item $t_E(t)$: Zeitpunkt ($\geq t$), zu dem die nächste Scheibe von $E(t)$
		      in $\sigma$ ausgeführt wird.
	\end{itemize}

	Wenn $\sigma \neq \sigma_{\text{EDF}}$, existiert ein Zeitpunkt $t$ mit
	$\sigma(t) \neq E(t)$. Das Vertauschen der Scheiben bei $t$ und $t_E$ kann die
	maximale Lateness nicht erhöhen: Wird eine Scheibe vorgezogen, bleibt die
	Machbarkeit offensichtlich erhalten. Wird eine Scheibe von $J_i$ auf $t_E$
	verschoben und ist $\sigma$ machbar, so gilt $(t_E + 1) \leq d_E$. Da
	$d_E \leq d_i$ für alle $i$, gilt $t_E + 1 \leq d_i$, was die Machbarkeit
	der verschobenen Scheibe garantiert.

	Durch Anwendung des Vertauschungsalgorithmus auf alle Zeitscheiben von $t=0$
	bis $t=D$ (dem spätesten Deadline-Zeitpunkt) wird $\sigma$ in $\sigma_{\text{EDF}}$
	umgewandelt, ohne die maximale Lateness zu erhöhen. Da EDF die maximale
	Lateness minimiert, findet EDF jeden machbaren Schedule. \hfill$\blacksquare$
\end{proof}

\subsection{Warum EDF scharf optimal ist}

\begin{definition}[Scharf optimal]
	Ein Scheduling-Algorithmus $\mathcal{A}$ heißt \textbf{scharf optimal}, wenn er
	für jede Task-Klasse $\mathfrak{S}$ gilt:
	\[
		\forall\, \Sigma \in \mathfrak{S}: \quad
		\Sigma \text{ ist schedulierbar} \iff \mathcal{A} \text{ findet einen machbaren Schedule.}
	\]
	Die Bedingung ist eine Äquivalenz, nicht nur eine Implikation.
\end{definition}

EDF ist scharf optimal, weil:
\begin{enumerate}
	\item EDF findet jeden machbaren Schedule (Satz~\ref{thm:edf_optimal}).
	\item Wenn EDF keinen machbaren Schedule findet, existiert keiner -- EDF scheitert
	      genau dann, wenn die Auslastung $U > 1$ ist oder die Deadline-Bedingungen
	      strukturell nicht erfüllbar sind.
	\item EDF minimiert $L_{\max}$ unter allen Algorithmen.
\end{enumerate}

Die Schärfe bedeutet: EDF nutzt die gesamte verfügbare Prozessorkapazität aus.
Die kürzeste Ausführungszeit aller Tasks findet EDF sofort, da EDF immer die Task
mit der frühesten Deadline priorisiert -- also diejenige, die am dringlichsten ist.

\subsection{Garantiebedingung für EDF}

Wenn Tasks dynamische Aktivierungszeiten haben, muss der Garantietest
dynamisch bei jeder neuen Task-Ankunft durchgeführt werden:

\begin{satz}[EDF-Garantietest]
	Sei $J$ die aktuell garantierte Task-Menge. Eine neu ankommende Task $J_{\text{new}}$
	kann garantiert werden, wenn für die Vereinigung $J' = J \cup \{J_{\text{new}}\}$
	gilt:
	\[
		\forall i = 1, \ldots, n: \quad \sum_{k=1}^{i} c_k(t) \leq d_i,
	\]
	wobei $c_k(t)$ die verbleibende Ausführungszeit von $J_k$ zum Zeitpunkt $t$ und
	$d_i$ die Deadline der $i$-ten Task (nach Deadline sortiert) ist.
\end{satz}

\section{Nicht-preemptives Scheduling}

Wenn Preemption nicht erlaubt ist und Tasks beliebige Ankunftszeiten haben,
werden Minimierung der maximalen Lateness und Suche nach einem machbaren Schedule
NP-schwer \cite{Buttazzo2011}.

EDF ist in diesem Modell \textbf{nicht mehr optimal}: Es kann Fälle geben, in
denen ein machbarer Schedule existiert, EDF diesen aber nicht findet. Dies
unterstreicht die Bedeutung der Preemptivität für die Optimalität von EDF.

\subsection{Bratley's Algorithmus}

Bratley's Algorithmus (1971) löst das nicht-preemptive Scheduling durch
Branch-and-Bound-Suche. Er verwendet Beschneidungstechniken:
\begin{itemize}
	\item Ein Ast wird abgebrochen, wenn jede Erweiterung des aktuellen Pfades
	      zu einer versäumten Deadline führt.
	\item Ein Ast wird abgebrochen, wenn ein machbarer Schedule gefunden wurde.
\end{itemize}
Die Komplexität ist $O(n \cdot n!)$ im schlimmsten Fall.

\section{Scheduling mit Vorrangrelationen}

\begin{theorem}[Lawler 1973 -- LDF Optimalität]
	Der Latest Deadline First (LDF) Algorithmus minimiert die maximale Lateness
	für eine Menge von Tasks mit Vorrangrelationen und synchronen Ankunftszeiten.
\end{theorem}

Für Tasks mit Vorrangrelationen und dynamischen Aktivierungszeiten existiert der
EDF*-Algorithmus von Chetto et al. (1990), der eine Transformation der
Zeitparameter durchführt, um EDF anwendbar zu machen.

%=============================================================
\chapter{Periodisches Task-Scheduling}
\label{chap:periodisch}
%=============================================================

\section{Einführung}

In vielen Echtzeit-Kontrollanwendungen stellen periodische Aktivitäten den
Hauptrechenaufwand dar. Dieses Kapitel behandelt vier grundlegende Algorithmen:
Timeline Scheduling, Rate Monotonic, Earliest Deadline First und Deadline Monotonic.

\textbf{Der zentrale Vergleich dieses Kapitels:}

\begin{quote}
\textbf{Wenn es einen machbaren Schedule gibt, findet RM diesen sofort. RM ist schwach optimal.}

\textbf{Wenn es einen machbaren Schedule gibt, findet EDF diesen sofort. EDF ist scharf optimal.}
\end{quote}

Diese Aussage wird durch die Schedulierbarkeitsanalyse präzisiert: RM kann nur
Auslastungen bis $\ln 2 \approx 69\%$ im schlimmsten Fall garantieren, während
EDF die volle Prozessorkapazität bis $U = 1$ ausschöpft.

\section{Notation und Annahmen}

\begin{itemize}
	\item $\Gamma$: Menge periodischer Tasks
	\item $\tau_i$: Generische periodische Task
	\item $T_i$: Periode von $\tau_i$
	\item $C_i$: Worst-Case Ausführungszeit von $\tau_i$
	\item $D_i = T_i$: Relative Deadline (Assumption A3)
	\item $U = \sum_{i=1}^{n} C_i/T_i$: Prozessorauslastung
\end{itemize}

Annahmen: A1 (konstante Periode), A2 (konstante WCET), A3 ($D_i = T_i$),
A4 (unabhängige Tasks).

\section{Timeline Scheduling}

Timeline Scheduling (Cyclic Executive) unterteilt die Zeitachse in gleich
lange Slots und weist Tasks deterministisch zu. Der \textbf{Minor Cycle} ist die
Länge eines Slots (GCD der Perioden), der \textbf{Major Cycle} ist der
Hyperperiod (LCM der Perioden).

\textbf{Vorteile:} Einfach zu implementieren, kein Overhead durch Context Switches,
deterministische Ausführungsreihenfolge.

\textbf{Nachteile:} Spröde bei Überlastbedingungen, sensibel gegenüber
Anwendungsänderungen, schwierig für aperiodische Tasks.

\section{Rate Monotonic Scheduling -- Detaillierte Analyse}

\subsection{Einführung und Grundidee}

Der Rate Monotonic (RM) Scheduling-Algorithmus wurde 1973 von Liu und Layland
eingeführt \cite{Liu1973} und gehört heute zu den meistgenutzten Algorithmen für
periodische Echtzeit-Tasks in eingebetteten Systemen. Die Grundidee ist denkbar
einfach: Tasks werden statische Prioritäten \emph{umgekehrt proportional} zu ihrer
Periode zugewiesen. Eine Task mit kurzer Periode (also hoher Anforderungsrate) erhält
eine höhere Priorität als eine Task mit langer Periode.

\begin{definition}[Rate Monotonic Prioritätszuweisung]
	Der \textbf{Rate Monotonic} Algorithmus weist jeder periodischen Task $\tau_i$
	eine Priorität $P_i$ zu, die umgekehrt proportional zur Periode ist:
	\[
		T_i < T_j \implies P_i > P_j.
	\]
	Tasks mit gleicher Periode erhalten beliebige, aber feste Prioritäten.
\end{definition}

Da die Perioden konstant sind, ist RM eine \textbf{feste Prioritätszuweisung}: Die
Priorität einer Task ändert sich während der Ausführung nicht. RM ist außerdem
\textbf{preemptiv}: Wird eine neue Instanz einer hochprioren Task aktiv, unterbricht
sie die laufende niederpriore Task sofort.

\begin{beispiel}[Einfache RM-Zuweisung]
	Gegeben sei ein Task-Set mit drei Tasks:
	\[
		\tau_1: T_1 = 4\,\text{ms}, \quad \tau_2: T_2 = 6\,\text{ms}, \quad \tau_3: T_3 = 12\,\text{ms}.
	\]
	RM ordnet die Prioritäten: $P_1 > P_2 > P_3$, da $T_1 < T_2 < T_3$.
	$\tau_1$ ist immer die dringlichste Task und kann $\tau_2$ und $\tau_3$ jederzeit
	unterbrechen. $\tau_2$ kann nur $\tau_3$ unterbrechen.
\end{beispiel}

\subsection{Kritischer Augenblick -- Das Fundament der RM-Analyse}

Das wichtigste Konzept für die Analyse von RM ist der \emph{kritische Augenblick}
(critical instant). Er gibt an, unter welchen Bedingungen eine Task die größte
Antwortzeit erfährt.

\begin{definition}[Kritischer Augenblick]
	Der \textbf{kritische Augenblick} einer Task $\tau_i$ ist derjenige
	Ankunftszeitpunkt, der die größte Antwortzeit von $\tau_i$ produziert.
	Die zugehörige \textbf{kritische Zeitzone} ist das Intervall zwischen dem
	kritischen Augenblick und der entsprechenden Antwortzeit.
\end{definition}

\begin{lemma}[Liu \& Layland 1973 -- Kritischer Augenblick]
	\label{lem:critical_instant}
	Der kritische Augenblick für eine Task $\tau_i$ unter RM tritt auf, wenn
	$\tau_i$ \emph{gleichzeitig} mit allen Tasks höherer Priorität aktiviert wird.
\end{lemma}

\begin{proof}
	Wir zeigen durch sukzessives Vorrücken der Ankunftszeitpunkte hochpriorer Tasks,
	dass die Antwortzeit von $\tau_i$ maximal wird, wenn alle gleichzeitig ankommen.

	Sei $\tau_n$ die Task mit der niedrigsten Priorität und $\tau_i$ eine Task mit
	höherer Priorität. Die Antwortzeit von $\tau_n$ wird durch die Interferenz von
	$\tau_i$ verzögert. Wie in Abbildung (a) dargestellt, verzögert eine Instanz von
	$\tau_i$ die Ausführung von $\tau_n$ um $C_i$. Das Vorziehen der Ankunftszeit von
	$\tau_i$ kann die Fertigstellungszeit von $\tau_n$ weiter erhöhen (Abbildung b),
	da eine zusätzliche Instanz von $\tau_i$ in das Interferenzintervall fallen kann.

	Durch wiederholtes Vorrücken aller hochprioren Tasks wird die Antwortzeit von
	$\tau_n$ maximiert, wenn alle gleichzeitig ankommen. Dieses Argument gilt für
	jede Task $\tau_n$ mit niedrigster Priorität, und durch Rekursion für jede
	Task im System. \hfill$\blacksquare$
\end{proof}

\begin{bemerkung}
	Lemma~\ref{lem:critical_instant} hat eine weitreichende praktische Konsequenz:
	Um zu prüfen, ob ein Task-Set unter RM schedulierbar ist, genügt es, den
	Schedule am kritischen Augenblick zu analysieren -- also den Fall, dass alle
	Tasks zum Zeitpunkt $t = 0$ gleichzeitig starten. Wenn alle Tasks in diesem
	ungünstigsten Fall ihre Deadlines einhalten, sind sie \emph{unter allen}
	möglichen Ankunftsbedingungen schedulierbar.
\end{bemerkung}

\subsection{Vollständiger Optimalitätsbeweis von RM}

\begin{definition}[Schwach optimal]
	\label{def:schwach_optimal}
	Ein Scheduling-Algorithmus $\mathcal{A}$ heißt \textbf{schwach optimal}
	innerhalb einer Klasse $\mathfrak{A}$ von Algorithmen, wenn gilt:
	\[
		\forall\, A \in \mathfrak{A}, \forall\, \Gamma: \quad
		A \text{ scheduliert } \Gamma \;\implies\; \mathcal{A} \text{ scheduliert } \Gamma.
	\]
\end{definition}

\begin{theorem}[Liu \& Layland 1973 -- RM Optimalität]
	\label{thm:rm_optimal}
	RM ist \textbf{schwach optimal} unter allen Algorithmen mit fester Priorität:
	Wenn ein Task-Set $\Gamma$ von \emph{irgendeiner} festen Prioritätszuweisung
	scheduliert werden kann, so kann es auch von RM scheduliert werden.
\end{theorem}

\begin{proof}
	Der Beweis wird durch ein direktes Austauschargument für zwei Tasks geführt
	und anschließend auf $n$ Tasks verallgemeinert.

	\textbf{Schritt 1: Zwei Tasks, $T_1 < T_2$.}

	Angenommen, die Prioritäten seien \emph{nicht} gemäß RM zugewiesen, d.\,h.\
	$\tau_2$ hat eine höhere Priorität als $\tau_1$. Am kritischen Augenblick
	($t = 0$) sieht der Schedule aus wie in Abbildung~\ref{fig:rm_non_rm_priority}:
	$\tau_2$ wird zuerst ausgeführt, dann $\tau_1$. Für Machbarkeit muss gelten:
	\begin{equation}
		C_1 + C_2 \leq T_1.
		\label{eq:non_rm_feasible}
	\end{equation}

	Nun weisen wir die Prioritäten gemäß RM zu, d.\,h.\ $\tau_1$ erhält die höhere
	Priorität ($T_1 < T_2$). Wir müssen zeigen, dass die Machbarkeitsbedingung
	(\ref{eq:non_rm_feasible}) auch unter RM ausreicht. Dazu unterscheiden wir
	zwei Fälle, je nachdem, wie $C_1$ und $T_2$ sich zueinander verhalten. Sei
	$F = \lfloor T_2 / T_1 \rfloor$ die Anzahl vollständiger $T_1$-Perioden
	innerhalb von $T_2$.

	\textbf{Fall 1:} $C_1 < T_2 - F T_1$.

	In diesem Fall werden alle $F+1$ Instanzen von $\tau_1$ abgeschlossen, bevor
	die zweite Instanz von $\tau_2$ ihre Deadline erreicht. Die Machbarkeitsbedingung
	für $\tau_2$ unter RM ist:
	\begin{equation}
		(F+1)C_1 + C_2 \leq T_2.
		\label{eq:rm_case1}
	\end{equation}
	Wir zeigen: (\ref{eq:non_rm_feasible}) $\implies$ (\ref{eq:rm_case1}). Aus
	(\ref{eq:non_rm_feasible}) folgt durch Multiplikation mit $F$:
	$F C_1 + F C_2 \leq F T_1$, und da $F \geq 1$:
	$F C_1 + C_2 \leq F T_1$.
	Addition von $C_1$ ergibt $(F+1)C_1 + C_2 \leq F T_1 + C_1$.
	Da wir in Fall 1 voraussetzen, dass $C_1 < T_2 - F T_1$, erhalten wir:
	$(F+1)C_1 + C_2 < F T_1 + (T_2 - F T_1) = T_2$,
	womit (\ref{eq:rm_case1}) erfüllt ist.

	\textbf{Fall 2:} $C_1 \geq T_2 - F T_1$.

	Hier überlappt die Ausführung von $\tau_1$ mit der zweiten Instanz von $\tau_2$.
	Die Machbarkeitsbedingung für $\tau_2$ unter RM ist:
	\begin{equation}
		F C_1 + C_2 \leq F T_1.
		\label{eq:rm_case2}
	\end{equation}
	Aus (\ref{eq:non_rm_feasible}) folgt durch Multiplikation mit $F$:
	$F C_1 + F C_2 \leq F T_1$, und da $F \geq 1$: $C_2 \leq F C_2$, also:
	$F C_1 + C_2 \leq F C_1 + F C_2 \leq F T_1$,
	womit (\ref{eq:rm_case2}) erfüllt ist.

	In beiden Fällen gilt: Wenn das Task-Set unter der nicht-RM-Zuweisung
	schedulierbar ist, ist es auch unter RM schedulierbar.

	\textbf{Schritt 2: Verallgemeinerung auf $n$ Tasks.}

	Sei $\Gamma = \{\tau_1, \ldots, \tau_n\}$ ein Task-Set, das von einer
	beliebigen festen Prioritätszuweisung $\pi$ scheduliert wird. Wenn $\pi$
	nicht die RM-Reihenfolge ist, existieren zwei Tasks $\tau_i, \tau_j$ mit
	$T_i < T_j$, aber $P_i^{\pi} < P_j^{\pi}$ (invertierte Priorität). Durch
	Anwendung des Zwei-Task-Argumentes aus Schritt 1 können die Prioritäten
	von $\tau_i$ und $\tau_j$ vertauscht werden, ohne die Schedulierbarkeit zu
	verletzen. Da das Task-Set endlich ist, führt die iterierte Anwendung dieser
	Transposition nach endlich vielen Schritten zur RM-Reihenfolge, ohne dass
	dabei eine Deadline verletzt wird. \hfill$\blacksquare$
\end{proof}

\subsection{Berechnung des Least Upper Bound für zwei Tasks}

Der Least Upper Bound $U_{\text{lub}}$ ist die kleinste Schranke der Prozessorauslastung,
unterhalb derer ein Task-Set unter RM \emph{garantiert} schedulierbar ist.

\subsubsection{Vorgehen}

Für $n = 2$ Tasks $\tau_1$ und $\tau_2$ mit $T_1 < T_2$ berechnen wir $U_{\text{lub}}$
wie folgt:
\begin{enumerate}
	\item Zuweisung gemäß RM: $\tau_1$ erhält die höhere Priorität.
	\item Erhöhung von $C_2$ bis der Prozessor vollständig ausgelastet ist.
	\item Minimierung der oberen Schranke $U_{\text{ub}}$ über alle Periodenrelationen.
\end{enumerate}

Sei wiederum $F = \lfloor T_2 / T_1 \rfloor$.

\subsubsection{Fall 1: $C_1 \leq T_2 - F T_1$}

Der größtmögliche Wert von $C_2$ beträgt $C_2 = T_2 - C_1(F+1)$, und die
obere Schranke der Auslastung ist:
\[
	U_{\text{ub}} = \frac{C_1}{T_1} + \frac{C_2}{T_2}
	= \frac{C_1}{T_1} + \frac{T_2 - C_1(F+1)}{T_2}
	= 1 + \frac{C_1}{T_2}\left(\frac{T_2}{T_1} - (F+1)\right).
\]
Da $T_2/T_1 \geq F + 1$ in diesem Fall nicht gilt (der Klammerausdruck ist negativ),
ist $U_{\text{ub}}$ monoton fallend in $C_1$. Das Minimum wird also bei dem größten
erlaubten $C_1$ erreicht, nämlich $C_1 = T_2 - F T_1$.

\subsubsection{Fall 2: $C_1 \geq T_2 - F T_1$}

Der größtmögliche Wert von $C_2$ beträgt $C_2 = (T_1 - C_1) F$, und:
\[
	U_{\text{ub}} = \frac{C_1}{T_1} + \frac{(T_1 - C_1)F}{T_2}
	= \frac{T_1}{T_2} F + \frac{C_1}{T_2}\left(\frac{T_2}{T_1} - F\right).
\]
Hier ist $U_{\text{ub}}$ monoton steigend in $C_1$. Das Minimum liegt also bei
$C_1 = T_2 - F T_1$.

\subsubsection{Gemeinsamer Minimierungspunkt}

In beiden Fällen wird das Minimum bei $C_1 = T_2 - T_1 F$ erreicht. Einsetzen liefert
mit der Abkürzung $G = T_2/T_1 - F$ (sodass $0 \leq G < 1$):
\[
	U_{\text{ub}} = \frac{T_1}{T_2}F + \frac{(T_2 - T_1 F)}{T_2}
	\cdot\frac{T_2/T_1 - F}{1}
	= \frac{F + G^2}{F + G}.
\]
Für $F = 1$ (den ungünstigsten Fall) vereinfacht sich das zu:
\[
	U = \frac{1 + G^2}{1 + G}.
\]
Minimierung über $G \in [0,1)$ durch Ableiten ergibt:
\[
	\frac{dU}{dG} = \frac{2G(1+G) - (1+G^2)}{(1+G)^2} = \frac{G^2 + 2G - 1}{(1+G)^2} = 0,
\]
also $G = -1 + \sqrt{2}$ (positive Lösung). Einsetzen liefert:
\[
	U_{\text{lub}}(2) = \frac{1 + (\sqrt{2}-1)^2}{1 + (\sqrt{2}-1)}
	= \frac{4 - 2\sqrt{2}}{2} = 2(\sqrt{2} - 1) \approx 0{,}828.
\]

\begin{bemerkung}
	Wenn $T_2$ ein ganzzahliges Vielfaches von $T_1$ ist, gilt $G = 0$ und
	$U_{\text{ub}} = 1$. Das bedeutet: Für harmonische Perioden ($T_2 = k T_1$)
	kann RM den Prozessor vollständig auslasten.
\end{bemerkung}

\subsection{Berechnung des Least Upper Bound für $n$ Tasks}

\begin{theorem}[Liu \& Layland 1973 -- RM Least Upper Bound]
	\label{satz:rm_lub}
	Für eine beliebige Menge von $n$ periodischen Tasks unter den Annahmen A1--A4
	ist der Least Upper Bound der Prozessorauslastung unter dem Rate Monotonic
	Algorithmus:
	\[
		U_{\text{lub}}(\text{RM}, n) = n\left(2^{1/n} - 1\right).
	\]
	Für $n \to \infty$ gilt $U_{\text{lub}} \to \ln 2 \approx 0{,}693$.
\end{theorem}

\begin{proof}
	Die schlechteste Bedingung für die Schedulierbarkeit, d.\,h.\ das globale
	Minimum der oberen Schranke $U_{\text{ub}}$ über alle Task-Sets, die den
	Prozessor vollständig auslasten, wird durch folgende Parameterwahl erreicht
	\cite{Liu1973}:
	\begin{equation}
		\begin{cases}
			T_1 < T_n < 2 T_1 \\
			C_1 = T_2 - T_1 \\
			C_2 = T_3 - T_2 \\
			\quad\vdots \\
			C_{n-1} = T_n - T_{n-1} \\
			C_n = 2T_1 - T_n.
		\end{cases}
		\label{eq:worst_case_params}
	\end{equation}
	Die resultierende Prozessorauslastung lässt sich mit den Abkürzungen
	$R_i = T_{i+1}/T_i$ schreiben als:
	\[
		U = \sum_{i=1}^{n-1} R_i + \frac{2}{R_1 R_2 \cdots R_{n-1}} - n.
	\]
	Minimierung über alle $R_i$ durch partielles Ableiten:
	\[
		\frac{\partial U}{\partial R_k} = 1 - \frac{2}{R_k^2 \cdot \prod_{i \neq k} R_i} = 0.
	\]
	Mit $P = \prod_{i=1}^{n-1} R_i$ erhält man $R_k P = 2$ für alle $k$, also:
	\[
		R_1 = R_2 = \cdots = R_{n-1} = 2^{1/n}.
	\]
	Einsetzen dieses gemeinsamen Wertes ergibt:
	\begin{align*}
		U_{\text{lub}}
		&= (n-1) \cdot 2^{1/n} + \frac{2}{2^{(n-1)/n}} - n \\
		&= (n-1) \cdot 2^{1/n} + 2^{1/n} - n \\
		&= n \cdot 2^{1/n} - n \\
		&= n(2^{1/n} - 1).
	\end{align*}

	\textbf{Grenzwert für $n \to \infty$:}

	Mit der Substitution $y = 2^{1/n} - 1$, also $n = \ln 2 / \ln(y+1)$, gilt:
	\[
		\lim_{n \to \infty} n(2^{1/n} - 1) = \ln 2 \cdot \lim_{y \to 0} \frac{y}{\ln(y+1)}.
	\]
	Nach der Regel von L'Hôpital:
	\[
		\lim_{y \to 0} \frac{y}{\ln(y+1)} = \lim_{y \to 0} \frac{1}{1/(y+1)} = 1.
	\]
	Daher gilt $\lim_{n\to\infty} U_{\text{lub}} = \ln 2 \approx 0{,}693$. \hfill$\blacksquare$
\end{proof}

\begin{table}[H]
	\centering
	\caption{Werte von $U_{\text{lub}}$ für RM als Funktion der Taskanzahl $n$}
	\label{tab:rm_lub}
	\begin{tabular}{cc|cc}
		\toprule
		$n$ & $U_{\text{lub}}$ & $n$ & $U_{\text{lub}}$ \\
		\midrule
		1   & 1{,}000          & 6   & 0{,}735          \\
		2   & 0{,}828          & 7   & 0{,}729          \\
		3   & 0{,}780          & 8   & 0{,}724          \\
		4   & 0{,}757          & 9   & 0{,}721          \\
		5   & 0{,}743          & 10  & 0{,}718          \\
		$\infty$ & $\ln 2 \approx 0{,}693$ & & \\
		\bottomrule
	\end{tabular}
\end{table}

\begin{bemerkung}
	Die Tabelle zeigt, dass der Least Upper Bound mit wachsender Taskanzahl
	monoton gegen $\ln 2 \approx 0{,}693$ fällt. Das bedeutet: Für sehr große
	Task-Sets kann RM nur Task-Mengen mit Auslastung unter 69\,\% garantiert
	einplanen. EDF hingegen garantiert Schedulierbarkeit bis $U = 1$, also eine
	um bis zu 31 Prozentpunkte höhere Nutzung der Prozessorkapazität.
\end{bemerkung}

\subsection{Praktische Anwendung: Hinreichende Schedulierbarkeitstest}

Der Liu-Layland-Bound liefert einen einfachen, hinreichenden Test:

\begin{satz}[RM Hinreichende Bedingung -- Liu \& Layland]
	\label{satz:rm_suf}
	Eine Menge von $n$ periodischen Tasks ist unter RM schedulierbar, wenn:
	\[
		U = \sum_{i=1}^{n} \frac{C_i}{T_i} \leq n(2^{1/n} - 1).
	\]
\end{satz}

\begin{beispiel}[Anwendung des Liu-Layland-Tests]
	Gegeben:
	\[
		\tau_1: C_1=1, T_1=4 \quad \tau_2: C_2=2, T_2=6 \quad \tau_3: C_3=3, T_3=12.
	\]
	Auslastung: $U = 1/4 + 2/6 + 3/12 = 0{,}25 + 0{,}333 + 0{,}25 = 0{,}833$.
	Schranke: $U_{\text{lub}}(3) = 3(2^{1/3}-1) \approx 0{,}780$.
	Da $U = 0{,}833 > 0{,}780$, gibt der Test \emph{keine Garantie} -- aber das
	Task-Set könnte dennoch schedulierbar sein.
\end{beispiel}

\begin{bemerkung}[Test ist hinreichend, nicht notwendig]
	Der Liu-Layland-Test ist \textbf{hinreichend}, aber \textbf{nicht notwendig}.
	Ein Task-Set mit $U > U_{\text{lub}}$ kann trotzdem mit RM schedulierbar sein,
	wenn die Perioden günstig zueinander liegen (z.\,B.\ bei harmonischen Perioden).
	Umgekehrt: Wenn $U \leq U_{\text{lub}}$, ist die Schedulierbarkeit garantiert.
	Zum exakten Test dient die Response Time Analysis (Abschnitt~\ref{sec:rta}).
\end{bemerkung}

\subsection{Hyperbolic Bound -- Engere Schranke für RM}

Der Hyperbolic Bound von Bini und Buttazzo \cite{Bini2001} ist eine schärfere
hinreichende Bedingung als der Liu-Layland-Bound. Er nutzt die individuellen
Auslastungsfaktoren $U_i = C_i/T_i$ statt ihrer Summe.

\begin{theorem}[Bini, Buttazzo \& Buttazzo 2001 -- Hyperbolic Bound]
	\label{thm:hyperbolic}
	Sei $\Gamma = \{\tau_1, \ldots, \tau_n\}$ mit $U_i = C_i/T_i$. Dann ist
	$\Gamma$ unter RM schedulierbar, wenn:
	\[
		\prod_{i=1}^{n} (U_i + 1) \leq 2.
	\]
\end{theorem}

\begin{proof}
	Der Beweis nutzt die gleiche Worst-Case-Parametrisierung (\ref{eq:worst_case_params}).
	Dort gilt $U_i = R_i - 1$ für $i = 1,\ldots,n-1$, also $R_i = U_i + 1$. Damit ist
	$\prod_{i=1}^{n-1} R_i = T_n/T_1$. Die Schedulierbarkeitsbed.~(\ref{eq:rm_case2})
	für $\tau_n$ lautet $U_n \leq 2T_1/T_n - 1$, d.\,h.\ $(U_n+1) \leq 2 / \prod_{i=1}^{n-1}(U_i+1)$,
	woraus $\prod_{i=1}^{n}(U_i+1) \leq 2$ folgt. \hfill$\blacksquare$
\end{proof}

\begin{bemerkung}[Vergleich der Schranken]
	Der Hyperbolic Bound ist stets mindestens so gut wie der Liu-Layland-Bound und
	häufig besser. Er akzeptiert Task-Sets, die der Liu-Layland-Bound ablehnen würde.
	Formell gilt:
	\[
		\prod_{i=1}^{n}(U_i + 1) \leq 2 \text{ ist eine striktere Bedingung als } \sum_{i=1}^n U_i \leq n(2^{1/n}-1)
	\]
	in dem Sinne, dass die zulässige Schedulierbarkeitsregion durch den Hyperbolic
	Bound größer ist. Der Gewinn wächst mit der Anzahl der Tasks und nähert sich
	$\sqrt{2}$ für $n \to \infty$.
\end{bemerkung}

\begin{beispiel}[Hyperbolic Bound vs. Liu-Layland]
	Gegeben: $\tau_1: U_1 = 0{,}4$, $\tau_2: U_2 = 0{,}4$.

	\textit{Liu-Layland:} $U = 0{,}8 \leq 2(\sqrt{2}-1) \approx 0{,}828$. Bestanden.

	\textit{Hyperbolic Bound:} $(1{,}4)(1{,}4) = 1{,}96 \leq 2$. Bestanden.

	Nun: $\tau_1: U_1 = 0{,}41$, $\tau_2: U_2 = 0{,}41$.

	\textit{Liu-Layland:} $U = 0{,}82 \leq 0{,}828$. Bestanden.

	\textit{Hyperbolic Bound:} $(1{,}41)^2 = 1{,}9881 \leq 2$. Bestanden.

	Nun: $\tau_1: U_1 = 0{,}43$, $\tau_2: U_2 = 0{,}43$.

	\textit{Liu-Layland:} $U = 0{,}86 > 0{,}828$. \emph{Kein Ergebnis!}

	\textit{Hyperbolic Bound:} $(1{,}43)^2 = 2{,}0449 > 2$. \emph{Kein Ergebnis!}

	Aber mit $U_1 = 0{,}42$, $U_2 = 0{,}40$:

	\textit{Liu-Layland:} $U = 0{,}82 \leq 0{,}828$. Bestanden.

	\textit{Hyperbolic Bound:} $(1{,}42)(1{,}40) = 1{,}988 \leq 2$. Bestanden.

	Dagegen mit $U_1 = 0{,}44$, $U_2 = 0{,}40$:

	\textit{Liu-Layland:} $U = 0{,}84 > 0{,}828$. \emph{Keine Garantie.}

	\textit{Hyperbolic Bound:} $(1{,}44)(1{,}40) = 2{,}016 > 2$. \emph{Keine Garantie.}
\end{beispiel}

\subsection{Response Time Analysis -- Exakter Schedulierbarkeitstest}
\label{sec:rta}

Der Liu-Layland-Bound und der Hyperbolic Bound sind hinreichend, aber nicht
notwendig. Für einen \emph{exakten} Test benötigt man die Response Time Analysis
(RTA), die von Audsley et al. \cite{Audsley1993} entwickelt wurde.

\subsubsection{Grundprinzip}

Die worst-case Antwortzeit $R_i$ einer Task $\tau_i$ setzt sich zusammen aus ihrer
eigenen Ausführungszeit $C_i$ und der Interferenz $I_i$ durch alle Tasks mit höherer
Priorität. Die Interferenz im Intervall $[0, R_i]$ ist die Summe aller Ausführungen
hochpriorer Tasks, die in dieses Intervall fallen:
\[
	I_i = \sum_{j \in hp(i)} \left\lceil \frac{R_i}{T_j} \right\rceil C_j,
\]
wobei $hp(i) = \{\tau_j \mid P_j > P_i\}$ die Menge der Tasks mit höherer Priorität
ist und $\lceil R_i/T_j \rceil$ die Anzahl der Instanzen von $\tau_j$ im Intervall
$[0, R_i]$ angibt.

\begin{definition}[Worst-Case Antwortzeit unter RM]
	Die \textbf{worst-case Antwortzeit} $R_i$ von $\tau_i$ ist die kleinste positive
	Lösung der impliziten Gleichung:
	\begin{equation}
		R_i = C_i + \sum_{j \in hp(i)} \left\lceil \frac{R_i}{T_j} \right\rceil C_j.
		\label{eq:rta}
	\end{equation}
\end{definition}

Da $R_i$ auf beiden Seiten von Gleichung~(\ref{eq:rta}) auftritt, gibt es keine
geschlossene Lösung. Die Berechnung erfolgt iterativ.

\subsubsection{Iterativer Berechnungsalgorithmus}

\begin{algorithmic}[1]
	\Require Task $\tau_i$, geordnet nach RM-Priorität
	\State $R_i^{(0)} \leftarrow \sum_{j=1}^{i} C_j$ \Comment{Startschätzung: frühester Fertigstellungszeitpunkt}
	\Repeat
	\State $R_i^{(k+1)} \leftarrow C_i + \sum_{j \in hp(i)} \left\lceil \dfrac{R_i^{(k)}}{T_j} \right\rceil C_j$
	\Until{$R_i^{(k+1)} = R_i^{(k)}$ \textbf{oder} $R_i^{(k+1)} > T_i$}
	\If{$R_i^{(k+1)} \leq T_i$}
	\State \Return SCHEDULIERBAR
	\Else
	\State \Return NICHT SCHEDULIERBAR
	\EndIf
\end{algorithmic}

Die Sequenz $R_i^{(0)}, R_i^{(1)}, R_i^{(2)}, \ldots$ ist monoton nichtfallend und
durch $T_i$ nach oben beschränkt. Sie konvergiert daher entweder zu einem Fixpunkt
(der Task schedulierbar) oder überschreitet $T_i$ (Task nicht schedulierbar).

\begin{beispiel}[Response Time Analysis -- Schritt für Schritt]
	Gegeben:
	\begin{center}
		\begin{tabular}{cccc}
			\toprule
			Task    & $C_i$ & $T_i$ & Priorität (RM) \\
			\midrule
			$\tau_1$ & 1     & 4     & 1 (höchste) \\
			$\tau_2$ & 2     & 6     & 2 \\
			$\tau_3$ & 3     & 10    & 3 (niedrigste) \\
			\bottomrule
		\end{tabular}
	\end{center}

	Auslastung: $U = 1/4 + 2/6 + 3/10 = 0{,}25 + 0{,}333 + 0{,}30 = 0{,}883$.
	Da $U_{\text{lub}}(3) \approx 0{,}780 < 0{,}883$, liefert der Liu-Layland-Test
	kein Ergebnis. Wir nutzen die RTA.

	\textbf{Task $\tau_1$} (keine höherpriore Task):
	$R_1 = C_1 = 1 \leq T_1 = 4$. \checkmark

	\textbf{Task $\tau_2$} ($hp(2) = \{\tau_1\}$):
	\begin{align*}
		R_2^{(0)} &= C_1 + C_2 = 3 \\
		R_2^{(1)} &= C_2 + \lceil 3/4 \rceil \cdot C_1 = 2 + 1 \cdot 1 = 3 = R_2^{(0)}.
	\end{align*}
	Konvergenz bei $R_2 = 3 \leq T_2 = 6$. \checkmark

	\textbf{Task $\tau_3$} ($hp(3) = \{\tau_1, \tau_2\}$):
	\begin{align*}
		R_3^{(0)} &= C_1 + C_2 + C_3 = 6 \\
		R_3^{(1)} &= 3 + \lceil 6/4 \rceil \cdot 1 + \lceil 6/6 \rceil \cdot 2 = 3 + 2 + 2 = 7 \\
		R_3^{(2)} &= 3 + \lceil 7/4 \rceil \cdot 1 + \lceil 7/6 \rceil \cdot 2 = 3 + 2 + 4 = 9 \\
		R_3^{(3)} &= 3 + \lceil 9/4 \rceil \cdot 1 + \lceil 9/6 \rceil \cdot 2 = 3 + 3 + 4 = 10 \\
		R_3^{(4)} &= 3 + \lceil 10/4 \rceil \cdot 1 + \lceil 10/6 \rceil \cdot 2 = 3 + 3 + 4 = 10 = R_3^{(3)}.
	\end{align*}
	Konvergenz bei $R_3 = 10 = T_3$. \checkmark (Gerade noch schedulierbar.)

	Das Task-Set ist unter RM schedulierbar, obwohl $U > U_{\text{lub}}$.
	Dies zeigt, dass der Liu-Layland-Bound zu pessimistisch war.
\end{beispiel}

\subsection{Schedulierbarkeit bei kritischem Augenblick -- Direkte Methode}

Alternativ zur Response Time Analysis kann die Schedulierbarkeit direkt durch
Konstruktion des Schedules am kritischen Augenblick überprüft werden.

\begin{satz}[Schedulierbarkeit am kritischen Augenblick]
	Ein Task-Set $\Gamma$ ist unter RM genau dann schedulierbar, wenn jede Task
	$\tau_i$ ihre Deadline an ihrem kritischen Augenblick einhält, d.\,h.\ wenn alle
	Tasks zum Zeitpunkt $t = 0$ gleichzeitig aktiviert werden und jede Task innerhalb
	ihrer ersten Periode fertiggestellt wird.
\end{satz}

\begin{beispiel}[Schedulierbarkeit direkt am kritischen Augenblick]
	Gegeben: $\tau_1: C_1=1, T_1=4$; $\tau_2: C_2=2, T_2=6$.

	Auslastung: $U = 1/4 + 2/6 = 0{,}583 \leq U_{\text{lub}}(2) \approx 0{,}828$. Garantiert.

	Konstruktion des Schedules ab $t=0$:
	\begin{itemize}
		\item $[0,1)$: $\tau_1$ läuft (höhere Priorität, $T_1 < T_2$).
		\item $[1,3)$: $\tau_2$ läuft (keine andere Task bereit).
		\item $[3,4)$: $\tau_2$ läuft weiter -- aber $\tau_1$ hat neue Instanz bei $t=4$!
		\item Bei $t=3$: $\tau_2$ ist fertig ($f_{2,1}=3 \leq d_{2,1}=6$). \checkmark
		\item $[3,4)$: Leerlauf.
		\item $[4,5)$: $\tau_1$, zweite Instanz.
		\item $[5,7)$: $\tau_2$, zweite Instanz. $f_{2,2} = 7 \leq d_{2,2} = 12$. \checkmark
	\end{itemize}
	Das Task-Set ist schedulierbar.
\end{beispiel}

\subsection{Workload-Analyse nach Lehoczky, Sha und Ding}

Lehoczky, Sha und Ding \cite{Lehoczky1989} entwickelten eine präzisere Methode
zur Schedulierbarkeitsanalyse unter fester Priorität, die auf dem Konzept der
Level-$i$-Arbeitslast basiert.

\begin{definition}[Level-$i$-Arbeitslast]
	Die \textbf{Level-$i$-Arbeitslast} $W_i(t)$ ist die kumulative Rechenzeit,
	die im Intervall $(0, t]$ von Task $\tau_i$ und allen Tasks mit höherer Priorität
	angefordert wird:
	\[
		W_i(t) = C_i + \sum_{h: P_h > P_i} \left\lfloor \frac{t}{T_h} \right\rfloor C_h.
	\]
\end{definition}

\begin{theorem}[Lehoczky, Sha, Ding 1989]
	Eine Menge vollständig preemptiver periodischer Tasks ist unter einem
	Algorithmus mit fester Priorität genau dann schedulierbar, wenn:
	\[
		\forall i = 1,\ldots,n: \quad \exists t \in (0, T_i]: \quad W_i(t) \leq t.
	\]
\end{theorem}

\subsection{Warum RM schwach optimal ist -- Präzisierung}

\begin{definition}[Schwach optimal -- Präzisierung]
	\label{def:schwach_optimal_praez}
	RM ist \textbf{schwach optimal} in folgendem Sinne:
	\begin{enumerate}
		\item RM ist optimal \emph{innerhalb} der Klasse aller Algorithmen mit
		      fester Prioritätszuweisung.
		\item Es gibt Task-Mengen mit $\ln 2 < U \leq 1$, die \emph{nicht} unter
		      RM schedulierbar sind, aber unter EDF schedulierbar sind.
		\item RM ist \emph{nicht} optimal über alle Scheduling-Algorithmen.
	\end{enumerate}
\end{definition}

\begin{satz}[Grenzen der Schwäche von RM]
	Für jede Auslastung $U$ mit $U_{\text{lub}}(\text{RM},n) < U \leq 1$ existiert
	ein Task-Set $\Gamma_U$ mit $n$ Tasks und Auslastung $U$, das:
	\begin{itemize}
		\item von EDF schedulierbar ist (da $U \leq 1$),
		\item von RM \emph{nicht} schedulierbar ist.
	\end{itemize}
\end{satz}

\begin{proof}
	Wähle $\Gamma_U$ gemäß den Worst-Case-Parametern (\ref{eq:worst_case_params})
	mit der gegebenen Auslastung $U$. Per Konstruktion ist $\Gamma_U$ nicht unter RM
	schedulierbar (da $U > U_{\text{lub}}$). Da $U \leq 1$, ist $\Gamma_U$ nach
	Satz~\ref{thm:edf_schedulierbar} unter EDF schedulierbar. \hfill$\blacksquare$
\end{proof}

\subsection{Vollständiges Beispiel: Vergleich RM und EDF bei hoher Auslastung}

\begin{beispiel}[RM versagt, EDF scheduliert]
	\label{bsp:rm_edf_vergleich}
	Gegeben:
	\[
		\tau_1: C_1 = 3, \quad T_1 = 5 \qquad \tau_2: C_2 = 4, \quad T_2 = 7.
	\]
	Auslastung: $U = 3/5 + 4/7 = 0{,}6 + 0{,}571 = 34/35 \approx 0{,}971$.

	\textbf{RM-Analyse:}
	$U_{\text{lub}}(2) \approx 0{,}828 < 0{,}971$. Der Liu-Layland-Test gibt keine
	Garantie. Exakter Test mittels RTA:
	\begin{align*}
		R_1 &= C_1 = 3 \leq 5. \quad\checkmark \\
		R_2^{(0)} &= 7 \\
		R_2^{(1)} &= 4 + \lceil 7/5 \rceil \cdot 3 = 4 + 2 \cdot 3 = 10 > T_2 = 7.
	\end{align*}
	$\tau_2$ ist unter RM \emph{nicht schedulierbar}. Am kritischen Augenblick
	(alle Tasks starten bei $t=0$) verpasst $\tau_2$ ihre erste Deadline bei $t=7$.

	\textbf{EDF-Analyse:}
	$U \approx 0{,}971 \leq 1$. Task-Set ist unter EDF schedulierbar. Der Schedule
	verläuft ohne Deadline-Verletzungen, da EDF zu jedem Zeitpunkt die Task mit der
	frühesten absoluten Deadline ausführt und dabei die volle Prozessorkapazität nutzt.

	Dieser direkte Vergleich illustriert die Aussage:
	\begin{quote}
		\textbf{Wenn es einen machbaren Schedule gibt, findet RM diesen sofort. RM ist schwach optimal.
		Wenn es einen machbaren Schedule gibt, findet EDF diesen sofort. EDF ist scharf optimal.}
	\end{quote}
	Für $\tau_2$ existiert ein machbarer Schedule (unter EDF), aber RM findet ihn nicht.
\end{beispiel}

\subsection{Übungsaufgaben zu Rate Monotonic}

\begin{beispiel}[Übungsaufgabe 1 -- Schedulierbarkeitstest]
	Bestimme, ob folgendes Task-Set unter RM schedulierbar ist:
	\[
		\tau_1: C=2, T=6 \quad \tau_2: C=2, T=8 \quad \tau_3: C=2, T=12.
	\]
	\textit{Lösung:} $U = 2/6 + 2/8 + 2/12 = 1/3 + 1/4 + 1/6 = 3/4 = 0{,}75$.
	$U_{\text{lub}}(3) = 3(2^{1/3}-1) \approx 0{,}780 > 0{,}75$. Schedulierbar. \checkmark
\end{beispiel}

\begin{beispiel}[Übungsaufgabe 2 -- Nicht schedulierbar]
	Bestimme, ob folgendes Task-Set unter RM schedulierbar ist:
	\[
		\tau_1: C=1, T=4 \quad \tau_2: C=2, T=6 \quad \tau_3: C=3, T=8.
	\]
	\textit{Lösung:} $U = 1/4 + 2/6 + 3/8 = 0{,}25 + 0{,}333 + 0{,}375 = 0{,}958$.
	$U_{\text{lub}}(3) \approx 0{,}780 < 0{,}958$. Kein Ergebnis aus Liu-Layland-Test.
	RTA für $\tau_3$:
	\begin{align*}
		R_3^{(0)} &= 1 + 2 + 3 = 6 \\
		R_3^{(1)} &= 3 + \lceil 6/4\rceil\cdot1 + \lceil 6/6\rceil\cdot2 = 3+2+2 = 7 \\
		R_3^{(2)} &= 3 + \lceil 7/4\rceil\cdot1 + \lceil 7/6\rceil\cdot2 = 3+2+4 = 9 \\
		R_3^{(3)} &= 3 + \lceil 9/4\rceil\cdot1 + \lceil 9/6\rceil\cdot2 = 3+3+4 = 10 > T_3 = 8.
	\end{align*}
	$\tau_3$ ist unter RM nicht schedulierbar. Da $U \approx 0{,}958 \leq 1$,
	ist das Task-Set unter EDF schedulierbar.
\end{beispiel}

\section{Earliest Deadline First (EDF) -- Scharf optimal}

\subsection{Algorithmus}

EDF ist eine dynamische Scheduling-Regel: Tasks mit früheren absoluten Deadlines
werden zuerst ausgeführt. EDF ist typischerweise preemptiv.

\begin{quote}
\textbf{EDF ist scharf optimal: Wenn ein machbarer Schedule existiert, findet EDF diesen sofort. EDF findet die kürzeste Ausführungszeit aller Tasks sofort.}
\end{quote}

\subsection{Schedulierbarkeitsanalyse für EDF}

\begin{theorem}[Liu \& Layland 1973 -- EDF Schedulierbarkeit]
	\label{thm:edf_schedulierbar}
	Eine Menge periodischer Tasks ist mit EDF schedulierbar genau dann, wenn:
	\[
		U = \sum_{i=1}^{n} \frac{C_i}{T_i} \leq 1.
	\]
\end{theorem}

\begin{proof}
	\textbf{``Nur wenn'' Teil:} Wenn $U > 1$, überschreitet der Gesamtrechenaufwand
	in der Hyperperiode $H$ die verfügbare Zeit $H$. Kein Algorithmus kann diesen
	Schedule machbar machen.

	\textbf{``Wenn'' Teil (durch Widerspruch):} Angenommen, $U \leq 1$, aber der
	Task-Satz ist nicht schedulierbar. Sei $t_2$ der erste Zeitpunkt, an dem eine
	Deadline versäumt wird, und $[t_1, t_2]$ das längste Intervall kontinuierlicher
	Auslastung vor $t_2$, in dem nur Instanzen mit Deadline $\leq t_2$ ausgeführt werden.

	Der gesamte Rechenzeitbedarf in $[t_1, t_2]$ ist:
	\[
		C_p(t_1, t_2) = \sum_{i=1}^{n} \left\lfloor \frac{t_2 - t_1}{T_i} \right\rfloor C_i \leq (t_2 - t_1) U.
	\]

	Da eine Deadline versäumt wird, muss $C_p(t_1, t_2) > (t_2 - t_1)$ gelten,
	also $U > 1$ -- Widerspruch. \hfill$\blacksquare$
\end{proof}

\subsection{Die Schärfe von EDF im Vergleich zu RM}

Der fundamentale Unterschied:

\begin{satz}[Schärfe von EDF gegenüber RM]
	\label{satz:edf_scharf}
	Es gibt Task-Mengen, die EDF schedulieren kann, RM aber nicht. Insbesondere gilt
	für alle $U$ mit $\ln 2 < U \leq 1$:
	\begin{itemize}
		\item EDF garantiert die Machbarkeit (nach Satz~\ref{thm:edf_schedulierbar}).
		\item RM kann die Machbarkeit \emph{nicht} garantieren (nach Satz~\ref{satz:rm_lub}).
	\end{itemize}
\end{satz}

\begin{beispiel}
	Betrachte eine periodische Task-Menge mit zwei Tasks:
	\[
		\tau_1: C_1 = 3, \quad T_1 = 5 \qquad \tau_2: C_2 = 4, \quad T_2 = 7.
	\]
	Die Prozessorauslastung ist:
	\[
		U = \frac{3}{5} + \frac{4}{7} = 0{,}6 + 0{,}571 = \frac{34}{35} \approx 0{,}97.
	\]
	Da $U = 0{,}97 > 2(\sqrt{2}-1) \approx 0{,}83$, kann RM die Schedulierbarkeit
	\emph{nicht} garantieren -- und erzeugt tatsächlich eine Deadline-Verletzung.
	Da $U \approx 0{,}97 \leq 1$, ist der Task-Satz unter EDF schedulierbar.
	Dies zeigt die Schwäche von RM und die Schärfe von EDF.
\end{beispiel}

\section{Deadline Monotonic Scheduling}

Deadline Monotonic (DM) ist eine Erweiterung von RM für Tasks mit Deadlines
kleiner als ihre Perioden ($D_i \leq T_i$). DM weist Prioritäten umgekehrt
proportional zur relativen Deadline zu.

DM ist optimal unter allen Algorithmen mit fester Priorität für Tasks mit
$D_i \leq T_i$ -- analog zu RM ist DM damit ebenfalls \emph{schwach optimal}
im Sinne dieser Klasse.

\section{EDF mit eingeschränkten Deadlines}

Für Tasks mit $D_i \leq T_i$ kann die Schedulierbarkeit unter EDF durch das
Processor Demand Criterion (Baruah et al. 1990) getestet werden:

\begin{satz}[Schedulierbarkeit unter EDF mit eingeschränkten Deadlines]
	Eine synchrone Menge periodischer Tasks mit $D_i \leq T_i$ ist mit EDF
	schedulierbar genau dann, wenn $U \leq 1$ und
	\[
		\forall t > 0: \quad \mathrm{dbf}(t) = \sum_{i=1}^{n} \left\lfloor \frac{t + T_i - D_i}{T_i} \right\rfloor C_i \leq t.
	\]
\end{satz}

\section{Vergleich RM und EDF}

Tabelle~\ref{tab:rm_edf} fasst die wesentlichen Unterschiede zusammen:

\begin{table}[H]
	\centering
	\caption{Vergleich RM und EDF}
	\label{tab:rm_edf}
	\begin{tabular}{p{4cm}p{5cm}p{5cm}}
		\toprule
		\textbf{Eigenschaft} & \textbf{RM} & \textbf{EDF} \\
		\midrule
		Priorität            & Fest (umgekehrt proportional zu $T_i$) & Dynamisch (absolutes Deadline) \\
		Optimalität          & \textbf{Schwach optimal} (innerhalb fester Priorität) & \textbf{Scharf optimal} (insgesamt) \\
		$U_{\text{lub}}$     & $\ln 2 \approx 0{,}693$ (worst case) & $1{,}000$ \\
		Schedulierbarkeitstest & Pseudo-polynomiell (RTA) & $O(n)$ (Auslastungstest) \\
		Überlastverhalten    & Niederpriore Tasks blockieren & Periodenskalierung (Cervin) \\
		Impl.-Aufwand & $O(1)$ (Multi-Level Queue) & $O(\log n)$ (Heap) \\
		Präemptionsanzahl    & Höher & Niedriger \\
		\bottomrule
	\end{tabular}
\end{table}

\begin{bemerkung}
	Dass RM schwach optimal ist, bedeutet: Wenn ein Task-Satz überhaupt
	mit \emph{irgendeiner} festen Prioritätszuweisung schedulierbar ist, so
	ist er auch mit RM schedulierbar. Das schließt jedoch nicht aus, dass
	ein dynamischer Algorithmus wie EDF mehr Task-Mengen schedulieren kann.

	EDF ist scharf optimal: EDF scheduliert genau die Task-Mengen, die
	überhaupt schedulierbar sind ($U \leq 1$). Es gibt keine Task-Menge
	mit $U \leq 1$, die EDF nicht schedulieren kann.
\end{bemerkung}

%=============================================================
\chapter{Fixed-Priority Server -- Aperiodische Tasks unter RM}
\label{chap:server}
%=============================================================

\section{Einführung und Problemstellung}

In realen Echtzeit-Systemen treten stets \emph{hybride Task-Sets} auf: harte
periodische Tasks (Sensor-Abfragen, Regelschleifen) müssen zusammen mit weichen
aperiodischen Tasks (Benutzereingaben, Diagnoseereignisse) eingeplant werden.

\textbf{Das zentrale Ziel} dieses Kapitels: Für ein hybrides System sollen
\begin{enumerate}
	\item die Deadlines aller harten periodischen Tasks garantiert werden, und
	\item die Antwortzeiten der weichen aperiodischen Tasks minimiert werden --
	      ohne die Garantien der periodischen Tasks zu gefährden.
\end{enumerate}

Alle Algorithmen dieses Kapitels setzen voraus:
\begin{itemize}
	\item Periodische Tasks werden unter Rate Monotonic (RM) eingeplant.
	\item Alle periodischen Tasks starten bei $t=0$ mit $D_i = T_i$.
	\item Aperiodische Ankunftszeiten sind a priori unbekannt.
\end{itemize}

\textbf{Grundidee Server:} Ein \emph{Server} ist eine künstliche periodische
Task mit Periode $T_s$ und Kapazität (Budget) $C_s$, die aperiodische Anfragen
im Rahmen ihres Budgets bedient. Der Server konkurriert mit den echten periodischen
Tasks um den Prozessor.

\section{Background Scheduling -- Einfachster Ansatz}

\subsection{Informelle Beschreibung}

Aperiodische Tasks werden nur dann ausgeführt, wenn der Prozessor von keiner
periodischen Task benötigt wird. Dies entspricht einer impliziten Priorität
unterhalb aller periodischen Tasks.

\begin{algorithm}[H]
	\caption{Background Scheduling}
	\begin{algorithmic}[1]
		\State Verwalte periodische Tasks in Hochprioritätswarteschlange (RM)
		\State Verwalte aperiodische Tasks in Niedrigprioritätswarteschlange (FCFS)
		\Loop
		\If{periodische Task bereit}
		\State Führe höchstpriore periodische Task aus (RM)
		\ElsIf{aperiodische Task bereit}
		\State Führe nächste aperiodische Task aus (FCFS)
		\Else
		\State Leerlauf
		\EndIf
		\EndLoop
	\end{algorithmic}
\end{algorithm}

\subsection{Korrektheit und Optimalitätsklassifikation}

\begin{satz}[Background Scheduling -- Korrektheit]
	Unter Background Scheduling wird die Schedulierbarkeit aller periodischen
	Tasks durch RM nicht beeinflusst. Die aperiodischen Tasks erhalten genau
	die verbleibende Prozessorzeit $(1-U_p)$.
\end{satz}

\begin{proof}
	Aperiodische Tasks werden nur in Leerlaufzeiten ausgeführt. Da sie keine
	periodische Task preemptieren oder verzögern können, ist der Schedule für
	die periodischen Tasks identisch mit dem RM-Schedule ohne aperiodische Tasks.
	Die RM-Schedulierbarkeitsanalyse bleibt unverändert gültig. \hfill$\blacksquare$
\end{proof}

\textbf{Optimalitätsklassifikation:} Background Scheduling ist bezüglich
aperiodischer Antwortzeiten \emph{nicht optimal}: Bei hoher periodischer Last
$U_p \approx \ln 2$ bleibt nur $1 - \ln 2 \approx 0{,}307$ für aperiodische
Tasks übrig, und diese Bandbreite wird ineffizient genutzt (Tasks müssen auf
große Leerlaufintervalle warten).

\textbf{Laufzeit:} $O(1)$ pro Scheduling-Entscheidung (Vergleich zweier
Warteschlangen-Kopfelemente).

\section{Polling Server -- Expliziter Server}

\subsection{Informelle Beschreibung}

Der Polling Server (PS) ist eine periodische Task $(\tau_s, C_s, T_s)$, die
regelmäßig aktiviert wird und beim Aufwachen aperiodische Anfragen bedient
-- aber nur bis zur Kapazitätsgrenze $C_s$. Wenn keine aperiodischen Anfragen
vorliegen, wird das Budget \emph{verworfen}.

\begin{algorithm}[H]
	\caption{Polling Server}
	\begin{algorithmic}[1]
		\State Server wird periodisch alle $T_s$ Zeiteinheiten aktiviert
		\State Setze $c_s \leftarrow C_s$ \Comment{Budget auffüllen}
		\While{$c_s > 0$ \textbf{und} aperiodische Anfrage vorhanden}
		\State Bediene nächste aperiodische Anfrage für $\min(C_{\text{ape}}, c_s)$ Einheiten
		\State $c_s \leftarrow c_s - $ verbrauchte Zeit
		\EndWhile
		\State Verwerfe verbleibendes Budget $c_s$; wird erst zur nächsten Periode aufgefüllt
	\end{algorithmic}
\end{algorithm}

\subsection{Schedulierbarkeitsnachweis}

\begin{satz}[PS Schedulierbarkeit]
	\label{satz:ps}
	Eine Menge von $n$ periodischen Tasks mit Auslastung $U_p$ und ein Polling
	Server mit Auslastung $U_s = C_s/T_s$ sind unter RM gemeinsam schedulierbar,
	wenn:
	\[
		U_p + U_s \leq (n+1)\left(2^{1/(n+1)} - 1\right).
	\]
\end{satz}

\begin{proof}
	Der Polling Server verhält sich \emph{exakt} wie eine periodische Task mit
	Periode $T_s$ und WCET $C_s$: Unabhängig von der Anzahl der aperiodischen
	Anfragen verbraucht PS in jeder Serverperiode maximal $C_s$ Zeiteinheiten.
	Damit ist $U_s = C_s/T_s$ die effektive Auslastung des Servers.

	Das erweiterte System $\{\tau_1,\ldots,\tau_n, \tau_s\}$ hat $n+1$ Tasks
	mit Gesamtauslastung $U_p + U_s$. Nach dem Liu-Layland-Bound (Satz~\ref{satz:rm_lub})
	ist dieses System unter RM schedulierbar, wenn
	$U_p + U_s \leq (n+1)(2^{1/(n+1)}-1)$. \hfill$\blacksquare$
\end{proof}

\subsubsection{Engerer Bound: Hyperbolic Test für PS}

Mit dem Hyperbolic Bound (Theorem~\ref{thm:hyperbolic}):
\[
	\prod_{i=1}^{n}(U_i + 1) \leq \frac{2}{U_s + 1}.
\]

\subsection{Worst-Case-Antwortzeit aperiodischer Anfragen}

Im schlimmsten Fall muss eine aperiodische Anfrage bis zur nächsten Serveraktivierung
warten (wenn der Server gerade sein Budget aufgebraucht hat oder keine Anfrage
vorlag). Für eine Anfrage $J_a$ mit $C_a \leq C_s$:
\[
	R_a^{\max} \leq 2T_s - \varepsilon + C_a.
\]
Für beliebige $C_a$ mit exakt $\lceil C_a/C_s \rceil$ Serverzugriffen:
\[
	R_a \leq T_s + \left\lceil \frac{C_a}{C_s} \right\rceil T_s.
\]

\subsection{Optimalitätsklassifikation und Laufzeit}

\textbf{Optimalitätsklassifikation:} PS ist \emph{schwach optimal} innerhalb
der Klasse der periodischen Server unter RM: Bei vorgegebenem $(C_s, T_s)$
garantiert PS dieselbe periodische Schedulierbarkeit wie jede äquivalente
periodische Task. PS ist jedoch \emph{nicht scharf optimal} bezüglich aperiodischer
Antwortzeiten: Das Budget wird verworfen, wenn keine Anfragen vorliegen.

\textbf{Nachteil:} Eine aperiodische Anfrage, die kurz nach dem Verwerfen des
Budgets ankommt, muss fast eine volle Serverperiode warten.

\textbf{Laufzeit:} $O(\log n)$ für Warteschlangenoperationen (RM-Heap),
$O(1)$ für Budget-Verwaltung pro Aktivierung.

\section{Deferrable Server -- Kapazitätserhaltung}

\subsection{Informelle Beschreibung}

Der Deferrable Server (DS) verbessert den PS, indem das Budget bei Inaktivität
\emph{nicht verworfen} wird, sondern bis zum Ende der Serverperiode erhalten
bleibt. Kommt eine aperiodische Anfrage in der laufenden Periode an, steht
das Budget sofort zur Verfügung.

\begin{algorithm}[H]
	\caption{Deferrable Server}
	\begin{algorithmic}[1]
		\State $c_s \leftarrow 0$ \Comment{Initialisierung}
		\State
		\State \textbf{Zu jedem Serverperiodenbeginn:}
		\State $c_s \leftarrow C_s$ \Comment{Budget auf Maximum auffüllen}
		\State
		\State \textbf{Wenn aperiodische Anfrage $J_a$ ankommt und $c_s > 0$:}
		\State Bediene $J_a$ sofort (wenn keine höherpriore Task aktiv)
		\State $c_s \leftarrow c_s - $ verbrauchte Zeit
		\State
		\State \textbf{Wenn Server aktiviert und $c_s > 0$, aber keine Anfrage:}
		\State Behalte Budget $c_s$ (nicht verwerfen!)
		\State Server schläft bis zur nächsten Anfrage oder nächsten Periode
	\end{algorithmic}
\end{algorithm}

\subsection{Warum DS KEINE äquivalente periodische Task ist}

Der DS verletzt Annahme A5 (Tasks suspendieren sich nicht), da er sich
suspendiert, obwohl er der höchstpriore Task wäre. Dies hat eine \emph{niedrigere}
Schedulierbarkeitsschranke zur Folge.

\begin{satz}[DS Schedulierbarkeitsschranke -- Strosnider, Lehoczky, Sha]
	\label{satz:ds}
	Die Schedulierbarkeitsschranke für $n$ periodische Tasks mit einem
	Deferrable Server mit Auslastung $U_s$ ist:
	\[
		U_p \leq n \left(\left(\frac{U_s+2}{2U_s+1}\right)^{1/n} - 1\right).
	\]
	Das Minimum über alle $U_s$ ergibt $U_p^* \approx 0{,}652$ bei $U_s^* \approx 0{,}186$.
\end{satz}

\begin{proof}
	Im schlimmsten Fall führt der DS dreimal innerhalb einer Periode $T_1$ der
	höchstprioren periodischen Task aus: am Ende seiner eigenen Periode (aufgehoben)
	und zu Beginn der nächsten Periode. Daraus folgt die Worst-Case-Parameterisierung:
	\[
		C_s = \frac{T_1 - T_s}{2}, \quad C_1 = T_2 - T_1, \ldots, \quad
		C_n = \frac{3T_s + T_1 - 2T_n}{2}.
	\]
	Minimierung der Auslastung über $R_i = T_{i+1}/T_i$ analog zum RM-Beweis
	(Abschnitt~\ref{satz:rm_lub}) ergibt:
	\[
		K = \frac{U_s + 2}{2U_s + 1}, \qquad U_{\text{lub}} = U_s + n(K^{1/n} - 1).
	\]
	Da DS die periodischen Tasks stärker beeinträchtigt als PS (wegen der
	Kapazitätserhaltung), gilt $U_{\text{lub}}^{DS} \leq U_{\text{lub}}^{PS}$
	für $U_s < 0{,}4$. \hfill$\blacksquare$
\end{proof}

\subsubsection{Vergleich PS vs.\ DS}

Für $U_p = 0{,}6$: $U_s^{\max}(\text{PS}) = 2/P - 1$ und $U_s^{\max}(\text{DS}) = (2-P)/(2P-1)$
mit $P = \prod_{i}(U_i+1)$. DS erlaubt typischerweise weniger Serverbandbreite
als PS, liefert aber deutlich bessere aperiodische Antwortzeiten.

\subsection{Aperiodische Antwortzeit unter DS}

Sei $c_s(t)$ das verbleibende Budget zur Ankunftszeit $r_a$ und $C_0 = \min(c_s(t), \text{next}(r_a) - r_a)$ die im aktuellen Intervall bedienbare Zeit:
\[
	R_a = \Delta_a + C_a - C_0 + F_a(T_s - C_s),
	\quad F_a = \left\lceil \frac{C_a - C_0}{C_s} \right\rceil.
\]

\textbf{Optimalitätsklassifikation:} DS ist \emph{nicht scharf optimal}
bezüglich aperiodischer Antwortzeiten, aber besser als PS. Die niedrigere
Schedulierbarkeitsschranke gegenüber PS ist der Preis für bessere aperiodische
Responsivität.

\textbf{Laufzeit:} $O(\log n)$ pro Ereignis. Budgetverwaltung: $O(1)$ pro
Aktivierung und Anfrage.

\section{Priority Exchange Server}

\subsection{Informelle Beschreibung}

Der Priority Exchange (PE) Server verbessert die Schedulierbarkeitsschranke
gegenüber DS, indem Budget nicht einfach erhalten, sondern \emph{getauscht} wird:
Wenn keine aperiodischen Anfragen vorliegen, tauscht der Server seine Ausführungszeit
gegen die Laufzeit der höchstprioren aktiven periodischen Task. Das Budget wird
dabei auf das Prioritätsniveau der getauschten Task degradiert.

\begin{algorithm}[H]
	\caption{Priority Exchange Server}
	\begin{algorithmic}[1]
		\State \textbf{Zu Beginn der Serverperiode:} $c_s \leftarrow C_s$ bei Serverpriorität
		\State
		\State \textbf{Wenn Server höchste Priorität hat und $c_s > 0$:}
		\If{aperiodische Anfrage vorhanden}
		\State Bediene Anfrage für $\min(C_{\text{ape}}, c_s)$ Einheiten
		\ElsIf{periodische Task $\tau_k$ aktiv}
		\State Führe $\tau_k$ für 1 Einheit aus (Tausch)
		\State $C_S^{d_k} \leftarrow C_S^{d_k} + 1$ (Budget auf $\tau_k$-Niveau)
		\State $c_s \leftarrow c_s - 1$ (Serverbudget reduzieren)
		\Else
		\State Leerlauf; $c_s \leftarrow c_s - 1$
		\EndIf
	\end{algorithmic}
\end{algorithm}

\subsection{Schedulierbarkeitsnachweis}

\begin{satz}[PE Schedulierbarkeit]
	\label{satz:pe}
	Ein Task-Set mit $n$ periodischen Tasks und einem PE-Server mit Auslastung
	$U_s$ ist unter RM schedulierbar, wenn:
	\[
		U_p \leq n\left(\left(\frac{2}{U_s+1}\right)^{1/n} - 1\right).
	\]
	Diese Schranke ist \emph{identisch} mit der des Polling Servers (Satz~\ref{satz:ps}).
\end{satz}

\begin{proof}
	Im schlimmsten Fall verhält sich PE wie eine periodische Task mit Periode $T_s$
	und Ausführungszeit $C_s$: Die Tauschoperationen verschieben Rechenzeit nur
	zwischen Prioritätsniveaus, ändern aber nicht die Gesamtausführungszeit. Das
	erweiterte System hat damit dieselbe effektive Auslastungsstruktur wie PS,
	und der Liu-Layland-Bound gilt analog. \hfill$\blacksquare$
\end{proof}

\subsection{Optimalitätsklassifikation und Vergleich}

PE ist damit gegenüber DS \emph{besser schedulierbar}, aber auf Kosten
komplexerer Implementierung. Für eine gegebene periodische Last $U_p$:
\begin{itemize}
	\item PE erlaubt mehr Serverbandbreite als DS.
	\item Die aperiodischen Antwortzeiten von PE sind vergleichbar mit DS.
\end{itemize}

\textbf{Laufzeit:} $O(\log n)$ pro Ereignis, zusätzlich $O(1)$ für Tauschoperationen.
Die Tracking-Verwaltung mehrerer Budget-Niveaus erfordert $O(n)$ Speicher.

\section{Sporadic Server -- Optimale Schedulierbarkeitsschranke unter RM}

\subsection{Informelle Beschreibung}

Der Sporadic Server (SS), von Sprunt, Sha und Lehoczky, hat die Schlüsseleigenschaft:
\emph{Das verbrauchte Budget wird erst $T_s$ Zeiteinheiten nach dem Verbrauch
wiederhergestellt}. Dadurch verhält sich SS -- im Gegensatz zu DS -- exakt wie
eine normale periodische Task in Bezug auf die Schedulierbarkeitsschranke.

\textbf{Replenishment-Regeln:}
\begin{itemize}
	\item \textbf{Wiederauffüllzeitpunkt} RT: Wird gesetzt, wenn SS aktiv wird
	      ($P_{\text{exe}} \geq P_s$) und $c_s > 0$. Sei $t_A$ dieser Zeitpunkt:
	      $\mathrm{RT} = t_A + T_s$.
	\item \textbf{Wiederauffüllmenge} RA: Wird gesetzt, wenn SS idle wird ($P_{\text{exe}} < P_s$)
	      oder $c_s = 0$. Sei $t_I$ dieser Zeitpunkt:
	      $\mathrm{RA} = $ verbrauchte Kapazität in $[t_A, t_I]$.
	\item Zum Zeitpunkt RT: $c_s \leftarrow c_s + \mathrm{RA}$.
\end{itemize}

\begin{algorithm}[H]
	\caption{Sporadic Server}
	\begin{algorithmic}[1]
		\State $c_s \leftarrow C_s$; RT $\leftarrow \infty$; $t_A \leftarrow -\infty$
		\State
		\State \textbf{Wenn SS aktiv wird ($P_{\text{exe}} \geq P_s$) und $c_s > 0$:}
		\If{RT $= \infty$}
		\State $t_A \leftarrow$ aktueller Zeitpunkt; RT $\leftarrow t_A + T_s$
		\EndIf
		\State
		\State \textbf{Wenn SS idle wird ($P_{\text{exe}} < P_s$) oder $c_s = 0$:}
		\State RA $\leftarrow c_s(t_A) - c_s(t_I)$ \Comment{Verbrauchtes Budget in $[t_A, t_I]$}
		\State Plane Wiederauffüllung RA zu Zeitpunkt RT
		\State RT $\leftarrow \infty$ \Comment{Nächster Aktivierungszeitpunkt wird neu gesetzt}
		\State
		\State \textbf{Zum geplanten Wiederauffüllzeitpunkt RT:}
		\State $c_s \leftarrow c_s + \mathrm{RA}$
	\end{algorithmic}
\end{algorithm}

\subsection{Hauptsatz: SS verhält sich wie eine periodische Task}

\begin{theorem}[Sprunt, Sha, Lehoczky -- SS Äquivalenz]
	\label{thm:ss}
	Ein Task-Set, das mit einer periodischen Task $\tau_s(C_s, T_s)$ schedulierbar
	ist, bleibt schedulierbar, wenn $\tau_s$ durch einen Sporadic Server mit
	denselben Parametern $(C_s, T_s)$ ersetzt wird.
\end{theorem}

\begin{proof}
	Wir analysieren das Verhalten von SS in den drei möglichen Fällen (Abbildung
	nach Buttazzo \cite{Buttazzo2011}):

	\textbf{Fall 1 -- Kein Verbrauch in $[t_A, t_I]$:}
	In diesem Intervall ist keine aperiodische Anfrage aktiv. SS behält sein Budget,
	aber keine Wiederauffüllung findet vor $t_I + T_s$ statt. Das bedeutet: In
	$[t_A, t_I + T_s]$ kann SS maximal $C_s$ Zeiteinheiten für aperiodische Tasks
	verbrauchen -- identisch mit einer periodischen Task $\tau_s$ mit verzögertem
	Release-Zeitpunkt $t_I$ statt $t_A$. Da verzögerte Releases die Antwortzeiten
	anderer Tasks nicht verschlechtern, bleibt Schedulierbarkeit erhalten.

	\textbf{Fall 2 -- Vollständiger Verbrauch:}
	$c_s$ wird in $[t_A, t_I]$ vollständig verbraucht ($\mathrm{RA} = C_s$).
	Wiederauffüllung bei $t_A + T_s$. SS verhält sich exakt wie $\tau_s$ mit
	Release-Zeit $t_A$.

	\textbf{Fall 3 -- Teilweiser Verbrauch:}
	$\mathrm{RA} = C_R$ (verbrauchter Teil) wird bei $t_A + T_s$ wiederaufgefüllt;
	verbleibendes Budget $C_s - C_R$ bleibt erhalten. SS verhält sich wie zwei
	periodische Tasks $\tau_x(C_R, T_s)$ und $\tau_y(C_s - C_R, T_s)$, wobei
	$\tau_x$ bei $t_A$ und $\tau_y$ bei $t_I$ released wird. Da $\tau_y$ verzögert
	ist, schadet es der Schedulierbarkeit nicht (analog zu Fall 1).

	In allen Fällen ist der Beitrag von SS zur Auslastung $U_s = C_s/T_s$, genau
	wie eine periodische Task. \hfill$\blacksquare$
\end{proof}

\subsection{Schedulierbarkeitsanalyse}

\begin{satz}[SS Schedulierbarkeitsschranke]
	\label{satz:ss_sched}
	Eine Menge von $n$ periodischen Tasks mit Auslastung $U_p$ und ein
	Sporadic Server mit Auslastung $U_s$ sind unter RM schedulierbar, wenn:
	\[
		U_p + U_s \leq (n+1)\left(2^{1/(n+1)} - 1\right).
	\]
	Mit dem Hyperbolic Bound:
	\[
		\prod_{i=1}^n (U_i + 1) \leq \frac{2}{U_s + 1}.
	\]
	Die maximale Serverauslastung, die die Schedulierbarkeit erhält, ist:
	\[
		U_s^{\max} = \frac{2-P}{P}, \quad P = \prod_{i=1}^n (U_i + 1).
	\]
\end{satz}

\begin{proof}
	Aus Theorem~\ref{thm:ss} folgt, dass SS wie eine periodische Task mit
	$U_s = C_s/T_s$ behandelt werden kann. Das erweiterte System
	$\{\tau_1,\ldots,\tau_n, \tau_s\}$ hat $n+1$ Tasks; der Hyperbolic Bound
	(Theorem~\ref{thm:hyperbolic}) für $n+1$ Tasks lautet $\prod_{i=1}^{n+1}(U_i+1) \leq 2$.
	Einsetzen von $U_{n+1} = U_s$ liefert die Behauptung. \hfill$\blacksquare$
\end{proof}

\subsection{Optimalitätsklassifikation und Laufzeit}

\textbf{Optimalitätsklassifikation:} SS ist \emph{schwach optimal} für den
Klasse fester Prioritätsserver: Es gibt keinen Server unter RM, der eine
höhere Schedulierbarkeitsschranke \emph{und} bessere aperiodische Antwortzeiten
gleichzeitig bietet (Theorem~\ref{thm:nooptserver}).

\textbf{Laufzeit:} Jede Ereignisbehandlung (Serveraktivierung, Idle-Wechsel,
Wiederauffüllung) erfordert $O(1)$ Operationen. Pro aperiodischer Anfrage
entstehen höchstens $\lceil C_a/C_s \rceil$ Wiederauffüllungen, daher
Gesamtoverhead pro Anfrage $O(\lceil C_a/C_s \rceil)$.

\section{Slack Stealing -- Maximale Responsivität}

\subsection{Informelle Beschreibung}

Der Slack Stealing Algorithmus (Lehoczky \& Ramos-Thuel 1992) nutzt nicht
eine feste Serverbandbreite, sondern ``stiehlt'' alle verfügbaren Slack-Zeiten
der periodischen Tasks für aperiodische Anfragen.

\textbf{Grundidee:} Wenn $\tau_i$ eine Slack-Zeit $\text{slack}_i(t) > 0$ hat,
bedeutet das, dass sie nicht \emph{sofort} ausgeführt werden muss. Der
Slack Stealer kann diese Zeit für aperiodische Anfragen nutzen, ohne die
Deadline von $\tau_i$ zu gefährden.

\begin{algorithm}[H]
	\caption{Slack Stealing (vereinfachte Darstellung)}
	\begin{algorithmic}[1]
		\State Berechne offline $A(0, t)$ = Slack-Funktion für alle $t \in [0, H]$
		\State
		\State \textbf{Wenn aperiodische Anfrage $J_a$ bei $t=r_a$ ankommt:}
		\State Berechne aktuellen Slack $A(r_a, t')$ für $t' > r_a$
		\State Finde frühestes $t'$ mit $A(r_a, t') \geq C_a$
		\State Plane $J_a$ so früh wie möglich innerhalb verfügbarer Slack-Zeiten
		\State
		\State \textbf{Wenn keine aperiodische Anfrage:}
		\State Normale RM-Einplanung der periodischen Tasks
	\end{algorithmic}
\end{algorithm}

\subsection{Slack-Funktion}

\begin{definition}[Slack-Funktion $A(s,t)$]
	$A(s, t)$ ist die maximale Rechenzeitdauer, die im Intervall $[s, t]$
	aperiodischen Tasks zugewiesen werden kann, ohne dass periodische Tasks
	ihre Deadlines verfehlen:
	\[
		A(s, t) = (t-s) - \sum_{i=1}^n \left\lceil \frac{t-s}{T_i} \right\rceil C_i.
	\]
	$A(s, \cdot)$ ist eine nichtfallende Treppenfunktion mit Sprüngen an den
	Zeitpunkten, an denen Slack entsteht.
\end{definition}

\begin{satz}[Slack Stealing -- Schedulierbarkeit]
	Wenn das periodische Task-Set unter RM schedulierbar ist ($U_p \leq U_{\text{lub}}$),
	gefährdet Slack Stealing die Schedulierbarkeit nie: Aperiodische Tasks werden
	ausschließlich in Zeiten eingeplant, in denen kein periodischer Task laufen muss.
\end{satz}

\begin{proof}
	Per Konstruktion: $A(s, t)$ gibt exakt die Zeitdauer an, die für periodische
	Tasks \emph{nicht} benötigt wird. Jede aperiodische Zuteilung innerhalb von
	$A(s, t)$ lässt mindestens genug Prozessorzeit für alle periodischen Tasks
	bis zu ihren Deadlines. \hfill$\blacksquare$
\end{proof}

\subsection{Nicht-Existenz optimaler Server -- Fundamentales Ergebnis}

\begin{theorem}[Tia, Liu, Shankar 1995 -- Kein optimaler Server]
	\label{thm:nooptserver}
	Für jeden festen Prioritätsscheduler und jede Ankunftsdisziplin für
	aperiodische Tasks existiert kein Online-Algorithmus, der die Antwortzeit
	\emph{jeder einzelnen} aperiodischen Anfrage minimiert.
\end{theorem}

\begin{proof}
	\textbf{Durch Gegenbeispiel:} Betrachte $\tau_1$ ($C=1, T=3$), $\tau_2$ ($C=1, T=4$),
	$\tau_3$ ($C=1, T=6$) unter RM. Slack tritt bei $t=2$ auf.

	\textbf{Ankunft $J_1$ bei $t=2$ ($C=1$):}
	\begin{itemize}
		\item Option A: $J_1$ sofort (bei $t=2$) einplanen $\Rightarrow$ $J_1$ fertig bei $t=3$, minimale Antwortzeit für $J_1$.
		\item Option B: $J_1$ auf $t=3$ verschieben $\Rightarrow$ $J_1$ fertig bei $t=4$.
	\end{itemize}

	\textbf{Ankunft $J_2$ bei $t=3$ ($C=1$):}
	\begin{itemize}
		\item Bei Option A (A sofort): Kein Slack in $[3,6]$; $J_2$ frühestens bei $t=7$.
		\item Bei Option B (A verschoben): Slack bei $t=4$; $J_2$ fertig bei $t=5$.
	\end{itemize}

	Option A minimiert $J_1$, aber nicht $J_2$.
	Option B minimiert $J_2$, aber nicht $J_1$.
	Kein Algorithmus kann beide gleichzeitig minimieren. \hfill$\blacksquare$
\end{proof}

\textbf{Laufzeit Slack Stealing:} Die Offline-Berechnung der Slack-Funktion
hat Komplexität $O(N)$ mit $N$ = Anzahl Deadlines in der Hyperperiode. Die
Online-Aktualisierung bei jeder aperiodischen Ankunft hat Komplexität $O(n)$.

\section{Dimensionierung von Servern}

\begin{satz}[Optimale Server-Dimensionierung]
	Für ein gegebenes periodisches Task-Set mit Hyperbolic-Produkt $P = \prod_i(U_i+1)$
	sind die maximalen Serverauslastungen:
	\begin{align*}
		U_s^{\max}(\text{PS/PE/SS}) &= \frac{2-P}{P}, \\
		U_s^{\max}(\text{DS}) &= \frac{2-P}{2P-1}.
	\end{align*}
	Die optimale Serverperiode (beste Antwortzeit bei gegebenem $U_s$) ist
	$T_s = T_1$ (höchste Priorität), mit $C_s = U_s \cdot T_s$.
\end{satz}

\section{Zusammenfassender Vergleich der Fixed-Priority Server}

\begin{table}[H]
	\centering
	\caption{Vergleich der Fixed-Priority Server -- Vollständige Übersicht}
	\resizebox{\textwidth}{!}{%
	\begin{tabular}{lllll}
		\toprule
		\textbf{Server} & \textbf{Budget-Mgmt.} & \textbf{Sched.-Schranke} & \textbf{Antwortzeit} & \textbf{Overhead} \\
		\midrule
		Background & Keine & Wie RM ($U_p \leq U_{\text{lub}}$) & Schlecht & $O(1)$ \\
		Polling (PS) & Verwerfen & Wie RM$+$1 Task & Mittel & $O(\log n)$ \\
		Deferrable (DS) & Erhalten & Niedriger als PS! & Gut & $O(\log n)$ \\
		Priority Exch.\ (PE) & Tauschen & Wie PS & Gut & $O(n \log n)$ \\
		Sporadic (SS) & Verzög.\ Auffüllen & \textbf{Wie PS} & Gut & $O(\log n)$ \\
		Slack Stealing & Slack aller Tasks & \textbf{Maximal} & \textbf{Opt.*} & $O(N)$ \\
		\bottomrule
	\end{tabular}%
	}

	\smallskip
	\noindent\small{*Kein optimaler Server existiert (Thm.~\ref{thm:nooptserver}).
	Slack Stealing minimiert Antwortzeiten soweit möglich.}
\end{table}

%=============================================================
\chapter{Dynamic Priority Server -- Dynamische Prioritäts-Server unter EDF}
\label{chap:dynserver}
%=============================================================

\section{Einführung und Motivation}

Gegenüber Fixed-Priority-Servern aus Kapitel~\ref{chap:server} bieten dynamische
Prioritäts-Server einen entscheidenden Vorteil: Da EDF den Prozessor bis $U=1$
auslasten kann, steht für aperiodische Tasks eine weit größere Bandbreite zur
Verfügung. Für eine Menge periodischer Tasks mit Auslastung $U_p$ verbleiben
$(1-U_p)$ für aperiodische Server -- gegenüber deutlich weniger unter RM.

\begin{bemerkung}[Vorteil dynamischer Server gegenüber RM-basierten Servern]
	Beispiel: $U_1 = U_2 = 0{,}3$, also $U_p = 0{,}6$.
	\begin{itemize}
		\item Unter RM+Sporadic Server: $U_s^{\max} = 2/P - 1$ mit
		      $P = (0{,}3+1)^2 = 1{,}69$, also $U_s^{\max} \approx 0{,}18$.
		\item Unter EDF: $U_s^{\max} = 1 - U_p = 0{,}40$.
	\end{itemize}
	EDF erlaubt damit mehr als doppelt so viel Bandbreite für aperiodische Tasks.
\end{bemerkung}

Alle in diesem Kapitel vorgestellten Server arbeiten unter folgenden Annahmen:
\begin{itemize}
	\item Periodische Tasks $\tau_i$ haben harte Deadlines mit $D_i = T_i$.
	\item Aperiodische Tasks $J_i$ haben keine Deadlines, sollen aber so früh wie
	      möglich fertiggestellt werden.
	\item Alle periodischen Tasks starten synchron bei $t=0$.
	\item Die Ausführungszeit jeder aperiodischen Task ist bekannt.
\end{itemize}

\section{Dynamic Priority Exchange Server (DPE)}

\subsection{Grundidee und Algorithmus}

Der Dynamic Priority Exchange Server (DPE), entwickelt von Spuri und Buttazzo
\cite{Spuri1996}, überträgt das Prinzip des Fixed-Priority-Exchange-Servers auf
EDF: Der Server tauscht seine Laufzeit mit niederprioren periodischen Tasks
(d.\,h.\ solchen mit späterer Deadline), anstatt sie zu verbrauchen. Dadurch
bleibt Kapazität erhalten, auch wenn keine aperiodischen Anfragen vorliegen.

\textbf{Informelle Beschreibung:}

\begin{itemize}
	\item Der DPE-Server hat eine Periode $T_s$ und Kapazität $C_s$.
	\item Zu Beginn jeder Serverperiode mit Deadline $d$ wird die aperiodische
	      Kapazität $C_S^d \leftarrow C_s$ gesetzt.
	\item Wenn die höchste Priorität im System eine aperiodische Kapazität $C$
	      hat, wird entweder eine aperiodische Task (falls vorhanden) oder die
	      periodische Task mit kürzester Deadline ausgeführt. Im letzteren Fall
	      wird deren Laufzeit in eine aperiodische Kapazität auf dem Prioritätsniveau
	      der periodischen Task übertragen (Tausch).
	\item Leerlauf verbraucht die Kapazität bis auf null.
\end{itemize}

\begin{algorithm}[H]
	\caption{DPE Server -- Ereignisbehandlung}
	\begin{algorithmic}[1]
		\State \textbf{Bei Serverperiodenbeginn mit Deadline $d$:}
		\State $C_S^d \leftarrow C_s$
		\State
		\State \textbf{Bei Ausführung der höchsten Priorität (aperiodische Kapazität $C$):}
		\If{aperiodische Anfrage $J$ vorhanden}
		\State Bediene $J$ bis Fertigstellung oder Kapazität $C$ erschöpft
		\ElsIf{periodische Task $\tau_k$ mit frühester Deadline bereit}
		\State Führe $\tau_k$ für 1 Zeiteinheit aus
		\State $C_S^{d_k} \leftarrow C_S^{d_k} + 1$; $C \leftarrow C - 1$
		\Comment{Tausch: Kapazität auf Niveau $d_k$ übertragen}
		\Else
		\State Leerlauf; $C \leftarrow C - 1$
		\EndIf
	\end{algorithmic}
\end{algorithm}

\subsection{Schedulierbarkeitsnachweis -- Korrektheit}

\begin{theorem}[Spuri \& Buttazzo -- DPE Schedulierbarkeit]
	\label{thm:dpe}
	Gegeben eine Menge periodischer Tasks mit Auslastung $U_p$ und ein DPE-Server
	mit Auslastung $U_s$, ist das gesamte System unter EDF genau dann schedulierbar,
	wenn
	\[
		U_p + U_s \leq 1.
	\]
\end{theorem}

\begin{proof}
	\textbf{Konstruktion des äquivalenten Schedules $\sigma'$:}

	Sei $\sigma$ ein von DPE produzierter Schedule. Wir konstruieren $\sigma'$
	wie folgt:
	\begin{itemize}
		\item Ersetze den DPE-Server durch eine periodische Task $\tau_s$
		      mit Periode $T_s$ und WCET $C_s$: In $\sigma'$ läuft $\tau_s$ genau dann,
		      wenn in $\sigma$ Serverkapazität verbraucht wird.
		\item Periodische Instanzen, die in $\sigma$ wegen Deadline-Tausch
		      vorgezogen wurden, werden in $\sigma'$ nach hinten verschoben.
	\end{itemize}

	\textbf{Eigenschaft von $\sigma'$:} Das resultierende $\sigma'$ ist das
	EDF-Schedule für die Menge $\{\tau_1,\ldots,\tau_n,\tau_s\}$ mit
	Gesamtauslastung $U_p + U_s$ -- unabhängig von aperiodischen Anfragen.

	\textbf{Äquivalenz $\sigma \leftrightarrow \sigma'$:} In $\sigma$ sind die
	Fertigstellungszeiten jeder periodischen Instanz $\leq$ denen in $\sigma'$,
	da DPE Instanzen vorziehen kann. Also: $\sigma'$ machbar $\Rightarrow$ $\sigma$
	machbar. Umgekehrt ist $\sigma'$ der DPE-Schedule ohne aperiodische Anfragen,
	also gilt auch: $\sigma$ machbar $\Rightarrow$ $\sigma'$ machbar.

	Da $\sigma'$ ein reines EDF-Schedule über periodische Tasks mit Auslastung
	$U_p + U_s$ ist, gilt nach Theorem~\ref{thm:edf_schedulierbar}:
	$\sigma'$ machbar $\iff$ $U_p + U_s \leq 1$. \hfill$\blacksquare$
\end{proof}

\subsection{Laufzeitanalyse}

\begin{satz}[DPE Laufzeit]
	Die Laufzeit von DPE pro Ereignis (Ankunft, Fertigstellung, Kapazitätserschöpfung)
	beträgt $O(\log n)$ für die Verwaltung der EDF-Prioritätswarteschlange als Heap.
	Pro Serverperiode entsteht ein Overhead von $O(C_s)$ durch die Tauschoperationen,
	da jede Zeiteinheit maximal eine Tauschoperation erzeugt.
\end{satz}

\subsection{Optimalitätsklassifikation von DPE}

DPE ist \textbf{nicht scharf optimal} bezüglich der Antwortzeit aperiodischer Tasks:
Es gibt Szenarien, in denen ein clairvoyanter Algorithmus bessere Antwortzeiten
liefern würde. DPE ist jedoch \textbf{schwach optimal} in dem Sinne, dass für jeden
Auslastungswert $U_p + U_s \leq 1$ die Schedulierbarkeit aller periodischen Tasks
garantiert wird -- genau wie ein äquivalenter rein-periodischer EDF-Schedule.

\section{Dynamic Sporadic Server (DSS)}

\subsection{Grundidee und Algorithmus}

Der Dynamic Sporadic Server (DSS), ebenfalls von Spuri und Buttazzo \cite{Spuri1996},
erweitert den Sporadic Server auf dynamische Prioritäten. Die Kapazität $C_s$ wird
nicht zu festen Zeiten aufgefüllt, sondern nach dem Verbrauch.

\textbf{Deadline- und Wiederauffüllungsregeln:}

\begin{itemize}
	\item Wenn der Server aktiv wird ($C_s > 0$ und aperiodische Anfrage vorhanden)
	      zum Zeitpunkt $t_A$: Setze $\mathrm{RT} = d_s = t_A + T_s$.
	\item Wenn der Server idle wird (Anfrage abgeschlossen oder $C_s = 0$)
	      zum Zeitpunkt $t_I$: Plane Wiederauffüllung $\mathrm{RA} = \text{(verbrauchte Kapazität in }[t_A, t_I])$ zu Zeitpunkt $\mathrm{RT}$.
\end{itemize}

\begin{algorithm}[H]
	\caption{DSS Server -- Kapazitätsverwaltung}
	\begin{algorithmic}[1]
		\State \textbf{Initialisierung:} $c_s \leftarrow C_s$, $d_s \leftarrow \infty$
		\State
		\State \textbf{Wenn aperiodische Anfrage $J$ ankommt und $c_s > 0$:}
		\State $t_A \leftarrow$ aktueller Zeitpunkt
		\State $\mathrm{RT} \leftarrow t_A + T_s$; $d_s \leftarrow t_A + T_s$
		\State Bediene $J$ mit Deadline $d_s$; decrementiere $c_s$ bei Ausführung
		\State
		\State \textbf{Wenn Server idle wird (Zeitpunkt $t_I$):}
		\State $\mathrm{RA} \leftarrow$ verbrauchte Kapazität in $[t_A, t_I]$
		\State Plane Wiederauffüllung $c_s \mathrel{+}= \mathrm{RA}$ bei Zeitpunkt $\mathrm{RT}$
	\end{algorithmic}
\end{algorithm}

\subsection{Schedulierbarkeitsnachweis}

\begin{lemma}[DSS verhält sich wie eine periodische Task]
	\label{lem:dss_periodic}
	In jedem Zeitintervall $[t_1, t_2]$, in dem $t_1$ der Aktivierungszeitpunkt
	des Servers ist, gilt:
	\[
		C_{\mathrm{ape}}(t_1, t_2) \leq \left\lfloor \frac{t_2 - t_1}{T_s} \right\rfloor C_s.
	\]
\end{lemma}

\begin{proof}
	Da die Wiederauffüllungsmenge stets gleich der verbrauchten Menge ist und
	das Wiederauffüllen frühestens $T_s$ nach dem Aktivierungszeitpunkt erfolgt,
	kann in jedem $T_s$-langen Intervall maximal $C_s$ Zeit für aperiodische
	Tasks verbraucht werden. Die Behauptung folgt durch Induktion. \hfill$\blacksquare$
\end{proof}

\begin{theorem}[Spuri \& Buttazzo -- DSS Schedulierbarkeit]
	\label{thm:dss}
	Gegeben eine Menge periodischer Tasks mit $U_p$ und ein DSS mit $U_s$, ist
	das System schedulierbar genau dann, wenn $U_p + U_s \leq 1$.
\end{theorem}

\begin{proof}
	\textbf{Hinreichend ($\Rightarrow$):} Angenommen, $U_p + U_s \leq 1$, aber
	es gibt einen Überlauf bei $t$. Dann gibt es ein $t' < t$ mit kontinuierlicher
	Auslastung. Für die Gesamtanforderung $C$ in $[t', t]$ gilt nach
	Lemma~\ref{lem:dss_periodic}:
	\[
		C \leq \sum_{i=1}^n \left\lfloor \frac{t-t'}{T_i} \right\rfloor C_i + C_{\mathrm{ape}}
		\leq (t-t') U_p + (t-t') U_s = (t-t')(U_p+U_s) \leq (t-t').
	\]
	Dies widerspricht $t-t' < C$. Widerspruch.

	\textbf{Notwendig ($\Leftarrow$):} Da DSS sich wie eine periodische Task mit
	$U_s$ verhält, folgt aus dem EDF-Schedulierbarkeitstheorem direkt $U_p + U_s \leq 1$. \hfill$\blacksquare$
\end{proof}

\subsection{Laufzeitanalyse und Optimalitätsklassifikation}

\begin{satz}[DSS Laufzeit]
	Jede Ereignisbehandlung (Ankunft, Fertigstellung, Wiederauffüllung) erfordert
	$O(\log n)$ für Heap-Operationen. Pro aperiodischer Anfrage entstehen maximal
	$\lceil C_a/C_s \rceil$ Wiederauffüllungen, sodass der Gesamtoverhead
	pro Anfrage in $O(\lceil C_a/C_s \rceil \cdot \log n)$ liegt.
\end{satz}

DSS ist wie DPE \textbf{nicht scharf optimal} bezüglich Antwortzeiten, aber
\textbf{schedulierbar-korrekt}: Periodische Tasks werden genau dann garantiert, wenn
$U_p + U_s \leq 1$.

\section{Total Bandwidth Server (TBS)}

\subsection{Grundidee und Algorithmus}

Der Total Bandwidth Server (TBS) ist der einfachste und effizienteste der hier
vorgestellten Server \cite{Spuri1996}. Statt eine eigene Periode zu verwalten,
weist er jeder aperiodischen Anfrage eine virtuelle Deadline zu, die die gesamte
zugeordnete Bandbreite $U_s$ sofort zuteilt.

\textbf{Deadlinezuweisung:} Wenn die $k$-te aperiodische Anfrage zum Zeitpunkt
$r_k$ mit Ausführungszeit $C_k$ ankommt:
\[
	d_k = \max(r_k,\; d_{k-1}) + \frac{C_k}{U_s},
\]
wobei $d_0 = 0$. Die Anfrage wird dann wie eine periodische Instanz mit Deadline
$d_k$ in die EDF-Warteschlange eingereiht.

\begin{algorithm}[H]
	\caption{Total Bandwidth Server}
	\begin{algorithmic}[1]
		\Require Serverauslastung $U_s$, initiale Deadline $d_0 \leftarrow 0$
		\State
		\State \textbf{Bei Ankunft der $k$-ten Anfrage $(r_k, C_k)$:}
		\State $d_k \leftarrow \max(r_k, d_{k-1}) + C_k / U_s$
		\State Füge $(C_k, d_k)$ in EDF-Warteschlange ein
		\State EDF-Schedule läuft weiter
	\end{algorithmic}
\end{algorithm}

\begin{beispiel}[TBS Deadlinezuweisung]
	$U_p = 3/4$, $U_s = 1/4$. Periodische Tasks $\tau_1$ ($C_1=3, T_1=6$),
	$\tau_2$ ($C_2=2, T_2=8$).

	Aperiodische Anfragen:
	\begin{itemize}
		\item $J_1$: $r_1=3$, $C_1=1$. $d_1 = 3 + 1/0{,}25 = 7$.
		\item $J_2$: $r_2=9$, $C_2=1$. $d_2 = \max(9,7) + 1/0{,}25 = 13$.
		\item $J_3$: $r_3=14$, $C_3=1$. $d_3 = \max(14,13) + 1/0{,}25 = 18$.
	\end{itemize}
	Jede Anfrage wird als EDF-Instanz mit den genannten Deadlines eingeplant.
\end{beispiel}

\subsection{Korrektheitsbeweis: Auslastungsbeschränkung}

\begin{lemma}[TBS Auslastungsschranke]
	\label{lem:tbs_util}
	In jedem Intervall $[t_1, t_2]$ ist die von TBS bediente Ausführungszeit
	$C_{\mathrm{ape}}$ durch $(t_2 - t_1) \cdot U_s$ nach oben beschränkt:
	\[
		C_{\mathrm{ape}} = \sum_{t_1 \leq r_k,\; d_k \leq t_2} C_k \leq (t_2 - t_1) U_s.
	\]
\end{lemma}

\begin{proof}
	Seien $J_{k_1}, \ldots, J_{k_2}$ die aperiodischen Tasks, die zu $[t_1, t_2]$
	beitragen (d.\,h.\ $r_{k_j} \geq t_1$ und $d_{k_j} \leq t_2$). Aus der
	TBS-Deadlineregel folgt durch Teleskopsumme:
	\[
		\sum_{k=k_1}^{k_2} C_k = \sum_{k=k_1}^{k_2} [d_k - \max(r_k, d_{k-1})] U_s
		\leq [d_{k_2} - \max(r_{k_1}, d_{k_1-1})] U_s \leq (t_2 - t_1) U_s.
	\]
	Hierbei wird genutzt, dass jede Klammer $[d_k - \max(\cdot)] \leq d_k - d_{k-1}$
	und die Summe teleskopisch kollabiert. \hfill$\blacksquare$
\end{proof}

\begin{theorem}[Spuri \& Buttazzo -- TBS Schedulierbarkeit]
	\label{thm:tbs}
	Gegeben periodische Tasks mit $U_p$ und ein TBS mit $U_s$, ist das System
	unter EDF genau dann schedulierbar, wenn $U_p + U_s \leq 1$.
\end{theorem}

\begin{proof}
	\textbf{Hinreichend:} Angenommen $U_p + U_s \leq 1$, aber es gibt einen
	Überlauf bei $t$. Sei $[t', t]$ das längste Intervall kontinuierlicher Auslastung
	vor $t$, in dem nur Tasks mit Deadline $\leq t$ ausgeführt werden. Die Gesamtanforderung:
	\[
		C \leq \sum_i \left\lfloor \frac{t-t'}{T_i} \right\rfloor C_i + C_{\mathrm{ape}}
		\leq (t-t') U_p + (t-t') U_s = (t-t')(U_p + U_s) \leq t-t'.
	\]
	Widerspruch zu $t-t' < C$.

	\textbf{Notwendig:} Wenn aperiodische Anfragen periodisch mit Periode $T_s$
	und Ausführungszeit $C_s = T_s U_s$ ankommen, verhält sich TBS wie eine
	periodische Task. Dann folgt aus EDF-Schedulierbarkeit $U_p + U_s \leq 1$. \hfill$\blacksquare$
\end{proof}

\subsection{Optimalitätsklassifikation von TBS}

TBS ist \textbf{nicht scharf optimal} im Sinne minimaler Antwortzeiten: Es kann
aperiodische Anfragen mit nichtminimaler Antwortzeit einplanen. Es gibt jedoch
eine Erweiterung TB* (Total Bandwidth, optimiert), die durch iterative Deadline-Verkürzung die minimale Antwortzeit annähert.

\begin{satz}[TB* Optimalität -- Buttazzo \& Sensini]
	Sei $\sigma$ ein machbarer EDF-Schedule und $f_k$ die Fertigstellungszeit
	einer aperiodischen Anfrage $J_k$ mit Deadline $d_k$. Wenn $f_k = d_k$, dann
	ist $f_k = f_k^*$ (minimale erreichbare Fertigstellungszeit).
\end{satz}

\textbf{Laufzeit von TBS:} $O(1)$ pro Ankunft (Deadlineberechnung) plus $O(\log n)$
für EDF-Heap-Einfügung. Gesamtoverhead vernachlässigbar.

\section{Constant Bandwidth Server (CBS)}

\subsection{Definition und Motivation}

Der Constant Bandwidth Server (CBS) von Abeni und Buttazzo \cite{Buttazzo2011}
ist die robusteste der vorgestellten Techniken: Er garantiert temporale Isolation,
d.\,h.\ eine Task kann den Server nicht über seine zugewiesene Bandbreite $U_s$
hinaus belasten -- selbst wenn sie mehr Rechenzeit beansprucht als erwartet.
Dies unterscheidet CBS grundlegend von TBS, der die Überschreitung von WCETs
nicht abfangen kann.

\begin{definition}[CBS-Parameter und Zustandsvariablen]
	Ein CBS ist durch ein Maximum-Budget $Q_s$ und eine Periode $T_s$ mit
	Bandbreite $U_s = Q_s/T_s$ charakterisiert. Zustandsvariablen:
	\begin{itemize}
		\item $c_s$: aktuelles Budget ($0 < c_s \leq Q_s$ stets)
		\item $d_{s,k}$: aktuelle Serverdeadline
		\item $k$: Zähler der erzeugten Chunks
	\end{itemize}
	Initialisierung: $d_{s,0} \leftarrow 0$, $c \leftarrow 0$, $n \leftarrow 0$, $k \leftarrow 0$.
\end{definition}

\subsection{CBS-Algorithmus (Pseudocode)}

\begin{algorithm}[H]
	\caption{Constant Bandwidth Server (CBS)}
	\begin{algorithmic}[1]
		\State \textbf{Bei Ankunft von Job $J_j$ zum Zeitpunkt $r_j$:}
		\State $n \leftarrow n+1$; füge $J_j$ in Serverwarteschlange ein
		\If{$n = 1$} \Comment{Server war idle}
		\If{$r_j + (c/Q_s) \cdot T_s \geq d_{s,k}$} \Comment{Regel 1: neues Deadline}
		\State $k \leftarrow k+1$; $a_k \leftarrow r_j$; $d_{s,k} \leftarrow r_j + T_s$; $c \leftarrow Q_s$
		\Else \Comment{Regel 2: alte Deadline}
		\State $k \leftarrow k+1$; $a_k \leftarrow r_j$; $d_{s,k} \leftarrow d_{s,k-1}$
		\EndIf
		\EndIf
		\State
		\State \textbf{Bei Fertigstellung von $J_j$:}
		\State $n \leftarrow n-1$
		\If{$n \neq 0$} starte nächsten Job mit Deadline $d_{s,k}$ \EndIf
		\State
		\State \textbf{Bei jeder Ausführungseinheit ($c = c - 1$):}
		\If{$c = 0$} \Comment{Regel 3: Budget erschöpft}
		\State $k \leftarrow k+1$; $a_k \leftarrow$ aktueller Zeitpunkt
		\State $d_{s,k} \leftarrow d_{s,k-1} + T_s$; $c \leftarrow Q_s$
		\EndIf
	\end{algorithmic}
\end{algorithm}

\subsection{Isolationseigenschaft -- formaler Beweis}

\begin{theorem}[CBS Isolationseigenschaft -- Abeni \& Buttazzo]
	\label{thm:cbs_isolation}
	Die CPU-Auslastung eines CBS $S$ mit Parametern $(Q_s, T_s)$ beträgt
	exakt $U_s = Q_s/T_s$, unabhängig von Ausführungszeiten und Ankunftsmustern.
\end{theorem}

\begin{proof}
	Durch Regel 3 wird das Budget $c_s$ bei jedem Verbrauch sofort auf $Q_s$
	aufgefüllt und die Serverdeadline um $T_s$ nach hinten verschoben. In jedem
	Intervall der Länge $T_s$ verbraucht der CBS genau $Q_s$ Zeiteinheiten für
	Serverausführung. Die Ausführungsrate ist damit konstant $Q_s/T_s = U_s$.

	Formaler: Betrachte zwei aufeinanderfolgende Serverdeadlines $d_{s,k}$ und
	$d_{s,k+1} = d_{s,k} + T_s$. In $[d_{s,k-1}, d_{s,k}]$ wird genau $Q_s$
	Budget verbraucht (sonst würde Regel 3 nicht auslösen). Damit ist die
	Bandbreite exakt $Q_s/T_s$ für jeden Beobachtungszeitraum. \hfill$\blacksquare$
\end{proof}

\begin{lemma}[Schedulierbarkeit unter CBS]
	Gegeben periodische harte Tasks mit $U_p$ und $m$ CBS-Server mit Gesamtbandbreite
	$U_s = \sum_{i=1}^m U_{s,i}$, ist das System unter EDF genau dann schedulierbar,
	wenn $U_p + U_s \leq 1$.
\end{lemma}

\begin{proof}
	Aus Theorem~\ref{thm:cbs_isolation} folgt, dass jeder CBS sich wie eine
	periodische Task mit Auslastung $U_{s,i}$ verhält. Die Gesamtauslastung
	beträgt $U_p + \sum_i U_{s,i} = U_p + U_s$, und EDF-Schedulierbarkeit
	gilt genau dann, wenn dieser Wert $\leq 1$. \hfill$\blacksquare$
\end{proof}

\subsection{Worst-Case-Antwortzeit und optimale CBS-Parameter}

Die worst-case Antwortzeit $R_i$ eines Jobs mit Ausführungszeit $C_i$ unter CBS mit
Budget $Q_s$ und Periode $T_s$ beträgt:
\[
	R_i = C_i + \left\lfloor \frac{C_i}{Q_s} \right\rfloor (T_s - Q_s).
\]

Das Budget $Q_s$, das die mittlere Antwortzeit minimiert (unter Berücksichtigung
eines Kontext-Switch-Overheads $\epsilon$), ergibt sich zu:
\[
	T_s^{\mathrm{opt}} = \frac{1}{U_s} \left( \epsilon + \sqrt{\frac{\epsilon \cdot C^{\mathrm{avg}}}{1-U_s}} \right).
\]

\subsection{Vergleich der dynamischen Server}

\begin{table}[H]
	\centering
	\caption{Vergleich dynamischer Prioritäts-Server}
	\label{tab:dynserver}
	\begin{tabular}{p{2.5cm}p{2cm}p{2cm}p{2cm}p{3cm}}
		\toprule
		\textbf{Server} & \textbf{Sched.\ korrekt} & \textbf{Antwort\-zeit} & \textbf{Overhead} & \textbf{Besonderheit} \\
		\midrule
		DPE  & $U_p+U_s\leq 1$ & Gut & Mittel & Deadline-Tausch \\
		DSS  & $U_p+U_s\leq 1$ & Gut & Mittel & Dynamische Wiederbefüllung \\
		TBS  & $U_p+U_s\leq 1$ & Sehr gut & \textbf{Minimal} & Einfachste Impl. \\
		EDL  & $U_p \leq 1$ & \textbf{Optimal} & Hoch & Idle-Zeit-Ausnutzung \\
		IPE  & $U_p \leq 1$ & Nahezu opt. & Mittel & Vorberechnete Zeiten \\
		CBS  & $U_p+U_s\leq 1$ & Gut & Gering & \textbf{Temporale Isolation} \\
		\bottomrule
	\end{tabular}
\end{table}

\section{EDL-Server -- Optimale aperiodische Antwortzeiten}

Der Earliest Deadline Late (EDL) Server nutzt die Leerlaufzeiten eines EDL-Schedules
(EDF-Variante, in der Tasks so spät wie möglich eingeplant werden) für aperiodische
Anfragen.

\begin{theorem}[Chetto \& Chetto -- EDL Optimalität]
	Für jede Online-Algorithmus $A$ und jeden aperiodischen Job $J_i$ gilt:
	\[
		f^{\mathrm{EDL-Server}}(J_i) \leq f^A(J_i).
	\]
	Der EDL-Server minimiert die Fertigstellungszeit jeder aperiodischen Anfrage.
\end{theorem}

\begin{proof}[Beweisidee]
	Da EDL in jedem Intervall $[0,t]$ die maximale Leerlaufzeit unter allen
	Schedulingalgorithmen garantiert (Theorem von Chetto \& Chetto), stehen
	für aperiodische Tasks stets die meisten Freiräume zur Verfügung. Jeder andere
	Algorithmus kann somit höchstens dieselben Freiräume nutzen. \hfill$\blacksquare$
\end{proof}

\begin{satz}[EDL-Server Laufzeit]
	Die Berechnung der Leerlaufzeiten unter EDL hat Komplexität $O(N \cdot n)$,
	wobei $n$ die Anzahl der periodischen Tasks und $N$ die Anzahl der
	Instanzen in der Hyperperiode ist. Im schlimmsten Fall kann $N$ sehr groß
	sein (bei nicht-harmonischen Perioden), was EDL für viele praktische Systeme
	ungeeignet macht.
\end{satz}

\section{Nicht-Existenz optimaler Server -- Fundamentales Ergebnis}

\begin{theorem}[Tia, Liu, Shankar 1995]
	Für jeden Task-Set mit fester Prioritätszuweisung und jede aperiodische
	Ankunftsdisziplin existiert kein Online-Algorithmus, der die Antwortzeit
	\emph{jeder} aperiodischen Anfrage minimiert.
\end{theorem}

\begin{proof}[Beweis durch Gegenbeispiel]
	Betrachte drei periodische Tasks: $\tau_1$ ($C=1, T=3$), $\tau_2$ ($C=1, T=4$),
	$\tau_3$ ($C=1, T=6$), scheduliert unter RM. Bei Ankunft von $J_1$ ($C=1$)
	bei $t=2$ hat jeder Algorithmus zwei Optionen:
	\begin{itemize}
		\item Option 1: Führe $J_1$ sofort aus. Wenn dann $J_2$ bei $t=3$ ankommt,
		      ist kein Slack bis $t=6$ verfügbar; $J_2$ terminiert frühestens bei $t=7$.
		\item Option 2: Verschiebe $J_1$ auf $t=3$. Dann hat $J_2$ einen freien
		      Slot bei $t=4$; $J_2$ terminiert bei $t=5$. Aber $J_1$ ist nicht minimal.
	\end{itemize}
	In keiner Option sind beide Antwortzeiten gleichzeitig minimal. \hfill$\blacksquare$
\end{proof}

\begin{bemerkung}
	Dieses fundamentale Ergebnis zeigt, dass die in diesem Kapitel präsentierten
	Server keine scharf optimalen Lösungen sind: Es gibt kein ``perfektes'' Server-
	Verfahren. Die vorgestellten Algorithmen sind stattdessen Kompromisse zwischen
	Antwortzeit, Overhead, Speicherbedarf und Implementierungskomplexität.
\end{bemerkung}

%=============================================================
\chapter{Ressourcenzugriffsprotokolle -- Prioritätsinversion und Lösungsstrategien}
\label{chap:ressourcen}
%=============================================================

\section{Das Prioritätsinversionsproblem}

\begin{definition}[Kritischer Abschnitt]
	Ein \textbf{kritischer Abschnitt} ist ein Codebereich, der unter gegenseitiger
	Ausschluss-Bedingung ausgeführt wird. Nur eine Task darf einen kritischen
	Abschnitt, der eine Ressource $R_k$ schützt, gleichzeitig betreten.
	Gesichert durch Semaphor $S_k$ mit \texttt{wait}($S_k$) beim Eintritt und
	\texttt{signal}($S_k$) beim Austritt.
\end{definition}

\begin{definition}[Prioritätsinversion]
	\textbf{Prioritätsinversion} tritt auf, wenn eine hochpriore Task $\tau_H$
	blockiert ist, während eine niederpriore Task $\tau_M$ (ohne gemeinsame
	Ressource mit $\tau_H$) ausgeführt wird. Die Blockierzeit ist dann nicht
	auf die Dauer des kritischen Abschnitts begrenzt, sondern hängt von der
	gesamten Ausführungszeit aller Tasks mit $P_M < P < P_H$ ab.
\end{definition}

\begin{beispiel}[Klassische Prioritätsinversion -- 3 Tasks]
	$\tau_1 > \tau_2 > \tau_3$ in Priorität. $\tau_1$ und $\tau_3$ teilen Ressource $R$.
	\begin{enumerate}
		\item $t_0$: $\tau_3$ startet, betritt kritischen Abschnitt auf $R$.
		\item $t_1$: $\tau_2$ kommt, preemptiert $\tau_3$ (höhere Priorität).
		\item $t_2$: $\tau_1$ kommt, preemptiert $\tau_2$ (höhere Priorität).
		\item $t_3$: $\tau_1$ versucht $R$ zu betreten -- blockiert (von $\tau_3$ gehalten).
		\item $t_4$: $\tau_2$ läuft weiter und preemptiert $\tau_3$ -- $\tau_1$ wartet!
	\end{enumerate}
	$\tau_1$ wartet also nicht auf $\tau_3$'s kritischen Abschnitt, sondern auf die
	\emph{gesamte} Ausführungszeit von $\tau_2$. Im schlimmsten Fall unbeschränkt.
\end{beispiel}

\textbf{Terminologie:}
\begin{itemize}
	\item $B_i$: maximale Blockierzeit von $\tau_i$
	\item $\delta_{j,k}$: Dauer des längsten kritischen Abschnitts von $\tau_j$ auf $S_k$
	\item $C(S_k) = \max\{P_i \mid \tau_i$ benutzt $S_k\}$: Prioritätsgrenze von $S_k$
	\item $\gamma_i$: Menge der kritischen Abschnitte, die $\tau_i$ blockieren können
\end{itemize}

\section{Non-Preemptive Protocol (NPP)}

\subsection{Algorithmus}

\begin{definition}[NPP]
	Bei \textbf{NPP} wird die Priorität einer Task beim Eintritt in jeden
	kritischen Abschnitt auf das systemweite Maximum angehoben:
	\[
		p_i(R_k) = \max_h \{P_h\}.
	\]
	Bei Verlassen des kritischen Abschnitts wird $p_i$ auf $P_i$ zurückgesetzt.
\end{definition}

\begin{algorithm}[H]
	\caption{Non-Preemptive Protocol (NPP)}
	\begin{algorithmic}[1]
		\State \textbf{wait($S_k$):}
		\State $p_i \leftarrow \max_h\{P_h\}$ \Comment{Priorität auf Maximum}
		\State Belege Ressource $R_k$
		\State
		\State \textbf{signal($S_k$):}
		\State Gib $R_k$ frei
		\State $p_i \leftarrow P_i$ \Comment{Nominal-Priorität wiederherstellen}
	\end{algorithmic}
\end{algorithm}

\subsection{Blockierzeit unter NPP}

Da unter NPP nie zwei Tasks gleichzeitig in kritischen Abschnitten sein können
(die Prioritätsanhebung verhindert Preemption), ist die maximale Blockierzeit:
\[
	B_i = \max\{\delta_{j,k} - 1 \mid P_j < P_i,\; k=1,\ldots,m\}.
\]

\textbf{Nachteil:} NPP verursacht \emph{unnötige} Blockierung: Wenn $\tau_1$ keine
gemeinsame Ressource mit $\tau_3$ hat, aber $\tau_3$ in einem langen kritischen
Abschnitt ist, wird $\tau_1$ dennoch blockiert.

\section{Priority Inheritance Protocol (PIP)}

\subsection{Protokolldefinition}

\begin{definition}[PIP -- Sha, Rajkumar, Lehoczky 1990]
	Im \textbf{Priority Inheritance Protocol} erbt eine Task $\tau_j$, die eine
	Ressource hält, vorübergehend die höchste Priorität aller Tasks, die auf diese
	Ressource warten:
	\[
		p_j(R_k) = \max\bigl\{P_j,\; \max\{P_h \mid \tau_h \text{ wartet auf } R_k\}\bigr\}.
	\]
	Diese Vererbung ist \emph{transitiv}: Wenn $\tau_3$ von $\tau_2$ blockiert wird
	und $\tau_2$ von $\tau_1$, erbt $\tau_3$ die Priorität von $\tau_1$.
\end{definition}

\begin{algorithm}[H]
	\caption{Priority Inheritance Protocol (PIP)}
	\begin{algorithmic}[1]
		\State \textbf{wait($S_k$):}
		\If{$R_k$ von $\tau_j$ gehalten ($P_j < P_i$)}
		\State $\tau_i$ wird blockiert
		\State $p_j \leftarrow \max(p_j, P_i)$ \Comment{Prioritätsvererbung}
		\State Setze $\tau_j$ fort mit neuer Priorität
		\Else
		\State $\tau_i$ belegt $R_k$
		\EndIf
		\State
		\State \textbf{signal($S_k$):}
		\State Gib $R_k$ frei; wecke höchstpriore wartende Task
		\If{keine Tasks mehr von $\tau_j$ blockiert}
		\State $p_j \leftarrow P_j$ \Comment{Nominal-Priorität}
		\Else
		\State $p_j \leftarrow \max\{P_h \mid \tau_j \text{ blockiert } \tau_h\}$
		\EndIf
	\end{algorithmic}
\end{algorithm}

\subsection{Korrektheit und Eigenschaften}

\begin{lemma}
	Ein Semaphor $S_k$ kann bei $\tau_i$ nur dann Push-Through-Blockierung
	auslösen, wenn $S_k$ sowohl von einer Task mit $P < P_i$ als auch von einer
	Task mit $P > P_i$ benutzt wird.
\end{lemma}

\begin{lemma}
	Transitive Prioritätsvererbung tritt nur bei \emph{geschachtelten}
	kritischen Abschnitten auf.
\end{lemma}

\begin{theorem}[Sha, Rajkumar, Lehoczky -- PIP Blockierzeit]
	\label{thm:pip}
	Unter PIP kann $\tau_i$ für höchstens $\alpha_i = \min(l_i, s_i)$ kritische
	Abschnitte blockiert werden, wobei $l_i$ die Anzahl der niederprioren Tasks
	und $s_i$ die Anzahl der eindeutigen Semaphore sind, die $\tau_i$ blockieren
	können.
\end{theorem}

\begin{proof}
	\textbf{Bound durch $l_i$ (Lemma~7.3 in Buttazzo):} $\tau_i$ kann durch
	$\tau_j$ nur blockiert werden, solange $\tau_j$ in einem kritischen Abschnitt
	ist. Sobald $\tau_j$ diesen verlässt, kann er $\tau_i$ nicht mehr blockieren.
	Pro niederpriorer Task also maximal ein Blockiervorgang.

	\textbf{Bound durch $s_i$ (Lemma~7.4 in Buttazzo):} Da Semaphore binär sind,
	kann genau eine Task pro Semaphor aktiv in einem kritischen Abschnitt sein.
	Nach dem Freigeben kann $S_k$ $\tau_i$ nicht mehr blockieren. Also maximal
	eine Blockierung pro Semaphor. Das Minimum beider Bounds gilt. \hfill$\blacksquare$
\end{proof}

\subsection{Blockierzeitberechnung unter PIP}

\begin{algorithmic}[1]
	\Require Tasks geordnet nach absteigender Priorität; keine Schachtelung
	\For{$i = 1$ \textbf{to} $n-1$}
	\State $B_i^l \leftarrow \sum_{j:\,P_j < P_i} \max_k\{\delta_{j,k}-1 \mid C(S_k) \geq P_i\}$
	\State $B_i^s \leftarrow \sum_{k=1}^m \max_{j:\,P_j < P_i}\{\delta_{j,k}-1 \mid C(S_k) \geq P_i\}$
	\State $B_i \leftarrow \min(B_i^l, B_i^s)$
	\EndFor
	\State $B_n \leftarrow 0$
\end{algorithmic}

Komplexität: $O(mn^2)$.

\section{Priority Ceiling Protocol (PCP)}

\subsection{Protokolldefinition}

\begin{definition}[PCP -- Sha, Rajkumar, Lehoczky 1990]
	Jeder Ressource $R_k$ wird eine \textbf{Prioritätsgrenze} (Priority Ceiling)
	zugewiesen:
	\[
		C(R_k) = \max\{P_i \mid \tau_i \text{ benutzt } R_k\}.
	\]
	Eine Task $\tau_i$ darf $R_k$ nur betreten, wenn ihre aktive Priorität
	\emph{strikt höher} ist als die Prioritätsgrenze aller aktuell gesperrten
	Ressourcen:
	\[
		p_i > \max_{R_h \text{ gesperrt}} C(R_h).
	\]
\end{definition}

\begin{algorithm}[H]
	\caption{Priority Ceiling Protocol (PCP)}
	\begin{algorithmic}[1]
		\State \textbf{wait($S_k$):}
		\If{$p_i > \max\{C(R_h) \mid R_h \text{ belegt}\}$}
		\State $\tau_i$ belegt $R_k$; $p_i(R_k) \leftarrow C(R_k)$
		\Comment{Priorität auf Ceiling}
		\Else
		\State $\tau_i$ wird blockiert (von Task mit höchster Ceiling-Priorität)
		\EndIf
		\State
		\State \textbf{signal($S_k$):}
		\State Gib $R_k$ frei; $p_i \leftarrow P_i$ \Comment{Nominal-Priorität}
		\State Wecke höchstpriore wartende Task
	\end{algorithmic}
\end{algorithm}

\subsection{Eigenschaften und Nachweis}

\begin{theorem}[PCP Eigenschaften]
	\label{thm:pcp}
	Unter PCP gilt:
	\begin{enumerate}
		\item Jede Task $\tau_i$ wird höchstens einmal blockiert.
		\item Deadlocks sind unmöglich.
		\item Gegenseitiger Ausschluss ist gewährleistet.
	\end{enumerate}
\end{theorem}

\begin{proof}
	\textbf{Zu 1 (Einmalige Blockierung):} Angenommen, $\tau_i$ wird zweimal
	blockiert -- einmal durch $\tau_a$ auf $R_a$ und einmal durch $\tau_b$ auf $R_b$.
	Damit beide gleichzeitig ausgeführt werden könnten, müsste eine von ihnen die
	andere innerhalb eines kritischen Abschnitts preemptiert haben. Aber nach PCP-Regel
	erfordert das $p > C(R_{\text{gehalten}})$, also $P_a > C(R_b)$ oder $P_b > C(R_a)$.
	Da $C(R_b) \geq P_i > P_a$ und $C(R_a) \geq P_i > P_b$ (beide haben niedrigere
	Priorität als $\tau_i$), führt das zu einem Widerspruch.

	\textbf{Zu 2 (Deadlock-Freiheit):} Ein Deadlock erfordert zyklisches Warten.
	Nach PCP kann $\tau_j$ $R_k$ nur belegen, wenn $p_j > C(R_k)$. Für einen Zyklus
	$\tau_a \to \tau_b \to \tau_a$ müsste $P_a > C(R_b) \geq P_b > C(R_a) \geq P_a$
	gelten -- Widerspruch. \hfill$\blacksquare$
\end{proof}

\subsection{Schedulierbarkeitsanalyse unter PCP}

Die Schedulierbarkeit unter RM+PCP kann mit folgender erweiterter Response Time Analysis
geprüft werden (Sha et al.):
\[
	R_i = C_i + B_i + \sum_{j \in hp(i)} \left\lceil \frac{R_i}{T_j} \right\rceil C_j,
\]
wobei $B_i = \max\{\delta_{j,k}-1 \mid P_j < P_i,\; C(R_k) \geq P_i\}$ die maximale
Blockierzeit ist. Da $\tau_i$ höchstens einmal blockiert wird, ist $B_i$ durch einen
einzigen kritischen Abschnitt beschränkt.

\section{Stack Resource Policy (SRP)}

\begin{definition}[Preemptionsgrenze (Preemption Level)]
	Unter SRP wird jedem Task $\tau_i$ ein statischer \textbf{Preemptionslevel}
	$\pi_i$ zugewiesen (unter RM gleich der Priorität). Jeder Ressource $R_k$ wird
	eine \textbf{Ressourcengrenze} zugewiesen:
	\[
		\Pi(R_k) = \max\{\pi_i \mid \tau_i \text{ benutzt } R_k\}.
	\]
	Ein \textbf{System Ceiling} $\Pi_{\text{sys}}(t)$ ist das Maximum aller
	Ressourcengrenzen der aktuell gesperrten Ressourcen.
\end{definition}

\begin{definition}[SRP-Zugriffsregel]
	$\tau_i$ darf den Prozessor belegen (preemptieren), wenn $\pi_i > \Pi_{\text{sys}}(t)$.
	$\tau_i$ darf $R_k$ betreten, ohne zusätzliche Prüfung -- die Zugangskontrolle
	erfolgt bereits beim Preemptionszeitpunkt.
\end{definition}

\begin{satz}[SRP Eigenschaften]
	Unter SRP gilt:
	\begin{itemize}
		\item Jede Task wird höchstens \emph{einmal} blockiert (bei der Aktivierung,
		      nicht beim Ressourcenzugriff).
		\item Stack-Sharing ist möglich: Alle Tasks können denselben Stack nutzen,
		      da niemals zwei Tasks gleichzeitig im kritischen Abschnitt sein können.
		\item SRP ist sowohl für feste als auch für dynamische Prioritäten anwendbar
		      (OSEK, AUTOSAR, POSIX-ähnliche Systeme).
	\end{itemize}
\end{satz}

\begin{proof}[Beweis der einmaligen Blockierung unter SRP]
	$\tau_i$ kann nur dann nicht aktiviert werden, wenn $\pi_i \leq \Pi_{\text{sys}}(t)$.
	Da $\Pi_{\text{sys}}$ monoton abnimmt (Ressourcen werden freigegeben), kann
	$\tau_i$ nach seiner Aktivierung nie wieder durch den System Ceiling blockiert
	werden. Damit ist genau eine Blockierphase möglich, nämlich vor der ersten
	Ausführung. \hfill$\blacksquare$
\end{proof}

\subsection{Vergleich der Protokolle}

\begin{table}[H]
	\centering
	\caption{Vergleich der Ressourcenzugriffsprotokolle}
	{\small
	\begin{tabular}{p{2.0cm}p{1.9cm}p{1.9cm}p{2.0cm}p{3.0cm}}
		\toprule
		\textbf{Protokoll} & \textbf{Max.\ Block.} & \textbf{Deadlock-frei} & \textbf{Stack-Sharing} & \textbf{Bemerkung} \\
		\midrule
		NPP    & 1 Abschnitt & Ja & Ja & Unnötige Blockierung \\
		HLP    & 1 Abschnitt & Ja & Nein & Ceiling bei Eintritt \\
		PIP    & $\min(l_i,s_i)$ & Nein! & Nein & Komplexe Analyse \\
		PCP    & 1 Abschnitt & Ja & Nein & Guter Kompromiss \\
		SRP    & 1 Abschnitt & Ja & \textbf{Ja} & Optimal für Einzel-CPU \\
		\bottomrule
	\end{tabular}
	}
\end{table}

%=============================================================
\chapter{Limited Preemptive Scheduling -- Begrenzte Preemption}
\label{chap:limited}
%=============================================================

\section{Motivation: Kosten und Nutzen der Preemption}

Vollständig preemptives Scheduling maximiert die Schedulierbarkeit, verursacht aber
erheblichen Overhead: Cache-Misses durch zerstörte Lokalität, Pipeline-Flushes,
Kontext-Switch-Kosten. In vielen eingebetteten Systemen (AUTOSAR, ISO~26262) ist
dieser Overhead signifikant.

\begin{definition}[Preemptionskosten-Komponenten]
	Jede Preemption verursacht vier Kostenkategorien:
	\begin{enumerate}
		\item \textbf{Scheduling-Overhead}: Suspendierung, Warteschlangenoperation, Dispatch.
		\item \textbf{Pipeline-Kosten}: Flush beim Unterbrechen, Refill beim Fortsetzen.
		\item \textbf{Cache-Kosten (CRPD)}: Cache-Related Preemption Delay --
		      Reload der Cache-Lines, die die preemptierende Task verdrängt hat.
		\item \textbf{Bus-Kosten}: Zusätzliche RAM-Zugriffe durch Cache-Misses.
	\end{enumerate}
	Empirisch: Auf PowerPC MPC7410 mit 2 MB L2-Cache kann CRPD bis zu 33\,\%
	der nicht-preemptiven WCET betragen \cite{Buttazzo2011}.
\end{definition}

Zentrales Ergebnis dieses Kapitels: Für ein präemptiv machbares Task-Set lässt sich
stets ein Preemptionsschema finden, das die Schedulierbarkeit erhält und gleichzeitig
die Preemptionsanzahl minimiert.

\section{Nicht-Preemptives Scheduling}

\subsection{Blockierzeit und Schedulierbarkeit}

Ohne Preemption ergibt sich für jede Task $\tau_i$ eine Blockierzeit durch die
längste niederpriore Task:
\begin{equation}
	B_i = \max_{j:\,P_j < P_i}\{C_j - 1\}.
	\label{eq:np_blocking}
\end{equation}

\begin{satz}[Nicht-preemptiv -- $U_{\text{lub}} = 0$]
	Unter nicht-preemptivem RM und EDF gilt $U_{\text{lub}} = 0$: Es gibt Task-Sets
	mit beliebig kleiner Auslastung, die unter nicht-preemptivem Scheduling
	nicht schedulierbar sind.
\end{satz}

\begin{proof}
	Betrachte $\tau_1$ ($C_1 = \varepsilon$, $T_1 = 2\varepsilon$) und
	$\tau_2$ ($C_2 = T_1 + \varepsilon = 3\varepsilon$, $T_2 = M$ für großes $M$).
	Auslastung: $U = \varepsilon/(2\varepsilon) + 3\varepsilon/M = 0{,}5 + 3\varepsilon/M$,
	beliebig nah an 0{,}5 für großes $M$. Aber: Nicht-preemptiv kann $\tau_2$
	nicht unterbrochen werden, also $\tau_1$ muss $C_2 = 3\varepsilon > T_1 = 2\varepsilon$
	warten -- Deadline-Verletzung. \hfill$\blacksquare$
\end{proof}

\subsection{Machbarkeitsanalyse für nicht-preemptives Scheduling}

Das \emph{Self-Pushing-Phänomen} (Bril et al.\ 2009) zeigt, dass unter
nicht-preemptivem Scheduling die worst-case Antwortzeit nicht beim ersten Job
(kritischer Augenblick), sondern bei einem späteren Job auftreten kann. Daher
muss die Analyse innerhalb der längsten Level-$i$-Aktivitätsperiode $L_i$
durchgeführt werden:
\begin{align}
	L_i^{(0)} &= B_i + C_i, \nonumber\\
	L_i^{(s)} &= B_i + \sum_{h:\,P_h \geq P_i} \left\lfloor \frac{L_i^{(s-1)}}{T_h} \right\rfloor C_h.
	\label{eq:level_i_period}
\end{align}

Für alle Jobs $k \in [1, K_i]$ mit $K_i = \lfloor L_i/T_i \rfloor$ wird der Startzeitpunkt
$s_{i,k}$ iterativ berechnet:
\[
	s_{i,k}^{(\ell+1)} = B_i + (k-1)C_i + \sum_{h:\,P_h > P_i} \left(\left\lfloor \frac{s_{i,k}^{(\ell)}}{T_h} \right\rfloor + 1\right) C_h.
\]

Die Antwortzeit ergibt sich zu $R_i = \max_{k \in [1,K_i]}\{s_{i,k} + C_i - (k-1)T_i\}$.

\section{Preemption Thresholds}

\subsection{Modell und Algorithmus}

\begin{definition}[Preemption Threshold -- Wang \& Saksena 1999]
	Jede Task $\tau_i$ erhält zwei Prioritätsniveaus:
	\begin{itemize}
		\item Nominale Priorität $P_i$ (für die Ready-Queue und Einplanung),
		\item Preemption Threshold $\theta_i \geq P_i$ (für die Ausführung).
	\end{itemize}
	$\tau_h$ kann $\tau_i$ \emph{während der Ausführung} nur preemptieren, wenn
	$P_h > \theta_i$.
\end{definition}

\begin{algorithm}[H]
	\caption{Assign Minimum Preemption Thresholds (Wang \& Saksena)}
	\begin{algorithmic}[1]
		\Require Task-Set $\mathcal{T}$ mit $\{C_i, T_i, D_i, P_i\}$, preemptiv machbar
		\For{$i \leftarrow n$ \textbf{downto} $1$} \Comment{Von niedrigster zu höchster Priorität}
		\State $\theta_i \leftarrow P_i$
		\State Berechne $R_i$ nach Gl.~\eqref{eq:rta_threshold}
		\While{$R_i > D_i$}
		\State $\theta_i \leftarrow \theta_i + 1$
		\If{$\theta_i > P_1$} \Return INFEASIBLE \EndIf
		\State Berechne $R_i$ nach Gl.~\eqref{eq:rta_threshold}
		\EndWhile
		\EndFor
		\State \Return FEASIBLE
	\end{algorithmic}
\end{algorithm}

\subsection{Machbarkeitsanalyse}

Blockierzeit unter Preemption Thresholds:
\[
	B_i = \max\{C_j - 1 \mid P_j < P_i \leq \theta_j\}.
\]
Antwortzeit (innerhalb der längsten Level-$i$-Aktivitätsperiode):
\begin{equation}
	R_i = S_i + C_i + \sum_{h:\,P_h > \theta_i}\left(\left\lfloor \frac{R_i}{T_h} \right\rfloor - \left\lfloor \frac{S_i}{T_h} \right\rfloor - 1\right) C_h,
	\label{eq:rta_threshold}
\end{equation}
wobei $S_i = B_i + \sum_{h:\,P_h > P_i}\left(\lfloor S_i/T_h \rfloor + 1\right) C_h$ der Startzeitpunkt ist.

\begin{satz}[Optimalität des Threshold-Algorithmus]
	Der Algorithmus ist optimal: Wenn irgendeine Threshold-Zuweisung das Task-Set
	machbar macht, findet er eine solche Zuweisung.
\end{satz}

\subsection{Optimalitätsklassifikation}

Preemption Thresholds sind \textbf{nicht scharf optimal}: Es gibt Task-Sets, die
weder vollständig preemptiv noch vollständig nicht-preemptiv, aber mit geeigneten
Thresholds schedulierbar sind. Innerhalb der Klasse der Threshold-Scheduler ist der
Algorithmus \textbf{schwach optimal} (findet Lösung, wenn eine existiert).

\section{Deferred Preemptions -- Aufgeschobene Preemption}

\subsection{Floating Model und Aktivierungsbasiertes Modell}

\begin{definition}[Maximales nicht-preemptives Intervall]
	Jede Task $\tau_i$ hat ein maximales nicht-preemptives Intervall $q_i$, in dem
	sie nicht unterbrochen werden kann. Im \textbf{Floating Model} platziert der
	Programmierer explizite Grenzen; im \textbf{Activation-Triggered Model}
	wird beim Eintreffen einer höherprioren Task ein Timer auf $q_i$ gesetzt.
\end{definition}

Blockierzeit: $B_i = \max_{j:\,P_j < P_i}\{q_j - 1\}$.

\subsection{Längste machbare nicht-preemptive Intervalle}

Die \emph{Blocking Tolerance} $\beta_i$ ist die maximale Blockierzeit, die $\tau_i$
tolerieren kann, ohne eine Deadline zu verpassen:
\[
	\beta_i = \max_{t \in TS_i}\{t - W_i(t)\},
\]
wobei $W_i(t)$ die Level-$i$-Arbeitslast und $TS_i$ die Testmenge ist.

\begin{theorem}[Yao, Buttazzo, Bertogna -- Längste nicht-preemptive Intervalle]
	Die längsten nicht-preemptiven Intervalle $Q_i$, die die Schedulierbarkeit
	erhalten, können rekursiv berechnet werden:
	\[
		Q_i = \min\{Q_{i-1},\; \beta_{i-1} + 1\},
	\]
	mit $Q_1 = \infty$ und $\beta_1 = D_1 - C_1$.
\end{theorem}

\begin{proof}
	$Q_i$ ist das minimale $\beta_k + 1$ über alle Tasks mit höherer Priorität als
	$\tau_i$. Durch Einsetzen der Rekursionsdefinition:
	\[
		\min_{k:\,P_k > P_i}\{\beta_k + 1\}
		= \min\Bigl\{\min_{k:\,P_k > P_{i-1}}\{\beta_k + 1\},\; \beta_{i-1} + 1\Bigr\}
		= \min\{Q_{i-1},\; \beta_{i-1} + 1\}. \quad\blacksquare
	\]
\end{proof}

\begin{algorithm}[H]
	\caption{Berechnung der längsten nicht-preemptiven Intervalle}
	\begin{algorithmic}[1]
		\State $\beta_1 \leftarrow D_1 - C_1$; $Q_1 \leftarrow \infty$
		\For{$i \leftarrow 2$ \textbf{to} $n$}
		\State $Q_i \leftarrow \min\{Q_{i-1},\; \beta_{i-1} + 1\}$
		\State Berechne $\beta_i$ aus $W_i(t)$
		\EndFor
	\end{algorithmic}
\end{algorithm}

\subsection{Laufzeitanalyse}

Berechnung aller $Q_i$: $O(n)$ Iterationen, jede mit $O(|\mathit{TS}_i|)$ Aufwand
für die Berechnung von $\beta_i$. Im ungünstigsten Fall ist $|\mathit{TS}_i|$
pseudopolynomiell (abhängig von den Perioden).

\section{Task Splitting -- Kooperatives Scheduling}

\subsection{Modell und optimale Preemptionspunkt-Auswahl}

\begin{definition}[Task Splitting]
	Task $\tau_i$ wird in $m_i$ nicht-preemptive \emph{Chunks} (Teilaufgaben) zerlegt.
	Preemption ist nur an Chunkgrenzen (Effective Preemption Points, EPPs) erlaubt.
\end{definition}

\begin{algorithm}[H]
	\caption{SelectEPP -- Optimale Preemptionspunktauswahl (Bertogna et al.\ 2011)}
	\begin{algorithmic}[1]
		\Require Task $\tau_i$ mit Basisblöcken $\delta_1,\ldots,\delta_N$, maximales NPR $Q$
		\State $\hat{C}_0 \leftarrow 0$; $\mathit{Prev}_0 \leftarrow \{\delta_0\}$
		\For{$k \leftarrow 1$ \textbf{to} $N$}
		\State Entferne aus $\mathit{Prev}_k$ alle $\delta_j$, die Bed.~\eqref{eq:npr_cond} verletzen
		\If{$\mathit{Prev}_k = \emptyset$} \Return INFEASIBLE \EndIf
		\State $\hat{C}_k \leftarrow \min_{\delta_j \in \mathit{Prev}_k}\left\{\hat{C}_{j-1} + \xi_{j-1} + \sum_{\ell=j}^{k} b_\ell\right\}$
		\State Speichere $\delta_{\mathrm{prev}}(\delta_k)$; $\mathit{Prev}_k \leftarrow \mathit{Prev}_k \cup \{\delta_k\}$
		\EndFor
		\State Rekonstruiere EPPs rückwärts von $\delta_{\mathrm{prev}}(\delta_N)$
		\State \Return FEASIBLE
	\end{algorithmic}
\end{algorithm}

Dabei ist die Bedingung für nicht-preemptive Regionen:
\begin{equation}
	\xi_{j-1} + \sum_{\ell=j}^{k} b_\ell \leq Q.
	\label{eq:npr_cond}
\end{equation}

\begin{theorem}[Bertogna et al.\ 2011 -- Optimalität von SelectEPP]
	Das von SelectEPP gefundene EPP-Muster minimiert den Preemptions-Overhead
	von $\tau_i$ unter der Bedingung, dass alle nicht-preemptiven Regionen die
	Länge $Q$ nicht überschreiten.
\end{theorem}

\begin{proof}
	Der Algorithmus implementiert eine dynamische Programmierung über die
	Basisblöcke. Für jedes $\delta_k$ berechnet er das global optimale $\hat{C}_k$
	(kleinste kumulative WCET der ersten $k$ Blöcke). Da $\hat{C}_k$ das Minimum
	über alle machbaren Vorgänger $\delta_j$ bildet, werden alle möglichen
	machbaren EPP-Platzierungen für die ersten $k$ Blöcke berücksichtigt.
	Durch Induktion über $k$ ist $\hat{C}_N = C_i^{\min}$ die minimale WCET
	des gesamten Tasks. \hfill$\blacksquare$
\end{proof}

\subsection{Laufzeitanalyse von SelectEPP}

\begin{satz}[Laufzeit SelectEPP]
	SelectEPP hat eine Laufzeit von $O(N^2)$, wobei $N$ die Anzahl der Basisblöcke
	ist. Für jeden der $N$ Basisblöcke wird über die Menge $\mathit{Prev}_k$ (höchstens
	$N$ Einträge) iteriert. Der Gesamtaufwand für das gesamte Task-Set mit $n$
	Tasks ist $O(\sum_i N_i^2)$.
\end{satz}

\section{Vergleich der Ansätze}

\begin{table}[H]
	\centering
	\caption{Bewertung der Limited-Preemption-Methoden}
	\begin{tabular}{p{3cm}p{2cm}p{2.5cm}p{2.5cm}p{2.5cm}}
		\toprule
		\textbf{Methode} & \textbf{Impl.-Aufwand} & \textbf{Vorhersage\-barkeit} & \textbf{Effektivität} & \textbf{Optimalität} \\
		\midrule
		Non-Preemptive  & Gering   & Hoch    & Niedrig  & $U_{\text{lub}} = 0$ \\
		Preem. Threshold & Gering  & Niedrig & Mittel   & Schwach optimal \\
		Deferred Preem. & Mittel   & Mittel  & Hoch     & Lokal optimal \\
		Task Splitting   & Mittel  & Hoch    & Hoch     & Scharf optimal \\
		\bottomrule
	\end{tabular}
\end{table}

%=============================================================
\chapter{Überlastbehandlung -- Scheduling unter Überlast}
\label{chap:uberlast}
%=============================================================

\section{Einführung: Arten und Ursachen der Überlast}

Überlast ($U > 1$) tritt auf, wenn die Gesamtrechenanforderung die verfügbare
Prozessorkapazität überschreitet. Ursachen sind:
\begin{itemize}
	\item \textbf{Schlechtes Systemdesign}: Keine Worst-Case-Analyse durchgeführt.
	\item \textbf{Gleichzeitige Ereignisse}: Peak-Load-Szenarien überschätzt.
	\item \textbf{Hardwarefehler}: Anomale Interrupt-Sequenzen sättigen CPU.
	\item \textbf{Umgebungsveränderungen}: Unvorhergesehene Lastspitzen.
\end{itemize}

\begin{definition}[Instantane Last]
	Die \textbf{instantane Last} zum Zeitpunkt $t$ ist:
	\[
		\rho(t) = \max_i \rho_i(t), \qquad
		\rho_i(t) = \frac{\sum_{d_k \leq d_i} c_k(t)}{d_i - t},
	\]
	wobei $c_k(t)$ die verbleibende Ausführungszeit von Job $J_k$ ist.
	Ein System ist überlastet, wenn $\rho(t) > 1$.
\end{definition}

\section{Aperiodische Überlast -- Value-Based Scheduling}

Bei aperiodischer Überlast ist es nicht möglich, alle Deadlines einzuhalten.
Dann ist das Ziel, den \emph{Gesamtwert} der rechtzeitig fertiggestellten Tasks
zu maximieren.

\begin{definition}[Value Function]
	Jede Task $J_i$ hat eine \textbf{Wertfunktion} $V_i(t)$, die den Nutzen angibt,
	wenn sie zum Zeitpunkt $t$ fertiggestellt wird:
	\[
		V_i(t) = \begin{cases} v_i & t \leq d_i \\ 0 & t > d_i \end{cases}
	\]
	(harte Wertfunktion) oder eine monoton fallende Funktion (weiche Wertfunktion).
\end{definition}

\subsection{Algorithmus EARLIEST DEADLINE LATE (EDL) / RED}

Unter Überlast versucht der \textbf{Robust Earliest Deadline (RED)} Algorithmus,
den Gesamtwert zu maximieren, indem Tasks mit niedrigem Verhältnis $v_i/C_i$
abgelehnt werden (Value Density):

\begin{algorithm}[H]
	\caption{Value-Based Scheduling unter Überlast (vereinfacht)}
	\begin{algorithmic}[1]
		\State Sortiere aktive Tasks nach $v_i/C_i$ absteigend
		\State $\rho \leftarrow \rho(t_{\text{aktuell}})$
		\While{$\rho > 1$}
		\State Verwerfe Task mit kleinstem $v_i/C_i$
		\State Berechne $\rho$ neu
		\EndWhile
		\State Plane verbleibende Tasks mit EDF ein
	\end{algorithmic}
\end{algorithm}

\section{Überlastbehandlung unter RM vs. EDF}

\subsection{Permanente Überlast unter RM}

Unter RM mit fester Priorität bei $U > 1$:
\begin{itemize}
	\item \textbf{Alle niederprioren Tasks werden vollständig blockiert.}
	\item Tasks mit $P_i < P_j$ für alle $j$ erhalten keinen Prozessor mehr.
	\item Das System degeneriert: Nur hochpriore Tasks laufen; alle anderen
	      verpassen ihre Deadlines.
\end{itemize}

\subsection{Permanente Überlast unter EDF -- Cervin-Theorem}

\begin{theorem}[Cervin 2003 -- EDF bei permanenter Überlast]
	\label{thm:cervin}
	Bei $n$ periodischen Tasks mit $U > 1$ unter EDF gilt im eingeschwungenen
	Zustand: Die effektive mittlere Periode jeder Task $\tau_i$ beträgt
	\[
		\bar{T}_i = T_i \cdot U.
	\]
	EDF führt also eine \textbf{automatische Periodenskalierung} durch.
\end{theorem}

\begin{proof}[Beweisidee]
	In einem EDF-Schedule wird jede Task mit relativer Häufigkeit $C_i/T_i$ ausgeführt
	(proportional zu ihrer Auslastung). Bei $U > 1$ wird die verfügbare Zeit $1$ auf
	die Tasks im Verhältnis ihrer Auslastungen aufgeteilt. Task $\tau_i$ erhält den
	Anteil $U_i/U = C_i/(T_i U)$ der Prozessorzeit, was einer effektiven Rate von
	$1/\bar{T}_i = 1/(T_i U)$ entspricht. \hfill$\blacksquare$
\end{proof}

\begin{satz}[Vergleich RM vs. EDF bei Überlast]
	\begin{itemize}
		\item \textbf{RM}: Niederpriore Tasks werden vollständig blockiert;
		      hochpriore Tasks laufen normal. Ungerecht bei Überlast.
		\item \textbf{EDF}: Alle Tasks skalieren gleichmäßig; keine Task wird
		      vollständig blockiert. Fairer Umgang mit Überlast.
	\end{itemize}
	Dies ist ein weiterer Vorteil von EDF gegenüber RM, zusätzlich zur höheren
	Auslastungsschranke.
\end{satz}

\section{Elastic Scheduling -- Dynamische Periodeneanpassung}

\begin{definition}[Elastic Task]
	Eine elastische periodische Task $\tau_i$ hat einen nominellen Periodenbereich
	$[T_i^{\min}, T_i^{\max}]$ und einen Elastizitätskoeffizienten $e_i \geq 0$.
	Bei Überlast kann die Periode erhöht (und damit die Auslastung reduziert) werden,
	mit einem ``Kostenfaktor'' proportional zu $e_i$.
\end{definition}

\begin{satz}[Elastic Scheduling Feasibility]
	Ein elastisches Task-Set ist unter EDF schedulierbar, wenn für die optimalen
	Perioden $T_i^*$ gilt:
	\[
		\sum_{i=1}^n \frac{C_i}{T_i^*} \leq 1, \quad T_i^{\min} \leq T_i^* \leq T_i^{\max}.
	\]
	Die optimalen Perioden können durch einen iterativen Algorithmus in $O(n^2)$
	berechnet werden.
\end{satz}

\section{$(\pi, \lambda)$-Algorithmen für harte Überlast}

Für Systeme, in denen nicht alle Tasks weich sind, werden \textbf{Overload Manager}
eingesetzt, die Entscheidungen über Akzeptanz oder Ablehnung neuer Tasks treffen.

\begin{definition}[Admission Control]
	Ein \textbf{Admission Controller} akzeptiert eine neue Anfrage $J_{\text{new}}$,
	wenn das erweiterte System $\mathcal{J} \cup \{J_{\text{new}}\}$ schedulierbar
	bleibt:
	\[
		\rho(\mathcal{J} \cup \{J_{\text{new}}\}, t) \leq 1.
	\]
	Andernfalls wird $J_{\text{new}}$ abgelehnt oder eine niederpriore Task
	verdrängt.
\end{definition}

\begin{satz}[Güte von EDF-basiertem Overload Management]
	EDF maximiert unter allen Online-Algorithmen bei harter Überlast den
	erwarteten Gesamtwert, wenn Tasks homogene Wertfunktionen haben (d.\,h.\
	$v_i/C_i = \mathrm{const}$).
\end{satz}

%=============================================================
\chapter{Der EDF\texorpdfstring{$^+$}{+} Algorithmus -- Optimales Scheduling mit Blocking-Penalty}
\label{chap:edfplus}
%=============================================================

\section{Motivation und Einordnung}

Dieses Kapitel präsentiert \textbf{EDF$^+$} (Earliest Deadline First Plus),
einen neuartigen Echtzeit-Scheduling-Algorithmus, der die klassische EDF-Strategie
um ein dynamisches Blocking-Penalty-Verfahren erweitert \cite{Epp2025}.

Der fundamentale Kernsatz gilt auch für EDF$^+$:

\begin{quote}
\textbf{Wenn es einen machbaren Schedule gibt, findet EDF diesen sofort.
Die kürzeste Ausführungszeit aller Tasks findet EDF sofort. EDF ist scharf optimal.}
\end{quote}

EDF$^+$ erweitert diese scharfe Optimalität auf Systeme mit stochastischen
Blocking-Ereignissen. Während klassisches EDF bei Blocking-Ereignissen seine
Optimalität verlieren kann, stellt EDF$^+$ sicher, dass nach jedem
Blocking-Ereignis der optimale EDF-Schedule für die verbleibenden Tasks
berechnet wird.

EDF$^+$ eignet sich besonders für eingebettete Systeme nach AUTOSAR und
ISO~26262, da der WCET-Overhead durch Neuplanungen analytisch begrenzt ist
und der Algorithmus selbststabilisierend und fair ist.

\section{Grundlegende Definitionen}

\begin{definition}[Task-System]
	Ein \textbf{periodisches Task-System} ist ein Tupel
	\[
		\Sigma = \bigl(\mathcal{T},\, \mathbf{C},\, \mathbf{D},\, \mathbf{P}\bigr)
	\]
	mit Task-Menge $\mathcal{T} = \{T_1,\ldots,T_n\}$, Worst-Case-Ausführungszeiten
	$\mathbf{C} \in \mathbb{R}_{>0}^n$, absoluten Deadlines $\mathbf{D} \in \mathbb{R}_{>0}^n$
	und Perioden $\mathbf{P} \in \mathbb{R}_{>0}^n$ mit $D_i \le P_i$.
\end{definition}

\begin{definition}[Blocking-Ereignis]
	Task $T_i$ ist zum Zeitpunkt $t$ \textbf{blockiert}, falls sie nicht
	fortschreiten kann, obwohl der Prozessor ihr zugeteilt ist. Die Blocking-Indikatorvariable $B_i \in \{0,1\}$ erfüllt
	\[
		P(B_i = 1) = p, \quad P(B_i = 0) = 1-p,
	\]
	wobei $p \in (0,1)$ die Blocking-Wahrscheinlichkeit ist.
\end{definition}

\begin{definition}[Effektive Deadline]
	Die \textbf{effektive Deadline} von $T_i$ nach einem Blocking-Ereignis ist
	\[
		D_i^{\text{eff}} = \begin{cases}
			D_i      & \text{falls } T_i \text{ nicht blockiert hat,} \\
			+\infty  & \text{falls } T_i \text{ blockiert hat.}
		\end{cases}
	\]
\end{definition}

\begin{definition}[Geometrische Verteilung der Blockings]
	Die Anzahl der Ausführungsversuche bis zur erfolgreichen Ausführung von $T_i$
	folgt einer geometrischen Verteilung mit Parameter $p$:
	\[
		P(X = k) = (1-p)^{k-1} \cdot p, \quad k = 1, 2, 3, \ldots
	\]
	mit Erwartungswert $E[X] = 1/p$.
\end{definition}

\section{Der EDF$^+$-Algorithmus}

\subsection{Algorithmus-Beschreibung}

EDF$^+$ arbeitet in fünf Phasen:

\textbf{Phase 1: EDF-Initialisierung}
\[
	\sigma^{(0)} = \mathrm{MergeSort}(\mathcal{T},\, \mathbf{D}).
\]

\textbf{Phase 2: Ausführung} -- Tasks werden in Reihenfolge $\sigma^{(k)}$ ausgeführt.

\textbf{Phase 3: Blocking-Erkennung und Penalty}

Blockiert $T_i$ in Schritt $k$:
\[
	D_i^{\text{eff}} \leftarrow +\infty \quad \text{(niedrigste Priorität)}.
\]

\textbf{Phase 4: Sofortige Neuplanung}
\[
	\sigma^{(k+1)} = \mathrm{MergeSort}\!\bigl(\mathcal{T},\, \mathbf{D}^{\text{eff}}\bigr).
\]

\textbf{Phase 5: Terminierung} -- Wiederholung bis alle Tasks abgearbeitet sind.

\subsection{Pseudocode}

\begin{algorithm}[H]
	\caption{EDF$^+$ Scheduling Algorithm}
	\begin{algorithmic}[1]
		\Require Task-Menge $\mathcal{T}=\{T_1,\ldots,T_n\}$, Deadlines $D$, Periode $\mathcal{P}$
		\Ensure Schedule $\pi$ über $[0,\mathcal{P}]$
		\State $\mathbf{D}^{\text{eff}} \leftarrow \mathbf{D}$
		\State $\sigma \leftarrow \textsc{MergeSort}(\mathcal{T}, \mathbf{D}^{\text{eff}})$
		\State $\mathcal{Q} \leftarrow \sigma$, $t \leftarrow 0$, $\pi \leftarrow \emptyset$
		\While{$\mathcal{Q} \neq \emptyset$ \textbf{and} $t < \mathcal{P}$}
		\State $T_i \leftarrow \textsc{Dequeue}(\mathcal{Q})$
		\State $t_{\text{start}} \leftarrow t$
		\State \textbf{Führe} $T_i$ aus bis Fertigstellung \textbf{oder} Blocking
		\If{$T_i$ \textbf{blockiert}}
		\State $D_i^{\text{eff}} \leftarrow +\infty$
		\State $\mathcal{Q} \leftarrow \textsc{MergeSort}(\mathcal{Q} \cup \{T_i\},\, \mathbf{D}^{\text{eff}})$
		\State \textbf{continue}
		\EndIf
		\State $\pi \leftarrow \pi \cup \{(T_i, t_{\text{start}}, t)\}$
		\State $t \leftarrow t + C_i$
		\EndWhile
		\State \Return $\pi$
	\end{algorithmic}
\end{algorithm}

\section{Optimalitätsbeweis für EDF$^+$}

\subsection{Fundament: EDF-Optimalitätstheorem}

Das gesamte Optimalitätsargument für EDF$^+$ beruht auf dem in Kapitel~\ref{chap:aperiodisch}
bewiesenen EDF-Optimalitätstheorem (Satz~\ref{thm:edf_optimal}):

\begin{satz}[EDF-Optimalitätstheorem -- Wiederholung]
	\label{satz:edf_opt_wieder}
	Sei $\Sigma$ ein präemptives Single-Processor-Task-System mit $U \leq 1$ und
	Blocking-Wahrscheinlichkeit $p = 0$. Dann findet EDF stets einen gültigen Schedule:
	\[
		U \leq 1, p = 0 \implies \text{EDF verpasst keine Deadline.}
	\]
\end{satz}

Dieses Theorem ist das zentrale Werkzeug: EDF$^+$ erzeugt nach jedem Blocking-Ereignis
eine Teilmenge von Tasks, die die Voraussetzungen von Satz~\ref{satz:edf_opt_wieder}
erfüllt.

\subsection{Hauptsatz: Optimalität von EDF$^+$}

\begin{satz}[Optimalität von EDF$^+$]
	\label{satz:edfplus_optimal}
	EDF$^+$ ist \textbf{optimal} für präemptive Single-Processor-Systeme mit
	stochastischen Blocking-Ereignissen: Wenn ein schedulearbares Task-System mit
	Blocking-Wahrscheinlichkeit $p \in [0,1)$ existiert, so findet EDF$^+$ einen
	gültigen Schedule, der alle erreichbaren Deadlines einhält.
\end{satz}

\begin{proof}
	\textbf{Teil 1: Basisfall -- kein Blocking ($p = 0$).}
	Ohne Blocking reduziert sich EDF$^+$ auf klassisches EDF:
	$D_i^{\text{eff}} = D_i$ für alle $i$. Das System erfüllt $U \leq 1$ und $p = 0$.
	Nach Satz~\ref{satz:edf_opt_wieder} gilt: EDF$^+$ ($\equiv$ EDF) verpasst keine Deadline.

	\medskip
	\textbf{Teil 2: Penalty-Schritt -- eine blockierte Task ($p > 0$, ein Blocking-Ereignis).}
	Blockiert $T_i$ zum Zeitpunkt $t^* < D_i$. EDF$^+$ setzt
	$D_i^{\text{eff}} \leftarrow +\infty$, was $T_i$ aus dem laufenden Wettbewerb
	ausschließt. Die verbleibende Task-Menge
	$\mathcal{T}' = \{T_j \in \mathcal{T} \mid D_j^{\text{eff}} < +\infty\}$
	hat Auslastung:
	\[
		U' = U - \frac{C_i}{P_i} < U \leq 1.
	\]
	Für alle $T_j \in \mathcal{T}'$ gilt zum Zeitpunkt $t^*$: $p_j = 0$.
	Damit erfüllt $\mathcal{T}'$ die Voraussetzungen von Satz~\ref{satz:edf_opt_wieder}:
	\[
		U' \leq 1, p' = 0 \implies \text{EDF findet auf } \mathcal{T}' \text{ einen Deadline-freien Schedule.}
	\]

	\medskip
	\textbf{Teil 3: Induktionsschritt -- $R$ Blocking-Ereignisse.}
	Per vollständiger Induktion über $R$: Beim $R$-ten Blocking von $T_{i_R}$
	erzeugt EDF$^+$ die Menge:
	\[
		\mathcal{T}^{(R)} = \bigl\{T_j \mid D_j^{\text{eff}} < +\infty\bigr\}, \quad
		U^{(R)} = U - \sum_{k=1}^{R} \frac{C_{i_k}}{P_{i_k}} < U \leq 1.
	\]
	Auf $\mathcal{T}^{(R)}$ sind keine Tasks blockiert: $p^{(R)} = 0$. Nach
	Satz~\ref{satz:edf_opt_wieder} verpasst EDF auf $\mathcal{T}^{(R)}$ keine Deadline.
	EDF$^+$ berechnet via Merge Sort exakt diesen optimalen Schedule. Per
	Induktionshypothese war die Behandlung der vorherigen $R-1$ Blockings bereits
	optimal. Die Gesamtaussage folgt durch vollständige Induktion. \hfill$\blacksquare$
\end{proof}

\subsection{Korollare}

\begin{korollar}[Dominanz über RMS]
	Für jedes präemptive Task-System gilt:
	\[
		\text{Miss-Rate}(\text{EDF}^+) \leq \text{Miss-Rate}(\text{RMS}).
	\]
	Die Ungleichung ist strikt für $U > \ln 2 \approx 0{,}693$.
\end{korollar}

\begin{proof}
	RMS ist nach Liu \& Layland \cite{Liu1973} nur für $U \leq \ln 2$
	deadline-safe. Für $U > \ln 2$ können RMS-Schedules Deadlines verfehlen.
	EDF$^+$ hingegen reduziert bei jedem Blocking auf eine neue EDF-Instanz
	$\mathcal{T}^{(k)}$ mit $U^{(k)} < U \leq 1$; auf diese wendet
	Satz~\ref{satz:edf_opt_wieder} an, der die Deadline-Freiheit garantiert.
	\hfill$\blacksquare$
\end{proof}

\begin{korollar}[Minimax-Eigenschaft]
	Unter allen Scheduling-Algorithmen, die bei Blocking-Ereignissen neu planen,
	minimiert EDF$^+$ die erwartete Anzahl verpasster Deadlines:
	\[
		E[\,\text{Misses}_{\text{EDF}^+}\,] = \min_{\mathcal{A}} \,E[\,\text{Misses}_{\mathcal{A}}\,].
	\]
\end{korollar}

\section{Laufzeitanalyse}

\begin{satz}[Laufzeit von EDF$^+$]
	Die Laufzeit von EDF$^+$ für $N = 40$ Tasks mit Blocking-Wahrscheinlichkeit $p$:
	\begin{itemize}
		\item \textbf{Initialisierung:} $\Theta(n \log n)$ durch Merge Sort.
		\item \textbf{Pro Blocking-Ereignis:} $\Theta(n \log n)$ für Neuplanung.
		\item \textbf{Erwartete Neuplanungen:} $E[R] = np/(1-p)$ nach geometrischer Verteilung.
		\item \textbf{Gesamte erwartete Laufzeit:}
		      $T_{\text{gesamt}} = \Theta\!\left(n \log n \cdot \frac{1}{1-p}\right)$.
	\end{itemize}
\end{satz}

Tabelle~\ref{tab:edfplus_blocking} zeigt die erwartete Anzahl von Neuplanungen
für $N = 40$ Tasks:

\begin{table}[H]
	\centering
	\caption{Erwartete Neuplanungen für EDF$^+$ bei $N=40$ Tasks}
	\label{tab:edfplus_blocking}
	\begin{tabular}{ccc}
		\toprule
		$p$ (Blocking) & $E[R]$ Neuplanungen & $T_{\text{gesamt}}$ (relativ zu EDF) \\
		\midrule
		0{,}01 & 0{,}40  & 1{,}01$\times$ \\
		0{,}05 & 2{,}11  & 1{,}05$\times$ \\
		0{,}10 & 4{,}44  & 1{,}11$\times$ \\
		0{,}20 & 10{,}0  & 1{,}25$\times$ \\
		0{,}50 & 40{,}0  & 2{,}00$\times$ \\
		\bottomrule
	\end{tabular}
\end{table}

\section{EDF$^+$ für dynamische Task-Mengen}

Die Erweiterung auf dynamisch wachsende Task-Mengen ist möglich. Bei Ankunft
einer neuen Task $T_{n+1}$ zum Zeitpunkt $a_{n+1}$ wird die laufende Ausführung
von $T_{\text{cur}}$ preemptiert, falls $D_{n+1} < D_{\text{cur}}$, und der
Schedule für die erweiterte Menge neu berechnet.

\begin{satz}[Optimalität von EDF$^+$ für dynamische Task-Mengen]
	EDF$^+$ ist optimal für dynamische Task-Mengen mit stochastischen
	Blocking-Ereignissen und dynamischen Task-Ankünften. Bei jedem
	Ankunftsereignis oder Blocking-Ereignis erzeugt EDF$^+$ eine Teilmenge
	aktiver Tasks, die die Voraussetzungen des EDF-Optimalitätstheorems
	(Satz~\ref{satz:edf_opt_wieder}) erfüllt.
\end{satz}

\section{Integration in AUTOSAR und ISO~26262}

EDF$^+$ ist direkt anwendbar in eingebetteten Systemen:

\begin{itemize}
	\item \textbf{AUTOSAR Classic Platform:} EDF$^+$ kann als OS Application
	      mit Tasks nach OSEK/VDX implementiert werden.
	\item \textbf{ISO~26262 (ASIL-D):} Der analytisch begrenzte WCET-Overhead
	      durch Neuplanungen ermöglicht eine zertifizierbare Implementierung.
	\item \textbf{Vorteil gegenüber RM:} Bei Auslastungen $U > \ln 2 \approx 0{,}69$
	      (typisch in modernen ECUs) übertrifft EDF$^+$ RM systematisch.
\end{itemize}

%=============================================================
\chapter{Kernel-Design-Aspekte für Echtzeit-Systeme}
\label{chap:kernel}
%=============================================================

\section{Struktur eines Echtzeit-Kernels}

Ein Echtzeit-Kernel unterscheidet sich fundamental von einem Timesharing-Kernel:
Statt Durchschnittsdurchsatz zu maximieren, muss er Vorhersagbarkeit und die
Einhaltung individueller Zeitschranken garantieren.

\begin{definition}[Kernelprimitiven für Echtzeit]
	Ein Echtzeit-Kernel muss folgende Primitive mit \emph{beschränkter Ausführungszeit}
	bereitstellen:
	\begin{itemize}
		\item \textbf{create}($\tau$, $C$, $T$, $D$): Task erstellen mit expliziten Zeitschranken.
		\item \textbf{activate}($\tau$, $t$): Task zum Zeitpunkt $t$ aktivieren.
		\item \textbf{delay}($\tau$, $\delta$): Task für $\delta$ Zeiteinheiten suspendieren.
		\item \textbf{wait}($S$): Semaphorzugriff mit Prioritätsprotokoll.
		\item \textbf{signal}($S$): Semaphor freigeben mit Prioritätsreset.
	\end{itemize}
	Alle Primitive müssen \emph{preemptierbar} sein und dürfen keine unbeschränkten
	Schleifen enthalten.
\end{definition}

\subsection{Prozesszustände in einem Echtzeit-Kernel}

\begin{figure}[H]
	\centering
	\begin{tabular}{|c|l|}
		\hline
		\textbf{Zustand} & \textbf{Beschreibung} \\
		\hline
		RUNNING   & Task wird aktuell ausgeführt \\
		READY     & Task bereit, wartet auf CPU-Zuteilung \\
		BLOCKED   & Task wartet auf Ressource oder Ereignis \\
		SUSPENDED & Task deaktiviert (bis zum nächsten Aktivierungszeitpunkt) \\
		IDLE      & Leerlauf-Task (niedrigste Priorität) \\
		\hline
	\end{tabular}
	\caption{Prozesszustände in einem Echtzeit-Kernel}
\end{figure}

\subsection{Warteschlangen-Implementierung}

\begin{satz}[Komplexität der Scheduling-Warteschlange]
	\begin{itemize}
		\item \textbf{Feste Prioritäten (RM)}: Multi-Level-Queue mit $P$ Prioritätsstufen.
		      Einfügen und Entnehmen: $O(1)$.
		\item \textbf{Dynamische Prioritäten (EDF)}: Heap (balancierter Binärbaum).
		      Einfügen und Entnehmen: $O(\log n)$.
	\end{itemize}
	Für kleine Task-Sets ($n < 100$) ist der $O(\log n)$-Overhead von EDF
	vernachlässigbar. RM bietet nur bei sehr großen Task-Sets ($n > 1000$) einen
	wesentlichen Vorteil.
\end{satz}

\section{Interrupt-Behandlung in Echtzeit-Systemen}

\begin{satz}[Interrupt-Ansatz C -- Empfohlene Methode für Echtzeit]
	Der beste Ansatz für Echtzeit-Systeme: Interrupts aktivieren eine dedizierte
	\emph{Task} zur Gerätebehandlung, anstatt das Gerät direkt im Interrupt-Handler
	zu bedienen. Der Handler reduziert sich auf:
	\begin{algorithmic}[1]
		\State \textbf{interrupt\_handler($E$):}
		\State activate($J_E$) \Comment{Aktiviere dedizierte Task $J_E$}
		\State return
	\end{algorithmic}
	Die Gerätebehandlungs-Task $J_E$ wird dann vom Scheduler mit normaler Priorität
	verwaltet und kann von höherprioren Tasks preemptiert werden. Damit ist die
	Ausführungszeit von Interrupt-Handlern vollständig in die Schedulierbarkeitsanalyse
	integrierbar.
\end{satz}

\section{Speicherverwaltung}

\textbf{Demand Paging ist inkompatibel mit harten Echtzeit-Anforderungen},
da Page Faults zu unvorhersagbaren Verzögerungen führen. Empfohlene Alternativen:
\begin{itemize}
	\item \textbf{Statische Speicherpartitionierung}: Feste Segmente für jede Task,
	      keine Verdrängung. Vorhersagbar, aber weniger flexibel.
	\item \textbf{Memory Locking}: Alle Seiten kritischer Tasks werden im RAM gesperrt
	      (z.\,B.\ \texttt{mlockall()} unter Linux/POSIX).
	\item \textbf{Stack-Sharing unter SRP}: Durch die SRP-Eigenschaft (keine zwei Tasks
	      gleichzeitig im kritischen Abschnitt) kann ein gemeinsamer Stack verwendet
	      werden, was den Speicherbedarf erheblich reduziert.
\end{itemize}

\section{Kernel-Overhead und Schedulierbarkeitsanalyse}

In der Praxis müssen Kernel-Overheads in der Schedulierbarkeitsanalyse berücksichtigt
werden. Sei $\Delta_{\text{cs}}$ die Kontextwechselzeit und $\Delta_{\text{int}}$ der
Interrupt-Handler-Overhead, dann wird jede WCET aufgestockt zu:
\[
	C_i^{\text{eff}} = C_i + n_i \cdot \Delta_{\text{cs}} + k_i \cdot \Delta_{\text{int}},
\]
wobei $n_i$ die maximale Anzahl von Kontextwechseln und $k_i$ die maximale
Anzahl von Interrupts während der Ausführung von $\tau_i$ sind.

%=============================================================
\chapter{Anwendungsdesign für Echtzeit-Systeme}
\label{chap:anwendung}
%=============================================================

\section{Zeitschrankenableitung aus Anwendungsanforderungen}

Die kritischste und schwierigste Phase beim Entwurf eines Echtzeit-Systems
ist die Ableitung von Task-Perioden und Deadlines aus den Anwendungsanforderungen.

\begin{definition}[Nyquist-Rate für Regelungsanwendungen]
	Für eine Regelschleife mit Systemdynamik bis zur Frequenz $f_{\max}$ (Hz) muss
	die Abtastrate (und damit die Task-Periode) nach dem Nyquist-Theorem mindestens
	$T_i \leq 1/(2 f_{\max})$ betragen. In der Praxis wird $T_i \leq 1/(10 f_{\max})$
	empfohlen, um Implementierungstoleranzen zu berücksichtigen.
\end{definition}

\section{Hierarchisches Entwurfsmuster}

\begin{definition}[Hierarchische Dekomposition]
	Ein Echtzeit-System wird in Ebenen unterteilt:
	\begin{enumerate}
		\item \textbf{Sensor-Ebene}: Datenerfassung, höchste Frequenz ($T \approx 1\,\text{ms}$).
		\item \textbf{Filter-/Vorverarbeitungs-Ebene}: Rauschunterdrückung, Kalman-Filter.
		\item \textbf{Regelungs-Ebene}: PID, MPC, mittlere Frequenz ($T \approx 10\,\text{ms}$).
		\item \textbf{Planungs-Ebene}: Trajektorienplanung, niedrige Frequenz ($T \approx 100\,\text{ms}$).
		\item \textbf{Überwachungs-Ebene}: Zustandsanzeige, Logging ($T \approx 1\,\text{s}$).
	\end{enumerate}
\end{definition}

\section{Robotersteuerungsbeispiel -- Vollständige Analyse}

\begin{table}[H]
	\centering
	\caption{Task-Set einer 6-DOF-Robotersteuerung}
	\begin{tabular}{lccccl}
		\toprule
		Task & $C_i$ (ms) & $T_i$ (ms) & $U_i$ & RM-Prio & Funktion \\
		\midrule
		$\tau_1$ Sensor     & 0{,}5  & 2    & 0{,}250 & 5 & Gelenk-Encoder \\
		$\tau_2$ Filter     & 1{,}0  & 5    & 0{,}200 & 4 & Kalman-Filter \\
		$\tau_3$ Regler     & 2{,}0  & 10   & 0{,}200 & 3 & Kaskadenregelung \\
		$\tau_4$ Kinematik  & 3{,}0  & 20   & 0{,}150 & 2 & Inverse Kinematik \\
		$\tau_5$ Monitor    & 2{,}0  & 100  & 0{,}020 & 1 & Überwachung \\
		\midrule
		$\Sigma$ & & & 0{,}820 & & \\
		\bottomrule
	\end{tabular}
\end{table}

\textbf{Schedulierbarkeitsanalyse:}

$U = 0{,}820$. Liu-Layland-Bound: $U_{\text{lub}}(5) = 5(2^{1/5}-1) \approx 0{,}743$.

Da $U = 0{,}820 > 0{,}743$, gibt der Liu-Layland-Test \emph{keine Garantie}.
Wir verwenden Response Time Analysis:

$R_1 = 0{,}5 \leq 2$. \checkmark

$R_2^{(0)} = 1{,}5$; $R_2^{(1)} = 1{,}0 + \lceil 1{,}5/2 \rceil \cdot 0{,}5 = 2{,}0 > 1{,}5$;
$R_2^{(2)} = 1{,}0 + \lceil 2{,}0/2 \rceil \cdot 0{,}5 = 2{,}0$. Konvergenz: $R_2 = 2{,}0 \leq 5$. \checkmark

Analog für $\tau_3, \tau_4, \tau_5$: Alle scheduleierbar unter RM.

Da $U = 0{,}82 \leq 1$: Das System ist auch unter EDF schedulierbar -- mit mehr
Spielraum für aperiodische Tasks (TBS mit $U_s = 1-0{,}82 = 0{,}18$ möglich).

\section{AUTOSAR-Anwendungsbeispiel}

In AUTOSAR Classic wird EDF üblicherweise nicht direkt verwendet; stattdessen
nutzt AUTOSAR feste Prioritäten (RM-basiert). Mit EDF$^+$ (Kapitel~\ref{chap:edfplus})
kann jedoch ein optimales Scheduling auch mit Blocking-Ereignissen erreicht werden.

\begin{table}[H]
	\centering
	\caption{AUTOSAR-Task-Konfiguration für ein Motorsteuergerät (ECU)}
	\begin{tabular}{lccl}
		\toprule
		\textbf{Task (OSEK-Kategorie)} & $C_i$ ($\mu$s) & $T_i$ (ms) & \textbf{Funktion} \\
		\midrule
		PreemptTask10ms & 200  & 10   & Einspritzregelung (ASIL-D) \\
		PreemptTask20ms & 500  & 20   & Zündzeitpunkt (ASIL-C) \\
		PreemptTask100ms & 2000 & 100  & Diagnose (ASIL-B) \\
		BasicTask1s     & 5000 & 1000 & NVM-Update (QM) \\
		\bottomrule
	\end{tabular}
\end{table}

Auslastung: $U = 0{,}02 + 0{,}025 + 0{,}020 + 0{,}005 = 0{,}070$. Weit unter $U_{\text{lub}}$:
komfortabler Sicherheitspuffer für Spitzenlasten und aperiodische Diagnose-Tasks.

%=============================================================
\chapter{Zusammenfassung, Schlussfolgerungen und Ausblick}
\label{chap:zusammenfassung}
%=============================================================

\section{Kernaussagen dieses Buches}

Die zentralen Ergebnisse lassen sich im fundamentalen Kernsatz zusammenfassen:

\begin{quote}
\textbf{Wenn es einen machbaren Schedule gibt, findet RM diesen sofort. RM ist schwach optimal.}

\textbf{Wenn es einen machbaren Schedule gibt, findet EDF diesen sofort.
Die kürzeste Ausführungszeit aller Tasks findet EDF sofort. EDF ist scharf optimal.}
\end{quote}

\section{Schwache vs.\ Scharfe Optimalität -- Gesamtüberblick}

\begin{definition}[Schwache Optimalität von RM]
	RM ist \textbf{schwach optimal}: optimal innerhalb fester Prioritäten,
	$U_{\text{lub}} = \ln 2 \approx 0{,}693$ im worst case.
\end{definition}

\begin{definition}[Scharfe Optimalität von EDF]
	EDF ist \textbf{scharf optimal}: $U_{\text{lub}} = 1$.
	EDF scheduliert \emph{jedes} Task-Set mit $U \leq 1$.
\end{definition}

\section{EDF$^+$ -- Erweiterung der Schärfe auf Blocking-Szenarien}

EDF$^+$ \cite{Epp2025} überträgt die scharfe Optimalität von EDF auf Systeme
mit stochastischen Blocking-Ereignissen. Bei jedem Blocking erzeugt EDF$^+$ eine
Teilmenge nicht-blockierter Tasks, auf die das EDF-Optimalitätstheorem anwendbar ist.
EDF$^+$ dominiert RM für alle $U > \ln 2$ und minimiert die erwartete Anzahl
verpasster Deadlines unter allen neuplanenden Algorithmen.

\section{Praktische Empfehlungen}

\begin{enumerate}
	\item \textbf{Verwende EDF statt RM} wenn $U > 0{,}693$ und dynamische Prioritäten
	      akzeptabel sind.
	\item \textbf{Verwende EDF$^+$ statt EDF} wenn Blocking-Ereignisse möglich sind.
	\item \textbf{Verwende RM} wenn $U < 0{,}693$ und einfache $O(1)$-Implementierung
	      wichtig ist (AUTOSAR Classic, OSEK).
	\item \textbf{Implementiere SRP} wenn gemeinsame Ressourcen vorhanden sind.
	\item \textbf{Prüfe Schedulierbarkeit offline} (RTA für RM, Auslastungstest für EDF).
	\item \textbf{Verwende CBS} bei variablen Ausführungszeiten und temporaler Isolation.
\end{enumerate}

%=============================================================
\section{Ausblick: Offene Forschungsfragen und Zukünftige Entwicklungen}
%=============================================================

\subsection{Mehrprozessor-Scheduling -- Das größte offene Problem}

Das fundamentalste ungelöste Problem der Echtzeit-Theorie ist die Erweiterung
der Uniprocessor-Ergebnisse auf Mehrprozessorsysteme. Der Dhall-Effekt zeigt,
dass EDF auf Mehrprozessoren nicht mehr optimal ist.

\subsubsection{Globales vs. Partitioniertes Scheduling}

\begin{definition}[Globales Scheduling]
	Beim \textbf{globalen Scheduling} können Tasks auf jedem der $m$ Prozessoren
	ausgeführt werden. Vorteil: maximale Flexibilität. Nachteil: Task-Migrationen
	und Cache-Invalidierungen erhöhen den Overhead erheblich.
\end{definition}

\begin{definition}[Partitioniertes Scheduling]
	Beim \textbf{partitionierten Scheduling} wird jede Task statisch einem
	Prozessor zugewiesen. Die Zuweisung selbst ist ein NP-schweres Bin-Packing-Problem.
	Auf jedem Prozessor läuft ein unabhängiger Uniprocessor-Scheduler.
\end{definition}

Vielversprechende Ansätze: \emph{Semi-Partitioned Scheduling} (bestimmte Tasks
dürfen auf zwei Prozessoren migrieren) und \emph{Clustered Scheduling} (Gruppen
von Prozessoren mit gemeinsamem EDF-Scheduler).

\subsubsection{Pfair -- Optimales Mehrprozessor-Scheduling}

\begin{definition}[Pfair-Scheduling (Baruah et al.\ 1996)]
	\textbf{Pfair} (Proportionate Fair) ist das erste optimal auf homogenen
	Mehrprozessorsystemen: Jede Task $\tau_i$ erhält in jedem Zeitquantum
	$[t, t+1)$ exakt $U_i$ Zeiteinheiten (proportional zur Auslastung).
	Abweichungen um mehr als 1 Einheit werden als Pfair-Verletzung gewertet.
\end{definition}

\textbf{Nachteil:} Enorme Anzahl von Kontext-Switches (proportional zu $n$
pro Zeitquantum) macht Pfair für praktische Systeme kaum einsetzbar.
LLREF (Least Laxity and Required Execution time First) und DP-Fair
reduzieren die Preemptionsrate bei Beibehaltung optimaler Eigenschaften
für spezifische Szenarien.

\subsection{Mixed-Criticality Scheduling}

\begin{definition}[Mixed-Criticality System (MCS)]
	In einem \textbf{MCS} haben Tasks unterschiedliche Kritikalitätsstufen
	(ASIL A--D nach ISO~26262, DAL A--E nach DO-178C). Jede Task $\tau_i$
	hat zwei WCET-Werte:
	\begin{itemize}
		\item $C_i^{\mathrm{LO}}$: WCET bei niedrigem Kritikalitätsniveau (gemessen).
		\item $C_i^{\mathrm{HI}} \gg C_i^{\mathrm{LO}}$: WCET bei hohem Niveau
		      (formal verifiziert, erheblich größer).
	\end{itemize}
	Das System wechselt vom Lo-Mode in den Hi-Mode, wenn eine Task $C_i^{\mathrm{LO}}$
	überschreitet.
\end{definition}

\subsubsection{Vestal-Modell und AMC}

Das Vestal-Modell (2007) formalisierte MCS. Der \textbf{Adaptive Mixed Criticality
(AMC)} Algorithmus von Burns \& Davis (2011) ist der meistzitierte Ansatz:
RM-basiert für Lo-Mode, Wechsel zu modifizierten Deadlines im Hi-Mode.
Bekannte Grenzen: Keine Garantien für niedrigkritische Tasks im Hi-Mode;
Modusübergänge können Deadline-Verletzungen bei mittelkritischen Tasks
verursachen.

\subsubsection{EDF-VD -- Theoretisch fundierter MCS-Algorithmus}

\textbf{EDF mit Virtual Deadlines (EDF-VD)} von Baruah et al.\ (2012) weist
hochkritischen Tasks virtuelle (kürzere) Deadlines zu. Scharf optimal für
zwei Kritikalitätsstufen auf Uniprocessoren. Erweiterungen auf $k > 2$ Stufen
und Mehrprozessoren sind aktive Forschungsthemen.

\textbf{Offene Frage:} Gibt es einen universell optimalen MCS-Algorithmus,
der für alle Systeme mit $U^{\mathrm{LO}} \leq 1$ und $U^{\mathrm{HI}} \leq 1$
machbar ist?

\subsection{Maschinelles Lernen im Echtzeit-Scheduling}

Klassische ML-Modelle sind probabilistisch und inkompatibel mit harten
Echtzeit-Garantien. Dennoch gibt es vielversprechende Ansätze:

\subsubsection{ML-basierte WCET-Schätzung}

\begin{itemize}
	\item \textbf{Extreme Value Theory (EVT)}: Statistische Analyse seltener
	      Ausführungspfade für probabilistische WCET-Schranken (pWCET).
	      Ermöglicht engere WCET-Schätzungen als statische Analyse.
	\item \textbf{Random Forest / Gradient Boosting}: Schätzen WCET aus
	      Programm-Features (Schleifenanzahl, Cachestruktur, Abhängigkeitstiefe).
	\item \textbf{Messdaten-Regression}: Nutzung historischer Ausführungsdaten
	      zur Verbesserung der Analysegenauigkeit bei seltenen Pfaden.
\end{itemize}

\subsubsection{Reinforcement Learning für adaptives Scheduling}

In Systemen mit nicht vollständig bekannten Task-Parametern (Robotik, autonome
Fahrzeuge) lernen RL-basierte Scheduler online optimale Scheduling-Policies.
Da harte Garantien fehlen, wird \emph{Safe Reinforcement Learning} untersucht:
Der RL-Agent schlägt Entscheidungen vor; ein formaler Verifikator prüft
Schedulierbarkeit. Bei Verletzung greift ein klassischer sicherer Algorithmus.

\subsection{Cyber-Physical Systems und das Jitter-Problem}

\subsubsection{Control-Aware Scheduling}

Jitter in der Ausführung eines Regelalgorithmus verschlechtert direkt die
Regelgüte. Klassische Echtzeit-Theorie minimiert nur Deadlines -- nicht
Jitter. \emph{Control-Aware Scheduling} (Cervin, Henriksson, Lincoln 2003--2010)
berücksichtigt regelungstechnische Gütekriterien direkt im Scheduling-Algorithmus
und optimiert gemeinsam Scheduling und Regelgüte.

\subsubsection{Time-Sensitive Networking (TSN)}

IEEE~802.1Qbv integriert Scheduling auf Netzwerkebene. In verteilten
Echtzeit-Systemen (Industrial IoT, Industrie 5.0) muss Task-Scheduling auf
CPU-Ebene mit Netzwerkscheduling koordiniert werden:
\begin{itemize}
	\item Uhrsynchronisation (PTP/IEEE~1588) als Grundlage.
	\item Verteiltes EDF über mehrere Knoten mit Netzwerkverzögerungen.
	\item Ressourcen-Orchestrierung: CPU, GPU, FPGA und Netzwerkbandbreite
	      in heterogenen Systemen gemeinsam planen.
\end{itemize}

\subsection{Formale Verifikation von Scheduling-Algorithmen}

\subsubsection{Modellprüfung}

Werkzeuge wie UPPAAL (zeitbehaftete Automaten) und TLA+ ermöglichen die
formale Verifikation:
\begin{itemize}
	\item Nachweis der Deadline-Freiheit für konkrete Task-Sets.
	\item Verifikation von PIP/PCP auf Korrektheit und Deadlock-Freiheit.
	\item Verifikation von EDF$^+$ für endliche Blocking-Szenarien.
\end{itemize}

\textbf{Herausforderung:} Zustandsexplosion bei $n > 10$ Tasks.
Abstraktion und parametrische Verifikation (für beliebiges $n$) sind
offene Forschungsfragen.

\subsubsection{Theorem-Prover}

Isabelle/HOL, Coq und Lean ermöglichen maschinengeprüfte Beweise:
\begin{itemize}
	\item Der Liu-Layland-Bound wurde von Paulin (2006) in Isabelle formalisiert.
	\item Das EDF-Optimalitätstheorem ist für Formalisierung zugänglich
	      (Austauschargument als induktiver Beweis).
	\item \textbf{Ziel:} Vollständig verifizierte Scheduling-Bibliothek für
	      sicherheitskritische Systeme (ASIL-D-Zertifizierung).
\end{itemize}

\subsection{Der Subgraph-Algorithmus im verteilten Echtzeit-Kontext}

Der \textbf{Subgraph-Algorithmus} $\sigma_j = \sum A_{ij} \cdot 2^i + j \cdot 2^n$
\cite{Epp2025} bietet mit seiner $O(n^3)$-Komplexität neue Möglichkeiten:

\begin{itemize}
	\item \textbf{Topologieanalyse von Task-Abhängigkeitsgraphen}: Effiziente
	      Erkennung von Subgraph-Beziehungen zwischen Task-Sets.
	\item \textbf{Deduplizierung}: Identifikation äquivalenter Task-Konfigurationen
	      in großen Scheduling-Problemen.
	\item \textbf{Verteilte Systeme}: Schnelle Topologieanpassung bei Knotenausfällen
	      oder Neukonfigurationen in Echtzeit-Netzwerken.
\end{itemize}

Zukünftige GPU-Parallelisierung (Tensoroperationen für Matrixberechnungen)
könnte die Komplexität auf $O(n^2 \log n)$ reduzieren.

\subsection{Hardware-Beschleunigung für Scheduling}

Für Systeme mit sehr vielen Tasks ($n > 10.000$) oder extrem kurzen Perioden
(Mikrosekunden-Bereich) werden Hardware-beschleunigte Scheduler untersucht:

\begin{itemize}
	\item \textbf{FPGA-basierter EDF-Scheduler}: Parallele Heap-Operationen
	      in $O(1)$ statt $O(\log n)$. Erste Prototypen zeigen 10-fachen
	      Speedup.
	\item \textbf{Content-Addressable Memory (CAM)}: Hardware-Implementierung
	      von Priority Queues durch assoziative Speicherung.
	\item \textbf{Near-Memory Computing}: Scheduling-Logik direkt im
	      Speichercontroller für datenintensive Echtzeit-Workloads.
\end{itemize}

\subsection{Sicherheitszertifizierung -- ISO~26262 und POSIX}

\subsubsection{AUTOSAR Adaptive und EDF$^+$}

ISO~26262 ASIL-D fordert Nachweis der Schedulierbarkeit, beschränkte
Blockierzeiten und temporale Isolation. EDF$^+$ \cite{Epp2025} erfüllt diese
Anforderungen durch analytisch beschränkten Neuplanungs-Overhead
($E[R] = np/(1-p)$). Dies macht EDF$^+$ zum aussichtsreichen Kandidaten
für ASIL-D-Zertifizierung in AUTOSAR Adaptive Platform.

\subsubsection{POSIX SCHED\_DEADLINE}

\texttt{SCHED\_DEADLINE} implementiert seit Linux~3.14 exakt den CBS-Algorithmus
(Kapitel~\ref{chap:dynserver}) und ist der erste produktionsreife EDF-Scheduler
im Linux-Kernel. Aktuelle Forschung erweitert ihn auf Mehrprozessoren
(GEDF, partitioniertes EDF).

\subsection{Zusammenfassung der Forschungsrichtungen}

\begin{table}[H]
	\centering
	\caption{Forschungsrichtungen im Echtzeit-Scheduling}
	\begin{tabular}{p{3.8cm}p{3.8cm}p{3.8cm}}
		\toprule
		\textbf{Richtung} & \textbf{Stand} & \textbf{Offen} \\
		\midrule
		Mehrprozessor & Pfair, LLREF, DP-Fair & Effizienz, Jitter \\
		Mixed-Criticality & EDF-VD (2 Stufen) & $k > 2$ Stufen, Multicore \\
		ML im Scheduling & pWCET, Safe RL & Harte Garantien \\
		Formale Verifikation & Isabelle, UPPAAL & Skalierung, $n$ parametrisch \\
		Hardware-Scheduler & FPGA-Prototypen & Kernel-Integration \\
		Autonome Systeme & Elastic + CBS & ASIL-D Zertifizierung \\
		Verteiltes EDF & TSN, PTP & Globale Schedulierbarkeit \\
		EDF$^+$ / Subgraph & Formal bewiesen & Multicore, Zertifizierung \\
		\bottomrule
	\end{tabular}
\end{table}

\subsection{Schlusswort}

Die Echtzeit-Scheduling-Theorie hat seit Liu \& Layland (1973) enorme
Fortschritte gemacht. Die hier behandelten Algorithmen -- von RM über EDF
bis zu EDF$^+$ -- bilden eine solide Grundlage, die täglich in realen
Systemen eingesetzt wird.

Dennoch bestehen zentrale Herausforderungen: Die Lücke zwischen theoretisch
optimalen Algorithmen (Pfair) und praktisch einsetzbaren Lösungen bleibt
offen. Autonome Systeme, vernetzte Plattformen und KI-Anwendungen stellen
die Echtzeit-Theorie vor fundamental neue Fragen.

Der Kernsatz bleibt die beständige Richtschnur:

\begin{quote}
\textbf{Wenn es einen machbaren Schedule gibt, findet EDF diesen sofort.
EDF ist scharf optimal.}
\end{quote}

Alle Erweiterungen -- EDF$^+$, CBS, MCS EDF-VD -- bauen auf dieser Grundwahrheit
auf. Die Echtzeit-Scheduling-Theorie ist lebendig und offen -- ein faszinierendes
Forschungsfeld mit direkten Auswirkungen auf Sicherheit und Zuverlässigkeit
technischer Systeme.

%=============================================================
\begin{thebibliography}{99}
%=============================================================

%--- Lehrbücher und Monographien ---

\bibitem{Buttazzo2011}
G.\,C.\ Buttazzo,
\textit{Hard Real-Time Computing Systems: Predictable Scheduling Algorithms
and Applications},
3.~Aufl., New York: Springer, 2011.

\bibitem{Liu2000}
J.\,W.\,S.\ Liu,
\textit{Real-Time Systems},
Upper Saddle River, NJ: Prentice Hall, 2000.

\bibitem{Burns2001}
A.\ Burns und A.\,J.\ Wellings,
\textit{Real-Time Systems and Programming Languages},
4.~Aufl., Harlow: Addison-Wesley, 2009.

\bibitem{Kopetz2011}
H.\ Kopetz,
\textit{Real-Time Systems: Design Principles for Distributed Embedded Applications},
2.~Aufl., New York: Springer, 2011.

\bibitem{Laplante2004}
P.\,A.\ Laplante,
\textit{Real-Time Systems Design and Analysis: Tools for the Practitioner},
4.~Aufl., Hoboken, NJ: IEEE Press/Wiley, 2011.

%--- Grundlagenarbeiten Scheduling ---

\bibitem{Liu1973}
C.\,L.\ Liu und J.\,W.\ Layland,
``Scheduling algorithms for multiprogramming in a hard-real-time environment,''
\textit{Journal of the ACM}, Bd.~20, Nr.~1, S.~46--61, 1973.

\bibitem{Horn1974}
W.\,A.\ Horn,
``Some simple scheduling algorithms,''
\textit{Naval Research Logistics Quarterly}, Bd.~21, Nr.~1, S.~177--185, 1974.

\bibitem{Dertouzos1974}
M.\,L.\ Dertouzos,
``Control robotics: The procedural control of physical processes,''
in \textit{Proc.\ IFIP Congress}, S.~807--813, 1974.

\bibitem{Jackson1955}
J.\,R.\ Jackson,
``Scheduling a production line to minimize maximum tardiness,''
Research Report 43, Management Science Research Project, UCLA, 1955.

\bibitem{Lawler1973}
E.\,L.\ Lawler,
``Optimal sequencing of a single machine subject to precedence constraints,''
\textit{Management Science}, Bd.~19, Nr.~5, S.~544--546, 1973.

\bibitem{Bratley1971}
P.\ Bratley, M.\ Florian und P.\ Robillard,
``On sequencing with earliest deadlines and latest completion times,''
\textit{Journal of Computer and System Sciences}, Bd.~4, Nr.~5, S.~268--275, 1971.

\bibitem{Graham1969}
R.\,L.\ Graham,
``Bounds on multiprocessing timing anomalies,''
\textit{SIAM Journal of Applied Mathematics}, Bd.~17, Nr.~2, S.~416--429, 1969.

\bibitem{Dhall1978}
S.\,K.\ Dhall und C.\,L.\ Liu,
``On a real-time scheduling problem,''
\textit{Operations Research}, Bd.~26, Nr.~1, S.~127--140, 1978.

%--- Rate Monotonic und Schedulierbarkeitsanalyse ---

\bibitem{Lehoczky1989}
J.\,P.\ Lehoczky, L.\ Sha und Y.\ Ding,
``The rate monotonic scheduling algorithm: Exact characterization and
average case behavior,''
in \textit{Proc.\ 10th IEEE Real-Time Systems Symposium (RTSS)},
S.~166--171, 1989.

\bibitem{Leung1982}
J.\,Y.-T.\ Leung und J.\ Whitehead,
``On the complexity of fixed-priority scheduling of periodic real-time tasks,''
\textit{Performance Evaluation}, Bd.~2, Nr.~4, S.~237--250, 1982.

\bibitem{Audsley1993}
N.\,C.\ Audsley, A.\ Burns, M.\ Richardson, K.\ Tindell und A.\,J.\ Wellings,
``Applying new scheduling theory to static priority pre-emptive scheduling,''
\textit{Software Engineering Journal}, Bd.~8, Nr.~5, S.~284--292, 1993.

\bibitem{Bini2001}
E.\ Bini, G.\,C.\ Buttazzo und G.\,M.\ Buttazzo,
``A hyperbolic bound for the rate monotonic algorithm,''
in \textit{Proc.\ 13th Euromicro Conference on Real-Time Systems},
S.~59--66, 2001.

\bibitem{Joseph1986}
M.\ Joseph und P.\ Pandya,
``Finding response times in a real-time system,''
\textit{The Computer Journal}, Bd.~29, Nr.~5, S.~390--395, 1986.

%--- EDF und Deadline-Scheduling ---

\bibitem{Baruah1990}
S.\,K.\ Baruah, L.\,E.\ Rosier und R.\,R.\ Howell,
``Algorithms and complexity concerning the preemptive scheduling of
periodic, real-time tasks on one processor,''
\textit{Real-Time Systems}, Bd.~2, Nr.~4, S.~301--324, 1990.

\bibitem{Chetto1990}
H.\ Chetto, M.\ Silly und T.\ Bouchentouf,
``Dynamic scheduling of real-time tasks under precedence constraints,''
\textit{Real-Time Systems}, Bd.~2, Nr.~3, S.~181--194, 1990.

\bibitem{Stankovic1988}
J.\,A.\ Stankovic,
``Misconceptions about real-time computing: A serious problem for
next-generation systems,''
\textit{IEEE Computer}, Bd.~21, Nr.~10, S.~10--19, 1988.

%--- Server-Algorithmen ---

\bibitem{Sprunt1989}
B.\ Sprunt, L.\ Sha und J.\,P.\ Lehoczky,
``Aperiodic task scheduling for hard-real-time systems,''
\textit{Real-Time Systems}, Bd.~1, Nr.~1, S.~27--60, 1989.

\bibitem{Strosnider1995}
J.\,K.\ Strosnider, J.\,P.\ Lehoczky und L.\ Sha,
``The deferrable server algorithm for enhanced aperiodic responsiveness
in hard real-time environments,''
\textit{IEEE Transactions on Computers}, Bd.~44, Nr.~1, S.~73--91, 1995.

\bibitem{Lehoczky1992}
J.\,P.\ Lehoczky und S.\ Ramos-Thuel,
``An optimal algorithm for scheduling soft-aperiodic tasks in
fixed-priority preemptive systems,''
in \textit{Proc.\ 13th IEEE RTSS}, S.~110--123, 1992.

\bibitem{Spuri1996}
M.\ Spuri und G.\,C.\ Buttazzo,
``Scheduling aperiodic tasks in dynamic priority systems,''
\textit{Real-Time Systems}, Bd.~10, Nr.~2, S.~179--210, 1996.

\bibitem{Abeni1998}
L.\ Abeni und G.\,C.\ Buttazzo,
``Integrating multimedia applications in hard real-time systems,''
in \textit{Proc.\ 19th IEEE RTSS}, S.~4--13, 1998.

\bibitem{Tia1995}
T.\,S.\ Tia, Z.\ Deng, M.\ Shankar, M.\ Storch, J.\ Sun, L.-C.\ Wu
und J.\,W.\,S.\ Liu,
``Probabilistic performance guarantee for real-time tasks with
varying computation times,''
in \textit{Proc.\ Real-Time Technology and Applications Symp.},
S.~164--173, 1995.

%--- Ressourcenverwaltung ---

\bibitem{Sha1990}
L.\ Sha, R.\ Rajkumar und J.\,P.\ Lehoczky,
``Priority inheritance protocols: An approach to real-time synchronization,''
\textit{IEEE Transactions on Computers}, Bd.~39, Nr.~9, S.~1175--1185, 1990.

\bibitem{Baker1991}
T.\,P.\ Baker,
``Stack-based scheduling of realtime processes,''
\textit{Real-Time Systems}, Bd.~3, Nr.~1, S.~67--99, 1991.

\bibitem{Chen1990}
M.\,I.\ Chen und K.\,J.\ Lin,
``Dynamic priority ceilings: A concurrency control protocol for
real-time systems,''
\textit{Real-Time Systems}, Bd.~2, Nr.~4, S.~325--346, 1990.

%--- Limited Preemptive Scheduling ---

\bibitem{Wang1999}
Y.\ Wang und M.\ Saksena,
``Scheduling fixed-priority tasks with preemption threshold,''
in \textit{Proc.\ 6th Intl.\ Conf.\ on Real-Time Computing Systems
and Applications (RTCSA)}, S.~328--335, 1999.

\bibitem{Bertogna2011}
M.\ Bertogna, O.\ Xhani, M.\ Marinoni, F.\ Esposito und G.\,C.\ Buttazzo,
``Optimal selection of preemption points to minimize preemption overhead,''
in \textit{Proc.\ 23rd Euromicro Conference on Real-Time Systems},
S.~217--227, 2011.

%--- Überlast und Elastic Scheduling ---

\bibitem{Cervin2003}
A.\ Cervin,
\textit{Integrated Control and Real-Time Scheduling},
Diss., Lunds Tekniska Högskola, 2003.

\bibitem{Buttazzo2002}
G.\,C.\ Buttazzo und L.\ Abeni,
``Elastic scheduling for flexible workload management,''
\textit{IEEE Transactions on Computers}, Bd.~51, Nr.~3, S.~289--302, 2002.

%--- Mehrprozessor-Scheduling ---

\bibitem{BaruahPfair1996}
S.\,K.\ Baruah, N.\ Cohen, C.\,G.\ Plaxton und D.\,A.\ Varvel,
``Proportionate progress: A notion of fairness in resource allocation,''
\textit{Algorithmica}, Bd.~15, Nr.~6, S.~600--625, 1996.

\bibitem{Anderson2005}
J.\,H.\ Anderson, V.\ Bud und U.\,C.\ Devi,
``An EDF-based scheduling algorithm for multiprocessor soft real-time systems,''
in \textit{Proc.\ 17th Euromicro Conference on Real-Time Systems},
S.~199--208, 2005.

%--- Mixed-Criticality ---

\bibitem{Vestal2007}
S.\ Vestal,
``Preemptive scheduling of multi-criticality systems with varying
degrees of execution time assurance,''
in \textit{Proc.\ 28th IEEE RTSS}, S.~239--243, 2007.

\bibitem{BaruahMCS2012}
S.\ Baruah, V.\ Bonifaci, G.\ D'Angelo, H.\ Li, A.\ Marchetti-Spaccamela,
N.\ Megow und L.\ Stougie,
``Scheduling real-time mixed-criticality jobs,''
\textit{IEEE Transactions on Computers}, Bd.~61, Nr.~8, S.~1140--1152, 2012.

%--- Formale Verifikation ---

\bibitem{Paulin2006}
C.\ Paulin-Mohring,
``Scheduling-based real-time analysis,''
in \textit{Proc.\ 4th Intl.\ Workshop on Formal Methods for Industrial
Critical Systems (FMICS)}, 2006.

%--- Eigene Arbeiten ---

\bibitem{Epp2025}
S.\ Epp,
``Optimaler Scheduling-Algorithmus EDF$^+$: Prioritätsstrafe und
stochastische Analyse bei 40 Tasks,''
Universität Bielefeld, 2025.

%--- Standards und Normen ---

\bibitem{AUTOSAR2023}
AUTOSAR Consortium,
\textit{AUTOSAR Classic Platform -- Specification of Operating System},
Release~R23-11, 2023.

\bibitem{ISO26262}
ISO,
\textit{ISO~26262: Road Vehicles -- Functional Safety},
2.~Aufl., 2018.

\bibitem{POSIX2017}
IEEE und The Open Group,
\textit{POSIX.1-2017 (IEEE Std 1003.1): Base Specifications},
Issue~7, 2017.

\bibitem{TSN2016}
IEEE,
\textit{IEEE Std 802.1Qbv-2015: Enhancements for Scheduled Traffic},
2016.

\end{thebibliography}
\end{document}
